% Amélioration de la santé de la nutrition — SVT 1ère A — Cours signé OnBuch+
% Programme officiel MINESEC. Compiler avec : tectonic NA01-amelioration-sante-nutrition.tex
\newcommand{\DOCMATIERE}{SVT}
\newcommand{\DOCNIVEAU}{1ère A}
\newcommand{\DOCMODULE}{Module 1 : Le Monde Vivant}
\newcommand{\DOCLECON}{NA01}
\newcommand{\DOCTITRE}{Amélioration de la santé de la nutrition}
% OnBuch+ — préambule « cours signés » ENRICHI (compilé avec tectonic / XeLaTeX)
% Système visuel « Pop » d'OnBuch+ : fond crème (jamais blanc), bordures encre
% épaisses, ombres « dures » décalées, titres Archivo Black en capitales.
% Version MATHÉMATIQUES : graphes (pgfplots/tikz), boîtes pédagogiques
% (définition, propriété, méthode, attention, le savais-tu, exercice résolu,
% exercice, TP, bilan), page de garde et signature OnBuch+ ancrée en bas.
%
% Macros attendues dans main.tex AVANT \input{preamble} :
%   \DOCMATIERE  \DOCNIVEAU  \DOCMODULE  \DOCLECON (numéro)  \DOCTITRE
\documentclass[11pt]{article}

\usepackage{fontspec}
\usepackage[french]{babel}
\usepackage[a4paper,margin=2cm,top=3.4cm,bottom=2.8cm,headheight=1.4cm,footskip=1.3cm]{geometry}
\usepackage{amsmath,amssymb,amsfonts}
\usepackage{mathtools} % \xmapsto, \coloneqq... souvent utilisés par les modèles
\usepackage{cancel} % simplifications barrées (bilans de forces)
% \abs{z} (module, valeur absolue) et \norm{u} (norme), utilisés par les modèles
\providecommand{\abs}{}\let\abs\relax\DeclarePairedDelimiter{\abs}{\lvert}{\rvert}
\providecommand{\norm}{}\let\norm\relax\DeclarePairedDelimiter{\norm}{\lVert}{\rVert}
\usepackage[table]{xcolor} % [table] : fournit aussi \rowcolors (alternance de lignes)
\usepackage{pagecolor}
\usepackage{tikz}
\usetikzlibrary{arrows.meta,positioning,calc,shapes.geometric,shapes.misc,shapes.arrows,decorations.pathmorphing,decorations.pathreplacing,decorations.markings,patterns,shadows,mindmap,trees,fit,backgrounds,matrix,intersections,angles,quotes}
\usepackage{pgfplots}
\usepackage[version=4]{mhchem} % équations nucléaires \ce{^{235}_{92}U}
\usepackage[european,siunitx]{circuitikz} % schémas électriques
\pgfplotsset{compat=1.18}
\usepgfplotslibrary{fillbetween,groupplots} % aires entre courbes (calcul intégral)
\usepackage{enumitem}
\usepackage{fancyhdr}
% [most] charge toutes les bibliothèques tcolorbox (listings, minted, raster,
% documentation...) et gonfle inutilement la mémoire TeX sur un document long
% (~300 boîtes) : on ne charge que ce qui est réellement utilisé.
\usepackage[skins,breakable]{tcolorbox}
\usepackage{graphicx}
\usepackage{colortbl}
\usepackage{booktabs}
\usepackage{tabularx}
\usepackage{diagbox} % case d'en-tête diagonale des tableaux à double entrée
% Colonnes extensibles centrée / gauche / droite (C, L, R), souvent utilisées par les modèles
\newcolumntype{C}{>{\centering\arraybackslash}X}
\newcolumntype{L}{>{\raggedright\arraybackslash}X}
\newcolumntype{R}{>{\raggedleft\arraybackslash}X}
\usepackage{array}
\usepackage{multicol}
\usepackage{siunitx}
% parse-numbers=false : accepte aussi \SI{2,5 \times 10^{-3}}{...}
\sisetup{locale=FR,output-decimal-marker={,},group-minimum-digits=4,parse-numbers=false}
\DeclareSIUnit\ohm{\ensuremath{\symup{\Omega}}} % U+2126 absent de la police
\DeclareSIUnit\division{div} % réglages d'oscilloscope (ms/div, V/div)
\DeclareSIUnit\tr{tr} % tours (vitesse de rotation, ex. \SI{1500}{\tr\per\minute})
\DeclareSIUnit\ch{ch} % cheval-vapeur (puissance moteur, unité usuelle hors SI)
\DeclareSIUnit\franccfa{FCFA} % franc CFA (coûts, contexte camerounais)
\DeclareSIUnit\dioptre{\ensuremath{\delta}} % vergence d'une lentille (dioptrie)
\usepackage{qrcode}
% \degree : commande fréquemment utilisée par les modèles pour un angle ou une
% température en dehors de \SI{}{} (non fournie par les paquets chargés).
\providecommand{\degree}{^{\circ}}
% \qed (fin de démonstration), utilisé par les modèles sans amsthm
\providecommand{\qed}{\hfill$\square$}
% \barre : double barre des valeurs interdites dans un tableau de variations (tkz-tab)
\providecommand{\barre}{\ensuremath{\|}}
% environnement proof (amsthm non chargé) utilisé par les modèles
\ifdefined\proof\else\newenvironment{proof}[1][Démonstration]{\par\noindent\textit{#1.}\ }{\qed\par}\fi
\usepackage{lastpage}
\usepackage{hyperref}
\hypersetup{hidelinks}

% ============================================================
% Palette « Pop » OnBuch+ — mêmes hex que la classe P (plus_ui.dart)
% ============================================================
\definecolor{poporange}{HTML}{F59321}
\definecolor{popdark}{HTML}{E07A0C}
\definecolor{popcream}{HTML}{FFF6E8}
\definecolor{popink}{HTML}{1C1714}
\definecolor{poppurple}{HTML}{7A5AE0}
\definecolor{popgreen}{HTML}{2E9E6A}
\definecolor{poppink}{HTML}{E0457B}
\definecolor{popblue}{HTML}{2F7DE1}
\definecolor{popgold}{HTML}{C98A12}
\definecolor{popmuted}{HTML}{8A7F76}
\definecolor{popline}{HTML}{EADFD2}
\definecolor{popgray}{HTML}{8A7F76}
\colorlet{popgrey}{popgray}
% Alias tolérants (le modèle nomme parfois les couleurs différemment)
\colorlet{popwhite}{white}
\colorlet{popblack}{popink}
\colorlet{popred}{poppink}
\colorlet{popyellow}{popgold}
% Teintes claires pour fonds de tableaux / zones
\colorlet{popblueL}{popblue!10!white}
\colorlet{poporangeL}{poporange!14!white}
\colorlet{popgreenL}{popgreen!12!white}
\colorlet{poppurpleL}{poppurple!12!white}
\colorlet{poppinkL}{poppink!12!white}
\colorlet{popgoldL}{popgold!14!white}

\pagecolor{popcream}

% ============================================================
% Typographie
% ============================================================
\defaultfontfeatures{Ligatures=TeX}
\setmainfont{PlusJakartaSans-Medium.ttf}[Path=../../../fonts/,BoldFont=PlusJakartaSans-Bold.ttf]
% Maths en sans-serif (Fira Math) avec chiffres et lettres droites de Plus
% Jakarta, pour que les formules se fondent dans le texte.
\usepackage[math-style=ISO,bold-style=ISO]{unicode-math}
\setmathfont{FiraMath-Regular.otf}[Path=../../../fonts/]
\setmathfont{PlusJakartaSans-Medium.ttf}[Path=../../../fonts/,range={up/num,up/{Latin,latin},\mathup}]
\usepackage{icomma} % virgule décimale sans espace en mode math
% Caractères mathématiques Unicode absents des polices de texte (Plus Jakarta,
% Archivo Black) : rendus par leur équivalent mathématique, en texte comme en
% formule — sinon ils s'affichent en carré vide (ex. « Arithmétique dans ℤ »).
\usepackage{newunicodechar}
\newunicodechar{ℝ}{\ensuremath{\mathbb{R}}}
\newunicodechar{ℤ}{\ensuremath{\mathbb{Z}}}
\newunicodechar{ℂ}{\ensuremath{\mathbb{C}}}
\newunicodechar{ℕ}{\ensuremath{\mathbb{N}}}
\newunicodechar{ℚ}{\ensuremath{\mathbb{Q}}}
\newunicodechar{𝔻}{\ensuremath{\mathbb{D}}}
\newunicodechar{α}{\ensuremath{\alpha}}
\newunicodechar{β}{\ensuremath{\beta}}
\newunicodechar{γ}{\ensuremath{\gamma}}
\newunicodechar{δ}{\ensuremath{\delta}}
\newunicodechar{ε}{\ensuremath{\varepsilon}}
\newunicodechar{θ}{\ensuremath{\theta}}
\newunicodechar{λ}{\ensuremath{\lambda}}
\newunicodechar{μ}{\ensuremath{\mu}}
\newunicodechar{σ}{\ensuremath{\sigma}}
\newunicodechar{τ}{\ensuremath{\tau}}
\newunicodechar{φ}{\ensuremath{\varphi}}
\newunicodechar{ψ}{\ensuremath{\psi}}
\newunicodechar{ω}{\ensuremath{\omega}}
\newunicodechar{Σ}{\ensuremath{\Sigma}}
% majuscules produites par \MakeUppercase dans les titres (α-aminés -> Α)
\newunicodechar{Α}{\ensuremath{\alpha}}
\newunicodechar{Β}{\ensuremath{\beta}}
\newunicodechar{Γ}{\ensuremath{\Gamma}}
\newunicodechar{Δ}{\ensuremath{\Delta}}
\newunicodechar{Ε}{\ensuremath{\varepsilon}}
\newunicodechar{Θ}{\ensuremath{\Theta}}
\newunicodechar{Λ}{\ensuremath{\Lambda}}
\newunicodechar{Μ}{\ensuremath{\mu}}
\newunicodechar{Τ}{\ensuremath{\tau}}
\newunicodechar{Φ}{\ensuremath{\Phi}}
\newunicodechar{Ψ}{\ensuremath{\Psi}}
\newunicodechar{Ω}{\ensuremath{\Omega}}
\newunicodechar{↦}{\ensuremath{\mapsto}}
\newunicodechar{∈}{\ensuremath{\in}}
\newunicodechar{∉}{\ensuremath{\notin}}
\newunicodechar{∧}{\ensuremath{\wedge}}
\newunicodechar{⇔}{\ensuremath{\Leftrightarrow}}
\newunicodechar{⟺}{\ensuremath{\Longleftrightarrow}}
\newunicodechar{⇒}{\ensuremath{\Rightarrow}}
\newunicodechar{⟹}{\ensuremath{\Longrightarrow}}
\newunicodechar{∘}{\ensuremath{\circ}}
\newunicodechar{∪}{\ensuremath{\cup}}
\newunicodechar{∩}{\ensuremath{\cap}}
\newunicodechar{≡}{\ensuremath{\equiv}}
\newunicodechar{⊂}{\ensuremath{\subset}}
\newunicodechar{⊥}{\ensuremath{\perp}}
\newunicodechar{∃}{\ensuremath{\exists}}
\newunicodechar{∀}{\ensuremath{\forall}}
\newunicodechar{∛}{\ensuremath{\sqrt[3]{\,}}}
\newunicodechar{∜}{\ensuremath{\sqrt[4]{\,}}}
\newunicodechar{⃗}{\ensuremath{{}^{\to}}}
\newunicodechar{ⁿ}{\ifmmode^{n}\else\textsuperscript{\ensuremath{n}}\fi}
\newunicodechar{ˣ}{\ifmmode^{x}\else\textsuperscript{\ensuremath{x}}\fi}
\newunicodechar{ᵇ}{\ifmmode^{b}\else\textsuperscript{\ensuremath{b}}\fi}
\newunicodechar{ʳ}{\ifmmode^{r}\else\textsuperscript{\ensuremath{r}}\fi}
\newunicodechar{ᵘ}{\ifmmode^{u}\else\textsuperscript{\ensuremath{u}}\fi}
\newunicodechar{ʸ}{\ifmmode^{y}\else\textsuperscript{\ensuremath{y}}\fi}
\newunicodechar{ᵃ}{\ifmmode^{a}\else\textsuperscript{\ensuremath{a}}\fi}
\newunicodechar{ᵉ}{\ifmmode^{e}\else\textsuperscript{\ensuremath{e}}\fi}
\newunicodechar{ᵏ}{\ifmmode^{k}\else\textsuperscript{\ensuremath{k}}\fi}
\newunicodechar{ᵖ}{\ifmmode^{p}\else\textsuperscript{\ensuremath{p}}\fi}
\newunicodechar{ᴺ}{\ifmmode^{N}\else\textsuperscript{\ensuremath{N}}\fi}
\newunicodechar{ᶜ}{\ifmmode^{c}\else\textsuperscript{\ensuremath{c}}\fi}
\newunicodechar{ʰ}{\ifmmode^{h}\else\textsuperscript{\ensuremath{h}}\fi}
\newunicodechar{ᵠ}{\ifmmode^{q}\else\textsuperscript{\ensuremath{q}}\fi}
\newunicodechar{ₙ}{\ifmmode_{n}\else\textsubscript{\ensuremath{n}}\fi}
\newunicodechar{ₐ}{\ifmmode_{a}\else\textsubscript{\ensuremath{a}}\fi}
\newunicodechar{ᵢ}{\ifmmode_{i}\else\textsubscript{\ensuremath{i}}\fi}
\newunicodechar{ᵣ}{\ifmmode_{r}\else\textsubscript{\ensuremath{r}}\fi}
\newunicodechar{ₖ}{\ifmmode_{k}\else\textsubscript{\ensuremath{k}}\fi}
\newunicodechar{ᵦ}{\ifmmode_{\beta}\else\textsubscript{\ensuremath{\beta}}\fi}
\newunicodechar{ₓ}{\ifmmode_{x}\else\textsubscript{\ensuremath{x}}\fi}
\newfontfamily\popdisplay{ArchivoBlack-Regular.ttf}[Path=../../../fonts/]
\newfontfamily\poplogo{SpaceGrotesk-Bold.ttf}[Path=../../../fonts/]
\newfontfamily\popsemi{PlusJakartaSans-SemiBold.ttf}[Path=../../../fonts/]
\newfontfamily\popxbold{PlusJakartaSans-ExtraBold.ttf}[Path=../../../fonts/]
\newfontfamily\popxboldls{PlusJakartaSans-ExtraBold.ttf}[Path=../../../fonts/,LetterSpace=3.0]

\setlist[enumerate]{leftmargin=1.7em,itemsep=5pt,topsep=4pt}
\setlist[itemize]{leftmargin=1.7em,itemsep=4pt,topsep=4pt,label={\color{poporange}\textbullet}}
\setlength{\parindent}{0pt}
\setlength{\parskip}{5pt}
\renewcommand{\arraystretch}{1.25}

% pgfplots : style Pop par défaut
\pgfplotsset{
  every axis/.append style={axis line style={popink,line width=1pt},tick style={popink},
    grid style={popline,line width=0.5pt},label style={font=\small},tick label style={font=\footnotesize}},
  popcurve/.style={poporange,line width=1.8pt,no markers,smooth},
}
% Même style disponible dans un \draw TikZ simple (hors pgfplots)
\tikzset{popcurve/.style={poporange,line width=1.8pt}}
\tikzset{popgrid/.style={popline,very thin}}
\colorlet{popcurve}{poporange} % le nom du style est parfois utilisé comme couleur

% ============================================================
% Pill / squiggle
% ============================================================
\newcommand{\poppill}[3]{%
  \tikz[baseline=-0.55ex]\node[draw=popink,line width=1.2pt,rounded corners=2.6pt,%
    fill=#2,inner xsep=6.5pt,inner ysep=3pt]{%
    \popxboldls\fontsize{7}{7}\selectfont\color{#3}\MakeUppercase{#1}};%
}
\newcommand{\popsquiggle}[1]{%
  \tikz[baseline=-0.4ex]{
    \draw[#1,line width=2.6pt,line cap=round]
      plot [smooth] coordinates {(0,0.09) (0.34,0.01) (0.68,0.21) (1.02,0.01) (1.36,0.10)};
  }%
}
% Mot-clé surligné (marqueur orange sous le texte)
\newcommand{\cle}[1]{{\popxbold\color{popdark}#1}}

% ============================================================
% En-tête / pied
% ============================================================
\pagestyle{fancy}
\fancyhf{}
\renewcommand{\headrulewidth}{0pt}
\renewcommand{\footrulewidth}{0pt}
\lhead{\raisebox{-0.15ex}{\poplogo\fontsize{13}{13}\selectfont\color{popink}OnBuch\color{poporange}+}}
\rhead{\poppill{\DOCMATIERE\ \textbullet\ \DOCNIVEAU\ \textbullet\ Leçon \DOCLECON}{white}{popink}}
\lfoot{\popxbold\fontsize{7.6}{7.6}\selectfont\color{popmuted}\MakeUppercase{Cours signé OnBuch+}\ \textbullet\ {\popsemi\DOCTITRE}}
\rfoot{\tikz[baseline=-0.6ex]\node[rounded corners=8.5pt,fill=popink,minimum height=17pt,minimum width=17pt,inner xsep=5pt,inner sep=0]{\popxbold\fontsize{8}{8}\selectfont\color{popcream}\thepage\,/\,\pageref*{LastPage}};}

% ============================================================
% Titres
% ============================================================
\newcommand{\chaptitre}[2]{%
  \begin{tcolorbox}[enhanced,colback=poporange,colframe=popink,boxrule=2.6pt,arc=11pt,
      drop shadow=popink,left=15pt,right=15pt,top=12pt,bottom=11pt,before skip=3pt,after skip=18pt]
    \poppill{#2}{popink}{popcream}\par\vspace{9pt}
    {\raggedright\hyphenpenalty=10000\exhyphenpenalty=10000\popdisplay\fontsize{22}{24}\selectfont\color{popcream}\MakeUppercase{#1}\par}
  \end{tcolorbox}
}
% Section numérotée : pastille orange avec le numéro + titre Archivo
\newcounter{popsec}
\newcommand{\coursec}[1]{%
  \stepcounter{popsec}\setcounter{popsub}{0}%
  \par\vspace{16pt}\noindent
  \begin{minipage}{\linewidth}
  \tikz[baseline=-0.7ex]\node[draw=popink,line width=1.6pt,fill=poporange,rounded corners=4pt,minimum size=22pt,inner sep=0]{\popdisplay\fontsize{12}{12}\selectfont\color{popcream}\thepopsec};%
  \hspace{8pt}{\popdisplay\fontsize{15}{17}\selectfont\color{popink}\MakeUppercase{#1}}\par
  \vspace{4pt}\noindent{\color{poporange}\rule{36pt}{3.6pt}}
  \end{minipage}\par\nopagebreak\vspace{8pt}
}
\newcounter{popsub}
\newcommand{\courssub}[1]{%
  \stepcounter{popsub}%
  \par\vspace{9pt}\noindent{\popxbold\fontsize{11.5}{13}\selectfont\color{popink}{\color{poporange}\thepopsec.\thepopsub}\hspace{6pt}#1}\par\nopagebreak\vspace{4pt}
}

% ============================================================
% Boîtes pédagogiques — même recette : fond blanc, bordure épaisse colorée,
% ombre dure encre, badge plein. Titre optionnel [#1].
% ============================================================
\tcbset{popbase/.style={breakable,enhanced,colback=white,boxrule=2.3pt,arc=9pt,
  shadow={3pt}{-3pt}{0pt}{popink},left=14pt,right=14pt,top=11pt,bottom=11pt,before skip=11pt,after skip=15pt,
  fontupper=\popsemi\fontsize{10.6}{14.5}\selectfont\color{popink},
  fontlower=\popsemi\fontsize{10.6}{14.5}\selectfont\color{popink}}}
% Coche dessinée (le glyphe ✓ n'existe pas dans les polices du document)
\newcommand{\popcheck}[1]{\tikz[baseline=-0.5ex]\draw[#1,line width=1.6pt,line cap=round,line join=round](-0.13,0.0)--(-0.04,-0.09)--(0.13,0.1);}
\providecommand{\coche}{\popcheck{popgreen}}
\providecommand{\resultat}[1]{\par\smallskip\noindent\fcolorbox{popgreen}{popgreenL}{\parbox{\dimexpr\linewidth-2\fboxsep-2\fboxrule\relax}{\textbf{Résultat.} #1}}\par\smallskip}
\newcommand{\popboxhead}[3]{\poppill{#1}{#2}{white}\ifx\relax#3\relax\else\hspace{6pt}{\popxbold\fontsize{10.6}{12}\selectfont\color{popink}#3}\fi\par\vspace{7pt}}

\newtcolorbox{aretenir}[1][]{popbase,colframe=popgreen,before upper={\popboxhead{À retenir}{popgreen}{#1}}}
\newtcolorbox{exemplebox}[1][]{popbase,colframe=poporange,before upper={\popboxhead{Exemple}{poporange}{#1}}}
\newtcolorbox{definition}[1][]{popbase,colframe=popblue,before upper={\popboxhead{Définition}{popblue}{#1}}}
\newtcolorbox{propriete}[1][]{popbase,colframe=poppurple,before upper={\popboxhead{Propriété}{poppurple}{#1}}}
\newtcolorbox{methode}[1][]{popbase,colframe=popgold,before upper={\popboxhead{Méthode}{popgold}{#1}}}
\newtcolorbox{attention}[1][]{popbase,colframe=poppink,before upper={\popboxhead{Attention}{poppink}{#1}}}
\newtcolorbox{experience}[1][]{popbase,colframe=popblue,colback=popblueL,before upper={\popboxhead{Expérience}{popblue}{#1}}}
% « Le savais-tu ? » : carte violette pleine (contexte, culture, Cameroun)
\newtcolorbox{savaistu}[1][]{popbase,colback=poppurple,colframe=popink,
  fontupper=\popsemi\fontsize{10.4}{14.2}\selectfont\color{white},
  before upper={\poppill{Le savais-tu\,?}{white}{poppurple}\ifx\relax#1\relax\else\hspace{6pt}{\popxbold\color{white}#1}\fi\par\vspace{7pt}}}
% Exercice résolu : énoncé en haut, \tcblower puis la solution sur fond crème
\newcounter{popexo}\newcounter{popexr}
\newtcolorbox{exoresolu}[1][]{popbase,colframe=poporange,colbacklower=popcream,
  segmentation style={popink,line width=1.2pt,dashed},
  before upper={\stepcounter{popexr}\popboxhead{Exercice résolu \thepopexr}{poporange}{#1}},
  before lower={\poppill{Solution}{popink}{popcream}\par\vspace{6pt}}}
% Exercice (non résolu en place) — la correction est dans la partie Corrigés
\newtcolorbox{exercice}[1][]{popbase,colframe=popink,
  before upper={\stepcounter{popexo}\popboxhead{Exercice \thepopexo}{popink}{#1}}}
% Corrigé
\newtcolorbox{corrige}[1][]{popbase,colframe=popgreen,colback=popgreenL,
  before upper={\popboxhead{Corrigé}{popgreen}{#1}}}
% Objectifs / prérequis / bilan
\newtcolorbox{objectifs}{popbase,colframe=popink,colback=poporangeL,
  before upper={\popboxhead{Objectifs de la leçon}{popink}{}}}
\newtcolorbox{prerequis}{popbase,colframe=popink,colback=popblueL,
  before upper={\popboxhead{Prérequis}{popblue}{}}}
\newtcolorbox{bilan}{popbase,colframe=popink,colback=white,boxrule=2.8pt,
  before upper={\popboxhead{Fiche bilan}{poporange}{L'essentiel à connaître pour l'examen}}}
% Figure encadrée avec légende
\newtcolorbox{popfigure}[1][]{enhanced,breakable=false,colback=white,colframe=popline,boxrule=1.4pt,arc=8pt,
  left=10pt,right=10pt,top=10pt,bottom=8pt,before skip=10pt,after skip=12pt,halign=center,
  lower separated=false,halign lower=center,
  fontlower=\popsemi\fontsize{9}{11.5}\selectfont\color{popmuted},
  #1}
\newcommand{\legende}[1]{\par\vspace{6pt}{\popsemi\fontsize{9}{11.5}\selectfont\color{popmuted}#1}}

% ============================================================
% Signature OnBuch+
% ============================================================
\newcommand{\poplogoinline}[1]{{\poplogo\fontsize{#1}{#1}\selectfont\color{popink}OnBuch\color{poporange}+}}
\newcommand{\signatureonbuch}{%
  \begin{tcolorbox}[enhanced,colback=popink,colframe=popink,boxrule=2.2pt,arc=10pt,
      drop shadow=poporange,left=14pt,right=14pt,top=10pt,bottom=10pt,before skip=0pt,after skip=0pt]
    \tikz[baseline=-0.7ex]\node[circle,fill=popgreen,draw=popcream,line width=1.3pt,minimum size=22pt,inner sep=0]{\popcheck{white}};%
    \hspace{9pt}\begin{minipage}[c]{0.62\linewidth}
      {\popxbold\fontsize{12}{13}\selectfont\color{popcream}Signé\ \textbullet\ {\poplogo OnBuch\color{poporange}+}}\par\vspace{2pt}
      {\raggedright\popsemi\fontsize{8}{10.5}\selectfont\color{popline}\MakeUppercase{Équipe pédagogique OnBuch+}\\\MakeUppercase{Conforme au programme officiel MINESEC}\par}
    \end{minipage}\hfill
    {\poplogo\fontsize{22}{22}\selectfont\color{popcream}OnBuch\color{poporange}+}
  \end{tcolorbox}%
}

% ============================================================
% Page de garde
% #1 titre · #2 matière · #3 niveau · #4 module · #5 n° leçon · #6 durée ·
% #7 description / objectifs (texte ou itemize)
% ============================================================
\newcommand{\pagegarde}[7]{%
  \thispagestyle{empty}
  \newgeometry{a4paper,margin=1.7cm,top=1.6cm,bottom=1.4cm}%
  \begin{tikzpicture}[remember picture,overlay]
    \fill[poporange!18] ([xshift=-3.2cm,yshift=-3.6cm]current page.north east) circle (4.6cm);
    \fill[poppurple!14] ([xshift=2.4cm,yshift=7.6cm]current page.south west) circle (3.8cm);
  \end{tikzpicture}%
  {\poplogo\fontsize{30}{30}\selectfont\color{popink}OnBuch\color{poporange}+}\hfill
  \raisebox{0.6ex}{\poppill{Cours signé \textbullet\ Premium}{popink}{popcream}}\par
  \vspace{1pt}\popsquiggle{poppurple}\par
  \vspace{8pt}
  \begin{tcolorbox}[enhanced,colback=poporange,colframe=popink,boxrule=2.8pt,arc=13pt,
      drop shadow=popink,left=18pt,right=18pt,top=12pt,bottom=13pt,after skip=10pt]
    \poppill{#2 \textbullet\ #3}{popink}{popcream}\hspace{5pt}\poppill{#4}{white}{popink}\par\vspace{8pt}
    {\popdisplay\fontsize{14}{14}\selectfont\color{popink}LEÇON #5}\par\vspace{4pt}
    {\raggedright\hyphenpenalty=10000\exhyphenpenalty=10000\popdisplay\fontsize{25}{27}\selectfont\color{popcream}\MakeUppercase{#1}\par}
    \vspace{8pt}
    {\popxbold\fontsize{9.5}{9.5}\selectfont\color{popink}\MakeUppercase{Durée conseillée : #6}}
  \end{tcolorbox}
  \begin{tcolorbox}[enhanced,colback=white,colframe=popink,boxrule=2.2pt,arc=10pt,
      drop shadow=popink,left=16pt,right=16pt,top=10pt,bottom=10pt,after skip=8pt]
    \poppill{Dans cette leçon}{poppurple}{white}\par\vspace{5pt}
    \popsemi\fontsize{9.3}{12.6}\selectfont\color{popink}#7
  \end{tcolorbox}
  \noindent\begin{minipage}[t]{0.487\linewidth}
    \begin{tcolorbox}[enhanced,colback=popgreen,colframe=popink,boxrule=2.2pt,arc=10pt,
        drop shadow=popink,left=11pt,right=11pt,top=10pt,bottom=10pt,height=3.1cm,valign=center]
      \tcbox[colback=white,colframe=popink,boxrule=1.3pt,arc=5pt,boxsep=0pt,nobeforeafter,
        left=3pt,right=3pt,top=3pt,bottom=3pt,valign=center]{\qrcode[height=1.9cm]{https://play.google.com/store/apps/details?id=com.luvvix.onbuch&hl=fr}}%
      \hspace{8pt}\begin{minipage}[c]{0.5\linewidth}
        {\popdisplay\fontsize{11}{12.5}\selectfont\color{white}\MakeUppercase{Télécharge OnBuch}}\par\vspace{4pt}
        {\raggedright\popsemi\fontsize{8.4}{11}\selectfont\color{white}Gratuit \textbullet\ Google Play\\Tuteur IA, annales, concours, résultats\par}
      \end{minipage}
    \end{tcolorbox}
  \end{minipage}\hfill
  \begin{minipage}[t]{0.487\linewidth}
    \begin{tcolorbox}[enhanced,colback=poppurple,colframe=popink,boxrule=2.2pt,arc=10pt,
        drop shadow=popink,left=11pt,right=11pt,top=10pt,bottom=10pt,height=3.1cm,valign=center]
      \tikz[baseline=-0.7ex]\node[circle,fill=white,draw=popink,line width=1.3pt,minimum size=1.6cm,inner sep=0]{\popdisplay\fontsize{24}{24}\selectfont\color{poppurple}+};%
      \hspace{8pt}\begin{minipage}[c]{0.6\linewidth}
        {\popdisplay\fontsize{11}{12.5}\selectfont\color{white}\MakeUppercase{Passe à OnBuch+}}\par\vspace{4pt}
        {\raggedright\popsemi\fontsize{8.4}{11}\selectfont\color{white}Tous les cours signés, Tuteur IA illimité, fiches PDF — onglet OnBuch+ de l'app\par}
      \end{minipage}
    \end{tcolorbox}
  \end{minipage}
  \vspace{8pt}
  \popcontenu
  \vfill
  \signatureonbuch
  \restoregeometry
  \newpage
}

% Rangée « Ce cours contient » (page de garde)
\newcommand{\poptile}[3]{% #1 couleur #2 gros texte #3 légende
  \begin{tcolorbox}[enhanced,colback=#1,colframe=popink,boxrule=2pt,arc=9pt,shadow={2.5pt}{-2.5pt}{0pt}{popink},
      left=6pt,right=6pt,top=9pt,bottom=9pt,halign=center,valign=center,height=1.75cm,width=\linewidth,nobeforeafter]
    {\popdisplay\fontsize{12.5}{13}\selectfont\color{white}\MakeUppercase{#2}}\par\vspace{3pt}
    {\centering\popsemi\fontsize{7.8}{9.5}\selectfont\color{white}#3\par}
  \end{tcolorbox}}
\newcommand{\popcontenu}{%
  \par\noindent{\popxboldls\fontsize{7.5}{7.5}\selectfont\color{popmuted}CE COURS SIGNÉ CONTIENT}\par\vspace{7pt}
  \noindent\begin{minipage}[t]{0.235\linewidth}\poptile{popblue}{Cours}{complet, illustré, pas à pas}\end{minipage}\hfill
  \begin{minipage}[t]{0.235\linewidth}\poptile{poporange}{Méthodes}{et exercices résolus}\end{minipage}\hfill
  \begin{minipage}[t]{0.235\linewidth}\poptile{poppink}{Exercices}{type examen + corrigés}\end{minipage}\hfill
  \begin{minipage}[t]{0.235\linewidth}\poptile{popgreen}{Fiche bilan}{l'essentiel à retenir}\end{minipage}\par}

% Sommaire
\newtcolorbox{sommaire}{enhanced,colback=white,colframe=popink,boxrule=2.2pt,arc=10pt,
  drop shadow=popink,left=18pt,right=18pt,top=13pt,bottom=13pt,before skip=6pt,after skip=20pt,
  before upper={\poppill{Sommaire}{poppurple}{white}\par\vspace{9pt}\popsemi\fontsize{10.6}{14.5}\selectfont\color{popink}}}

% Signature de fin de cours — poussée tout en bas de la dernière page
\newcommand{\findecours}{%
  \par\vfill\nopagebreak
  \begin{center}\popsquiggle{poporange}\end{center}\vspace{4pt}
  \signatureonbuch
}
\AtBeginDocument{\def\llbracket{\mathopen{[\mkern-3.5mu[}}\def\rrbracket{\mathclose{]\mkern-3.5mu]}}} % intervalles d'entiers (glyphe ⟦ absent de la police)
% Glyphes absents de Fira Math : redéfinis à partir de glyphes disponibles
\AtBeginDocument{%
  \def\setminus{\mathbin{\backslash}}\let\smallsetminus\setminus
  \def\vdots{\mathord{\vbox{\baselineskip3.2pt\lineskiplimit0pt\kern4pt\hbox{.}\hbox{.}\hbox{.}}}}%
  \def\ddots{\mathinner{\mkern1mu\raise6.5pt\hbox{.}\mkern2mu\raise3.6pt\hbox{.}\mkern2mu\raise0.7pt\hbox{.}\mkern1mu}}%
  \def\triangle{\mathord{\scalebox{0.9}{$\symup{\Delta}$}}}%
  \def\nabla{\mathord{\raisebox{\depth}{\scalebox{1}[-1]{$\symup{\Delta}$}}}}%
  \def\checkmark{\popcheck{popdark}}%
  \def\star{\mathbin{\ast}}\def\bigstar{\mathord{\ast}}%
  \def\diamond{\mathbin{\scalebox{0.6}{$\Diamond$}}}%
  \def\bigcup{\mathop{\mathchoice{\raisebox{-0.3em}{\scalebox{1.7}{$\cup$}}}{\scalebox{1.25}{$\cup$}}{\cup}{\cup}}\displaylimits}%
  \def\bigcap{\mathop{\mathchoice{\raisebox{-0.3em}{\scalebox{1.7}{$\cap$}}}{\scalebox{1.25}{$\cap$}}{\cap}{\cap}}\displaylimits}%
  \def\lesseqqgtr{\mathrel{\lessgtr}}\def\bigoplus{\mathop{\oplus}\displaylimits}%
}
\providecommand{\ssi}{si et seulement si} % abréviation fréquente
\AtBeginDocument{\providecommand{\wideparen}{\overparen}} % arc de cercle

%% FIGFIX-BEGIN v5 — correctifs globaux des figures (docs/onbuch-generation/tools/figfix)
\usepackage{adjustbox}\usepackage{etoolbox}\tcbuselibrary{raster}
% toute figure TikZ/pgfplots plus large que la ligne est réduite à la largeur disponible
\BeforeBeginEnvironment{tikzpicture}{\begin{adjustbox}{max width=\linewidth}}
\AfterEndEnvironment{tikzpicture}{\end{adjustbox}}
% légendes pgfplots lisibles par-dessus les courbes
\pgfplotsset{every axis/.append style={legend style={fill=white,fill opacity=0.92,text opacity=1,draw=black!15,inner sep=3pt}}}
% étiquettes d'arêtes (syntaxe edge["..."]) avec fond blanc
\tikzset{every edge quotes/.append style={fill=white,inner sep=1.2pt}}
%% FIGFIX-END
\begin{document}

\pagegarde{Amélioration de la santé de la nutrition}{SVT}{1ère A}{Module 1 : Le Monde Vivant}{NA01}{3 h}{Cette leçon étudie les constituants de la matière vivante et des aliments, la classification des aliments simples et leurs rôles dans l'organisme. Elle aboutit à la notion d'alimentation équilibrée et au calcul de rations alimentaires adaptées à chaque individu.
\vspace{4pt}
\begin{multicols}{2}\small
\begin{itemize}
\item À la fin de cette leçon, je sais citer les principaux constituants minéraux et organiques de la matière vivante.
\item À la fin de cette leçon, je sais distinguer un aliment simple d'un aliment composé et classer les aliments simples en groupes.
\item À la fin de cette leçon, je sais mettre en évidence un aliment simple (glucides, lipides, protéines) par des tests caractéristiques.
\item À la fin de cette leçon, je sais décrire les caractéristiques physicochimiques des grands groupes d'aliments simples.
\item À la fin de cette leçon, je sais expliquer le rôle énergétique, plastique et fonctionnel des aliments.
\item À la fin de cette leçon, je sais définir l'alimentation équilibrée et citer ses conditions (âge, état physiologique, masse corporelle, état de santé).
\end{itemize}\end{multicols}}

\begin{sommaire}
\begin{multicols}{2}
\begin{enumerate}
\item Les constituants de la matière vivante
\item Aliments simples et aliments composés
\item Caractéristiques physicochimiques et tests d'identification
\item Rôle des aliments et alimentation équilibrée
\item Valeur énergétique des aliments et calcul d'un repas
\item Rations alimentaires et élaboration de menus équilibrés
\item Activité d'intégration
\item Méthodes et exercices résolus
\item Exercices
\item Corrigés des exercices
\item Fiche bilan
\end{enumerate}
\end{multicols}
\end{sommaire}

\begin{objectifs}
\begin{itemize}
\item À la fin de cette leçon, je sais citer les principaux constituants minéraux et organiques de la matière vivante.
\item À la fin de cette leçon, je sais distinguer un aliment simple d'un aliment composé et classer les aliments simples en groupes.
\item À la fin de cette leçon, je sais mettre en évidence un aliment simple (glucides, lipides, protéines) par des tests caractéristiques.
\item À la fin de cette leçon, je sais décrire les caractéristiques physicochimiques des grands groupes d'aliments simples.
\item À la fin de cette leçon, je sais expliquer le rôle énergétique, plastique et fonctionnel des aliments.
\item À la fin de cette leçon, je sais définir l'alimentation équilibrée et citer ses conditions (âge, état physiologique, masse corporelle, état de santé).
\item À la fin de cette leçon, je sais calculer la valeur énergétique d'un repas en kilojoules et en kilocalories.
\item À la fin de cette leçon, je sais élaborer un menu équilibré adapté à une situation donnée (enfant, adulte, femme enceinte, sportif).
\end{itemize}
\end{objectifs}

\begin{prerequis}
\begin{itemize}
\item Notions de chimie de base : molécules, éléments chimiques (C, H, O, N), réactions chimiques simples (Seconde).
\item Notion de cellule et de matière vivante vues en classe de Seconde.
\item Opérations de calcul : proportionnalité, produit en croix, conversions d'unités.
\item Vocabulaire courant de la nutrition : repas, ration, calorie (vie quotidienne).
\end{itemize}
\end{prerequis}

\newpage

\coursec{Les constituants de la matière vivante}

\textbf{Situation d'accroche.} À Yaoundé, la maman de Junior prépare chaque matin son ndolé avec des arachides, du miondo et un morceau de poisson fumé. Un jour, le médecin du district lui annonce que son petit frère de 3 ans souffre de \emph{malnutrition protéino-énergétique}, alors que la famille mange « à sa faim » chaque jour. Comment expliquer qu'un enfant qui mange beaucoup puisse être mal nourri ? Quels aliments faut-il associer, et en quelles quantités, pour couvrir les besoins d'un écolier, d'une femme enceinte ou d'un sportif ?

Pour répondre à ces questions, il faut d'abord répondre à une question plus fondamentale : \textbf{de quoi est fait notre corps ?} Si nous savons de quoi la matière vivante est constituée, nous comprendrons quels matériaux l'alimentation doit fournir chaque jour. C'est l'objet de cette première section.

\courssub{Observation préalable : de quoi est faite une cellule ?}

Prenons une situation simple de la vie courante. Un morceau de manioc frais laissé plusieurs jours au soleil devient léger et dur : il a perdu quelque chose. Une feuille de ndolé jetée dans un feu noircit, dégage une odeur de brûlé, puis il ne reste qu'une poudre grise très légère. Ces deux observations quotidiennes suggèrent déjà que la matière vivante contient :
\begin{itemize}
\item une substance qui s'évapore à la chaleur modérée : \cle{l'eau} ;
\item une substance qui brûle : la \cle{matière organique} ;
\item un résidu qui ne brûle pas : les \cle{minéraux} (les cendres).
\end{itemize}

\begin{aretenir}[Idée fondamentale]
La matière vivante, qu'il s'agisse d'une cellule, d'un organe ou d'un organisme entier, est constituée de \textbf{trois grandes catégories de substances} : l'\cle{eau}, les \cle{sels minéraux} et la \cle{matière organique}. L'eau et les sels minéraux forment les \textbf{constituants minéraux} ; la matière organique forme les \textbf{constituants organiques}.
\end{aretenir}

\courssub{Les constituants minéraux de la matière vivante}

\begin{definition}[Constituants minéraux et constituants organiques]
On appelle \cle{constituants minéraux} de la matière vivante l'ensemble formé par \textbf{l'eau} et les \textbf{sels minéraux} (calcium, phosphore, fer, sodium, potassium, iode, etc.). On appelle \cle{constituants organiques} l'ensemble des substances carbonées propres au vivant : les \textbf{glucides} (sucres), les \textbf{lipides} (graisses), les \textbf{protéines} (protides), les \textbf{vitamines} et les \textbf{acides nucléiques}.
\end{definition}

\textbf{1. L'eau, constituant le plus abondant.}

Observe un nouveau-né : sa peau est souple, rebondie, « pleine d'eau ». Un vieillard, au contraire, a la peau ridée et desséchée. Ce n'est pas un hasard : la proportion d'eau diminue avec l'âge. Chez l'adulte, l'eau représente environ \textbf{65 à 75 \% de la masse corporelle} : un homme de \SI{70}{kg} contient donc environ \SI{45}{kg} d'eau, soit près de 45 litres !

\begin{propriete}[Abondance de l'eau (admise)]
L'eau est le constituant le plus abondant de la matière vivante : elle représente environ les \textbf{deux tiers de la masse corporelle humaine} (65 à 75 \% selon l'âge, le sexe et l'état physiologique). Aucune vie n'est possible sans eau.
\end{propriete}

L'eau joue un double rôle :
\begin{itemize}
\item un \textbf{rôle plastique} (de construction) : elle entre dans la constitution des cellules et leur donne leur volume et leur turgescence ;
\item un \textbf{rôle fonctionnel} : elle est le solvant dans lequel se déroulent toutes les réactions chimiques de la cellule ; elle transporte les nutriments et les déchets (sang, lymphe) ; elle participe à la régulation de la température du corps (transpiration).
\end{itemize}

\textbf{2. Les sels minéraux.}

Les sels minéraux ne représentent qu'environ 4 à 5 \% de la masse corporelle, mais ils sont indispensables. Voici les principaux :

\begin{center}
\begin{tabularx}{\linewidth}{|l|X|X|}
\hline
\rowcolor{popblueL} \textbf{Sel minéral} & \textbf{Rôle principal} & \textbf{Conséquence d'une carence} \\
\hline
Calcium (Ca) & Construction des os et des dents ; contraction musculaire & Rachitisme chez l'enfant, fragilité osseuse \\
\hline
Phosphore (P) & Constitution des os, des membranes et de l'ATP (énergie) & Fatigue, troubles de la croissance \\
\hline
Fer (Fe) & Constitution de l'hémoglobine des globules rouges & Anémie (fatigue, pâleur) \\
\hline
Sodium (Na) et Potassium (K) & Équilibre hydrique ; transmission de l'influx nerveux & Troubles nerveux et musculaires \\
\hline
Iode (I) & Fabrication des hormones thyroïdiennes & Goitre, retard de croissance \\
\hline
\end{tabularx}
\end{center}

On distingue, selon les quantités présentes dans l'organisme :
\begin{itemize}
\item les \cle{macro-éléments} (ou éléments majeurs), présents en grande quantité : calcium, phosphore, potassium, sodium, magnésium, chlore, soufre ;
\item les \cle{oligo-éléments}, présents à l'état de traces mais tout aussi indispensables : fer, iode, zinc, cuivre, fluor, sélénium...
\end{itemize}

\begin{attention}[Ne pas confondre « sels minéraux » et « sel de cuisine »]
Le sel de cuisine (chlorure de sodium, \ce{NaCl}) n'est qu'\textbf{un exemple} de sel minéral parmi des dizaines. Dire « les sels minéraux, c'est le sel » est une erreur classique au Probatoire. Le calcium des os, le fer du sang ou l'iode de la thyroïde sont aussi des sels minéraux, et ils ne se trouvent pas dans la salière !
\end{attention}

\begin{savaistu}[Le fer et l'anémie au Cameroun]
Le fer est un constituant indispensable de l'\cle{hémoglobine}, la molécule des globules rouges qui transporte le dioxygène dans le sang. Sa carence provoque l'\textbf{anémie ferriprive} : fatigue, pâleur, essoufflement. Au Cameroun, cette carence est fréquente chez la femme enceinte (dont les besoins en fer augmentent fortement) et chez le jeune enfant. C'est pourquoi les centres de santé recommandent une alimentation riche en fer (feuilles de ndolé, folong, gombo, haricots, viande, poisson) et distribuent parfois des suppléments de fer et d'acide folique pendant la grossesse.
\end{savaistu}

\courssub{Les constituants organiques de la matière vivante}

La matière organique est la matière \emph{carbonée} : toutes ses molécules sont bâties autour d'atomes de carbone. On distingue cinq grandes familles.

\textbf{1. Les glucides (sucres).} Ce sont les « carburants » de la cellule : ils fournissent l'énergie utilisable immédiatement. Exemples : glucose du sang, amidon du manioc et du miondo, saccharose du sucre de canne, glycogène des muscles.

\textbf{2. Les lipides (graisses).} Ce sont des réserves énergétiques concentrées, mais aussi des constituants des membranes cellulaires. Exemples : huile de palme, beurre de karité, graisses animales.

\textbf{3. Les protéines (protides).} Ce sont les « briques » du corps : elles construisent les muscles, la peau, les cheveux, mais remplissent aussi des fonctions (enzymes, anticorps, hémoglobine). Exemples : protéines du poisson fumé, des arachides, des haricots, de la viande, des œufs. C'est précisément la carence en protéines, associée au déficit énergétique, qui explique la malnutrition du petit frère de Junior : manger « à sa faim » de miondo seul ne suffit pas, car le manioc est presque exclusivement glucidique.

\textbf{4. Les vitamines.} Substances organiques nécessaires en très petites quantités, sans valeur énergétique, mais indispensables au bon fonctionnement de l'organisme (vitamines A, B, C, D...). Leur carence provoque des maladies : scorbut (vitamine C), cécité nocturne (vitamine A), rachitisme (vitamine D).

\textbf{5. Les acides nucléiques.} Ce sont les molécules de l'hérédité : l'\cle{ADN}, support de l'information génétique, et l'\cle{ARN}, qui participe à la fabrication des protéines. Ils constituent le matériel génétique de chaque cellule.

\begin{exemplebox}[Deux repas, deux destins]
Comparons deux déjeuners d'écolier à Yaoundé. \emph{Repas A} : une grande assiette de miondo seul. \emph{Repas B} : une portion de miondo, du ndolé aux arachides et un morceau de poisson fumé. Les deux repas « remplissent le ventre », mais le repas A n'apporte presque que des glucides : l'enfant reçoit de l'énergie, mais pas les protéines nécessaires à sa croissance ni suffisamment de sels minéraux et de vitamines. Le repas B, lui, associe glucides (miondo), protéines (poisson, arachides), lipides (arachides), sels minéraux et vitamines (feuilles de ndolé). C'est toute la différence entre « manger beaucoup » et « être bien nourri ».
\end{exemplebox}

\courssub{Les éléments chimiques fondamentaux de la matière vivante}

Si l'on descend au niveau des atomes, on constate que la matière vivante n'utilise qu'un petit nombre d'éléments chimiques, puisés dans l'environnement (air, eau, sol, aliments).

Quatre éléments forment à eux seuls environ \textbf{96 \% de la masse de la matière vivante} :
\begin{itemize}
\item le \textbf{carbone} (C), charpente de toutes les molécules organiques ;
\item l'\textbf{hydrogène} (H) et l'\textbf{oxygène} (O), constituants de l'eau et des molécules organiques ;
\item l'\textbf{azote} (N), constituant indispensable des protéines et des acides nucléiques.
\end{itemize}
Viennent ensuite, en plus petites quantités, le phosphore (P), le soufre (S), le calcium (Ca), le potassium (K), puis les oligo-éléments (Fe, Zn, I, Cu...).

\begin{aretenir}[À retenir absolument]
Les quatre éléments fondamentaux de la matière vivante sont \textbf{C, H, O, N} (environ 96 \% de la masse). On les retrouve dans les molécules organiques : glucides (C, H, O), lipides (C, H, O), protéines (C, H, O, N, parfois S), acides nucléiques (C, H, O, N, P).
\end{aretenir}

\begin{popfigure}
\begin{center}
\begin{tikzpicture}[scale=1.1]
% Camembert : eau 70, protéines 15, lipides 12, glucides 2, sels minéraux 1
\fill[popblue] (0,0) -- (90:2) arc (90:90-252:2) -- cycle;
\fill[poporange] (0,0) -- (90-252:2) arc (90-252:90-306:2) -- cycle;
\fill[poppurple] (0,0) -- (90-306:2) arc (90-306:90-349.2:2) -- cycle;
\fill[popgreen] (0,0) -- (90-349.2:2) arc (90-349.2:90-356.4:2) -- cycle;
\fill[popgold] (0,0) -- (90-356.4:2) arc (90-356.4:90-360:2) -- cycle;
\draw[popdark, thick] (0,0) circle (2);
% Étiquettes
\draw[thin, popmuted] (35:2) -- (35:2.6) node[right, align=left]{\textbf{Eau : 70 \%}};
\draw[thin, popmuted] (190:2) -- (190:2.7) node[left, align=right]{\textbf{Protéines : 15 \%}};
\draw[thin, popmuted] (245:2) -- (245:2.7) node[below left, align=right]{\textbf{Lipides : 12 \%}};
\draw[thin, popmuted] (300:2) -- (300:2.8) node[below right, align=left]{\textbf{Glucides : 2 \%}};
\draw[thin, popmuted] (358:2) -- (358:2.9) node[right, align=left]{\textbf{Sels minéraux : 1 \%}};
\end{tikzpicture}
\end{center}
\legende{Composition moyenne d'une cellule (en pourcentage de la masse). L'eau domine largement ; parmi la matière organique, les protéines sont les plus abondantes.}
\end{popfigure}

\courssub{Mise en évidence expérimentale des constituants de la matière vivante}

Comment un scientifique prouve-t-il ce que nous venons d'affirmer ? Suivons la démarche d'investigation complète.

\textbf{Problème :} un échantillon de matière vivante (graine de haricot, morceau de viande, feuille) contient-il de l'eau, de la matière organique et des minéraux ?

\textbf{Hypothèse :} la matière vivante contient de l'eau, de la matière organique combustible et des minéraux incombustibles.

\begin{experience}[Analyse d'un échantillon biologique]
\textbf{Matériel :} balance de précision, étuve réglée à \SI{100}{\celsius}, creuset, bec Bunsen ou four à moufle, échantillon (par exemple une graine de haricot écrasée).

\textbf{Protocole :}
\begin{enumerate}
\item \textbf{Pesée initiale} : on pèse l'échantillon frais ; soit $m_0 = \SI{10}{g}$.
\item \textbf{Dessiccation} : on place l'échantillon à l'étuve à \SI{100}{\celsius} jusqu'à masse constante ; nouvelle pesée : $m_1 = \SI{6}{g}$. La perte de masse correspond à l'\textbf{eau} évaporée.
\item \textbf{Calcination (incinération)} : on porte le résidu sec au rouge dans un creuset. La matière organique \textbf{brûle} (elle noircit, fume, dégage une odeur de brûlé) et part sous forme de gaz (\ce{CO2}, vapeur d'eau). Après refroidissement, pesée des cendres : $m_2 = \SI{0,4}{g}$.
\item \textbf{Analyse des cendres} : les cendres gris-blanc, incombustibles, contiennent les \textbf{sels minéraux}, que l'on peut identifier par des tests chimiques.
\end{enumerate}

\textbf{Observations et interprétation :}
\begin{itemize}
\item Perte d'eau : $m_0 - m_1 = \SI{4}{g}$, soit 40 \% de l'échantillon ;
\item Matière organique brûlée : $m_1 - m_2 = \SI{5,6}{g}$, soit 56 \% ;
\item Minéraux (cendres) : $m_2 = \SI{0,4}{g}$, soit 4 \%.
\end{itemize}

\textbf{Conclusion :} l'échantillon contient bien de l'eau, de la matière organique et des sels minéraux : l'hypothèse est vérifiée.
\end{experience}

\begin{popfigure}
\begin{center}
\begin{tikzpicture}[scale=1, every node/.style={font=\small}]
% Étape 1
\draw[fill=popblueL, draw=popblue, thick, rounded corners] (0,0) rectangle (2.6,1.4);
\node[align=center] at (1.3,0.7) {Échantillon frais\\ pesée : $m_0$};
% Flèche 1
\draw[-{Stealth}, very thick, poporange] (2.7,0.7) -- (4.1,0.7) node[midway, above]{Dessiccation} node[midway, below]{\SI{100}{\celsius}};
% Étape 2
\draw[fill=poporangeL, draw=poporange, thick, rounded corners] (4.2,0) rectangle (6.8,1.4);
\node[align=center] at (5.5,0.7) {Résidu sec\\ pesée : $m_1$};
% Flèche 2
\draw[-{Stealth}, very thick, poppurple] (6.9,0.7) -- (8.3,0.7) node[midway, above]{Incinération} node[midway, below]{chauffage fort};
% Étape 3
\draw[fill=poppurpleL, draw=poppurple, thick, rounded corners] (8.4,0) rectangle (11,1.4);
\node[align=center] at (9.7,0.7) {Cendres\\ pesée : $m_2$};
% Bilans
\node[align=center, text=popblue] at (1.3,-0.9) {eau perdue\\ $= m_0 - m_1$};
\draw[thin, popmuted] (3.4,0) -- (2.2,-0.6);
\node[align=center, text=poppurple] at (5.5,-0.9) {matière organique\\ $= m_1 - m_2$};
\draw[thin, popmuted] (7.6,0) -- (6.2,-0.6);
\node[align=center, text=popdark] at (9.7,-0.9) {sels minéraux\\ $= m_2$};
\draw[thin, popmuted] (9.7,0) -- (9.7,-0.6);
\end{tikzpicture}
\end{center}
\legende{Protocole de mise en évidence des constituants de la matière vivante : la dessiccation élimine l'eau, l'incinération brûle la matière organique, et les cendres révèlent les sels minéraux.}
\end{popfigure}

\begin{methode}[Analyser la composition d'un échantillon biologique]
Pour mettre en évidence les constituants d'un échantillon biologique, on réalise \textbf{successivement} :
\begin{enumerate}
\item la \textbf{dessiccation} (chauffage modéré à l'étuve) : la perte de masse mesure la quantité d'\textbf{eau} ;
\item la \textbf{calcination} (chauffage intense) : la nouvelle perte de masse mesure la \textbf{matière organique}, qui brûle ;
\item l'\textbf{analyse des cendres} : le résidu incombustible correspond aux \textbf{sels minéraux}.
\end{enumerate}
L'ordre est impératif : on ne peut pas brûler directement un échantillon frais, car on ne distinguerait pas l'eau de la matière organique.
\end{methode}

\begin{exoresolu}[Composition d'un échantillon de viande]
Un laboratoire du lycée analyse un morceau de viande fraîche de masse $m_0 = \SI{50}{g}$. Après dessiccation complète, la masse est $m_1 = \SI{14}{g}$. Après incinération, il reste $m_2 = \SI{1}{g}$ de cendres. Déterminer le pourcentage d'eau, de matière organique et de sels minéraux dans cette viande.
\tcblower
\textbf{Solution détaillée.}

\emph{Étape 1 : masse d'eau.} La dessiccation élimine l'eau :
\[ m_{\text{eau}} = m_0 - m_1 = 50 - 14 = \SI{36}{g} \quad \text{soit} \quad \frac{36}{50}\times 100 = 72\ \%. \]

\emph{Étape 2 : masse de matière organique.} L'incinération brûle la matière organique du résidu sec :
\[ m_{\text{org}} = m_1 - m_2 = 14 - 1 = \SI{13}{g} \quad \text{soit} \quad \frac{13}{50}\times 100 = 26\ \%. \]

\emph{Étape 3 : sels minéraux.} Les cendres représentent les minéraux :
\[ m_{\text{min}} = m_2 = \SI{1}{g} \quad \text{soit} \quad \frac{1}{50}\times 100 = 2\ \%. \]

\emph{Vérification :} $72 + 26 + 2 = 100\ \%$. La viande contient donc 72 \% d'eau, 26 \% de matière organique et 2 \% de sels minéraux, ce qui est conforme à la composition connue de la matière vivante.
\end{exoresolu}

\begin{attention}[Pièges fréquents au Probatoire]
\begin{itemize}
\item Ne jamais écrire que la matière vivante est « uniquement organique » : l'eau et les sels minéraux en sont des constituants à part entière.
\item Ne pas inverser les étapes : la \textbf{dessiccation} élimine l'eau (chauffage doux), la \textbf{calcination} brûle la matière organique (chauffage fort). Confondre les deux termes coûte des points.
\item Les cendres ne sont \emph{pas} de la matière organique : elles prouvent au contraire la présence des minéraux, seuls incombustibles.
\end{itemize}
\end{attention}

\begin{exercice}[Pour s'entraîner]
Une graine de maïs fraîche pèse \SI{20}{g}. Après séjour à l'étuve, sa masse tombe à \SI{17}{g}. Après incinération complète, il reste \SI{0,3}{g} de cendres.
\begin{enumerate}
\item Calculer le pourcentage d'eau, de matière organique et de sels minéraux de cette graine.
\item Comparer ces proportions à celles d'une cellule animale (figure du cours) et expliquer la différence observée pour l'eau.
\end{enumerate}
\end{exercice}

\begin{corrige}[Exercice 1]
\begin{enumerate}
\item Eau : $m_0 - m_1 = 20 - 17 = \SI{3}{g}$, soit $\frac{3}{20}\times 100 = 15\ \%$. Matière organique : $m_1 - m_2 = 17 - 0{,}3 = \SI{16,7}{g}$, soit $83{,}5\ \%$. Sels minéraux : $\SI{0,3}{g}$, soit $1{,}5\ \%$. Vérification : $15 + 83{,}5 + 1{,}5 = 100\ \%$.
\item La graine contient beaucoup moins d'eau (15 \%) qu'une cellule animale (environ 70 \%) : c'est un organe de \emph{réserve et de survie}, dont la vie est ralentie (dormance). La matière organique (amidon, protéines, lipides de réserve) y domine, ce qui permet à la graine de germer plus tard en puisant dans ses réserves.
\end{enumerate}
\end{corrige}

\textbf{Conclusion de la section.} Nous savons désormais que la matière vivante est faite d'eau, de sels minéraux et de matière organique (glucides, lipides, protéines, vitamines, acides nucléiques), elle-même bâtie essentiellement à partir de quatre éléments : C, H, O et N. L'expérience de dessiccation-incinération le prouve quantitativement. Reste une question : si notre corps est fait de ces substances, comment l'alimentation les lui fournit-elle ? C'est l'objet de la section suivante, consacrée aux aliments simples et composés.

\coursec{Aliments simples et aliments composés}

Dans la section précédente, nous avons montré que la matière vivante est constituée d'eau, de sels minéraux et de composés organiques : glucides, lipides et protéines. Une question se pose alors naturellement : si notre corps est bâti de ces matériaux, où les trouve-t-il ? La réponse tient en un mot : dans les \cle{aliments}. Mais tous les aliments se ressemblent-ils ? Un morceau de sucre et une assiette de ndolé apportent-ils la même chose à l'organisme ? C'est à cette question, en apparence simple, que cette section répond.

\courssub{Qu'est-ce qu'un aliment ?}

\textbf{Situation concrète.} Au marché Mokolo de Yaoundé, une mère de famille achète du manioc, de l'huile de palme, du poisson fumé et des feuilles de ndolé. Pourquoi achète-t-elle des denrées si différentes au lieu de se contenter d'une seule ? Parce que son corps a besoin, chaque jour, d'\emph{énergie} (pour marcher, travailler, maintenir sa température à 37 °C) et de \emph{matériaux} (pour fabriquer et réparer ses cellules). Aucun produit du marché ne fournit tout, seul.

\begin{definition}[Aliment]
Un \cle{aliment} est une substance d'origine végétale, animale ou minérale, consommée par un être vivant, qui fournit à l'organisme de l'\textbf{énergie} et des \textbf{matériaux} de construction et de fonctionnement.
\end{definition}

Chaque aliment apporte un ou plusieurs \cle{nutriments}, c'est-à-dire les substances nutritives utilisables par l'organisme : glucides, lipides, protéines, eau, sels minéraux et vitamines. Classer les aliments revient donc à répondre à une question simple : \emph{que contient cet aliment ?}

\courssub{Aliments simples et aliments composés}

\textbf{Observation.} Comparons deux produits familiers : un morceau de sucre de canne et une tasse de lait caillé (kindirmou). Le sucre, chauffé, fond et se caramélise sans laisser de résidu : il semble fait d'une seule substance. Le lait caillé, lui, laisse deviner de l'eau, une matière grasse, une partie solide qui coagule... Il semble fait de plusieurs substances. Cette observation est le point de départ d'une classification fondamentale.

\begin{definition}[Aliment simple, aliment composé]
Un \cle{aliment simple} est un aliment constitué essentiellement d'\textbf{un seul type de nutriment dominant}. Exemples : le sucre (glucides), l'huile (lipides), le blanc d'œuf (protéines).\par
Un \cle{aliment composé} est un aliment contenant \textbf{plusieurs types de nutriments} en proportions notables. Exemples : le lait, l'œuf entier, la viande, le ndolé, l'eru.
\end{definition}

\begin{propriete}[Admise]
La plupart des aliments naturels sont des \textbf{aliments composés}. Les aliments simples quasi purs sont rares dans la nature : ce sont essentiellement des produits extraits ou raffinés par l'Homme (sucre extrait de la canne, huile extraite du palmier ou de l'arachide, amidon extrait du manioc).
\end{propriete}

\begin{attention}[Piège fréquent]
Le miondo, le foutou ou le bâton de manioc ne sont \textbf{pas} des aliments simples au sens strict : ils contiennent aussi de l'eau, des sels minéraux et de petites quantités d'autres nutriments. On les \emph{classe} dans le groupe des glucides parce que c'est leur \textbf{nutriment dominant}. Dire « le miondo ne contient que des glucides » est une erreur ; dire « le miondo est classé parmi les aliments glucidiques car les glucides y dominent » est exact.
\end{attention}

\courssub{Démarche d'investigation : le lait est-il un aliment simple ou composé ?}

\begin{experience}[Analyse qualitative du lait]
\textbf{Problème :} le lait, aliment de base du nourrisson, contient-il un ou plusieurs nutriments ?\par
\textbf{Hypothèse :} le lait est un aliment composé.\par
\textbf{Protocole et observations :}
\begin{enumerate}
\item On chauffe doucement du lait dans un bécher : de la vapeur s'échappe $\rightarrow$ le lait contient de l'\textbf{eau}.
\item On ajoute quelques gouttes d'acide (vinaigre) : le lait \textbf{coagule} (formation de grumeaux de caséine) $\rightarrow$ présence de \textbf{protéines}.
\item Sur le liquide surnageant, le réactif de Fehling chauffé donne un précipité rouge brique $\rightarrow$ présence d'un sucre réducteur, le \textbf{lactose} (glucide).
\item Une goutte de lait déposée sur papier laisse une \textbf{tache translucide} persistante $\rightarrow$ présence de \textbf{lipides}.
\item L'incinération du résidu sec laisse des cendres $\rightarrow$ présence de \textbf{sels minéraux} (dont le calcium).
\end{enumerate}
\textbf{Interprétation et conclusion :} le lait contient cinq types de nutriments. C'est un \textbf{aliment composé}, même quasi complet : c'est pourquoi il suffit, seul, à nourrir le nouveau-né pendant ses premiers mois. L'hypothèse est validée.
\end{experience}

\begin{exemplebox}[Le lait, aliment composé quasi complet]
Le lait contient : de l'\textbf{eau} (environ 87 \%), des \textbf{protéines} (caséine, environ 3,5 \%), des \textbf{lipides} (environ 3,5 \%), un \textbf{glucide} (le lactose, environ 4,5 \%) et des \textbf{sels minéraux} (calcium, phosphore). Aucun autre aliment courant ne réunit seul autant de nutriments en proportions aussi équilibrées.
\end{exemplebox}

\courssub{Classification des aliments simples en trois grands groupes}

Les aliments simples se répartissent en \textbf{trois grands groupes}, correspondant aux trois grandes familles de nutriments énergétiques et plastiques, auxquels on associe l'eau, les sels minéraux et les vitamines (nutriments non énergétiques).

\textbf{1. Les glucides} (ou hydrates de carbone) : composés de carbone, d'hydrogène et d'oxygène. Ce sont les principaux \textbf{aliments énergétiques}. Exemples camerounais : manioc (bâton, miondo, gari), igname, maïs, riz, banane plantain, patate douce, sucre.

\textbf{2. Les lipides} (ou matières grasses) : composés de carbone, d'hydrogène et d'oxygène (avec beaucoup plus de carbone et d'hydrogène que les glucides). Ce sont les aliments \textbf{énergétiques de réserve}, les plus concentrés en énergie. Exemples camerounais : huile de palme, huile d'arachide, arachide, avocat, karité.

\textbf{3. Les protéines} : composés de carbone, d'hydrogène, d'oxygène, d'\textbf{azote} (et souvent de soufre) — c'est la présence d'azote qui les distingue des deux autres groupes. Ce sont les \textbf{aliments de construction} (bâtisseurs des cellules, des muscles, du sang). Exemples camerounais : poisson, viande, œufs, haricots niébé, soja, arachide (riche à la fois en protéines et en lipides).

\begin{popfigure}
\begin{center}
\begin{tikzpicture}[
  grow=right,
  level 1/.style={sibling distance=3.4cm, level distance=3.2cm},
  level 2/.style={sibling distance=1.15cm, level distance=3.6cm},
  every node/.style={font=\small},
  edge from parent/.style={draw=popmuted, thick}
]
\node[draw=popdark, fill=popdark, text=white, rounded corners, font=\bfseries] {Aliments}
  child { node[draw=poporange, fill=poporangeL, rounded corners] {Composés}
    child { node[align=left] {lait, œuf entier,\\ viande, ndolé, eru} }
  }
  child { node[draw=popblue, fill=popblueL, rounded corners] {Simples}
    child { node[draw=popgreen, fill=popgreenL, rounded corners, align=left] {Protéines\\ \footnotesize poisson, niébé, soja} }
    child { node[draw=poppink, fill=poppinkL, rounded corners, align=left] {Lipides\\ \footnotesize huile de palme, arachide, avocat} }
    child { node[draw=popgold, fill=popgoldL, rounded corners, align=left] {Glucides\\ \footnotesize manioc, igname, maïs, riz, plantain} }
  };
\end{tikzpicture}
\end{center}
\legende{Arbre de classification des aliments. Les aliments simples se rangent en trois groupes selon leur nutriment dominant ; les aliments composés contiennent plusieurs nutriments en proportions notables.}
\end{popfigure}

\begin{center}
\rowcolors{1}{}{popcream}
\begin{tabularx}{\linewidth}{|l|X|X|X|}
\hline
\rowcolor{popblueL} \textbf{Groupe} & \textbf{Éléments chimiques} & \textbf{Exemples (Cameroun)} & \textbf{Fonction principale} \\
\hline
Glucides & C, H, O & manioc, igname, maïs, riz, plantain, sucre & Fournir l'énergie immédiate (aliments énergétiques) \\
\hline
Lipides & C, H, O & huile de palme, arachide, avocat, karité & Réserve énergétique concentrée ; protection des organes \\
\hline
Protéines & C, H, O, N (S) & poisson, viande, œuf, niébé, soja & Construction et réparation des tissus (aliments bâtisseurs) \\
\hline
\end{tabularx}
\end{center}

\begin{aretenir}[L'essentiel]
Un aliment fournit énergie et matériaux. Un aliment \textbf{simple} est dominé par un seul nutriment ; un aliment \textbf{composé} en contient plusieurs. Les aliments simples se classent en trois groupes : \textbf{glucides} (énergétiques), \textbf{lipides} (réserve énergétique), \textbf{protéines} (construction), auxquels s'ajoutent l'eau, les sels minéraux et les vitamines.
\end{aretenir}

\begin{savaistu}[Le soja, protéine végétale complète]
Le soja, cultivé notamment dans le Nord-Cameroun (régions de l'Adamaoua et du Nord), est l'une des \textbf{rares sources végétales de protéines complètes} : il contient, comme la viande, tous les acides aminés indispensables à l'Homme. C'est pourquoi la farine de soja est utilisée dans les programmes de lutte contre la malnutrition infantile (enrichissement des bouillies de sevrage), là où la viande et le poisson restent coûteux.
\end{savaistu}

\begin{exoresolu}[Classer les aliments d'un repas camerounais]
Énoncé. Un élève déjeune d'une assiette composée de foutou de maïs, de ndolé (feuilles + arachides + viande) et d'un verre de jus d'avocat sucré. Pour chaque composant du repas, indique s'il s'agit d'un aliment simple ou composé, et précise le(s) groupe(s) de nutriments dominant(s).
\tcblower
Solution détaillée.
\begin{itemize}
\item \textbf{Foutou de maïs} : aliment classé parmi les \textbf{glucides} (nutriment dominant : l'amidon). Attention : il contient aussi de l'eau et des minéraux, donc il n'est « simple » que par son nutriment dominant.
\item \textbf{Ndolé} : \textbf{aliment composé} par excellence — les feuilles apportent vitamines et sels minéraux, les arachides des lipides et des protéines, la viande des protéines, l'huile de cuisson des lipides.
\item \textbf{Jus d'avocat sucré} : \textbf{aliment composé} — l'avocat est riche en lipides, le sucre ajouté apporte des glucides, l'eau du jus complète.
\end{itemize}
Conclusion : ce repas associe les trois grands groupes d'aliments simples (glucides, lipides, protéines) plus les vitamines et minéraux : c'est un repas varié, première condition d'une alimentation équilibrée que nous étudierons dans la suite de la leçon.
\end{exoresolu}

\begin{exercice}[À toi de jouer]
Classe les aliments suivants en « simple » (en précisant le groupe) ou « composé » : sucre de canne, œuf entier, huile de palme, eru, blanc d'œuf, lait caillé, bâton de manioc. Justifie chaque réponse en une phrase.
\end{exercice}

\begin{corrige}[Exercice : classification d'aliments courants]
Sucre de canne : simple (glucides, quasi pur). Œuf entier : composé (protéines + lipides + eau + minéraux). Huile de palme : simple (lipides quasi purs). Eru : composé (feuilles, huile, viande ou poisson). Blanc d'œuf : simple (protéines dominantes, albumine). Lait caillé : composé (eau, protéines, lipides, lactose, calcium). Bâton de manioc : classé glucidique par nutriment dominant, mais contient aussi eau et minéraux.
\end{corrige}

\coursec{Caractéristiques physicochimiques et tests d'identification}

Dans la section précédente, nous avons distingué les \cle{nutriments} (glucides, lipides, protéines, sels minéraux, vitamines, eau) des \cle{aliments} qui les contiennent. Le programme classe les aliments en \cle{aliments simples} (apportant essentiellement un seul type de nutriment : sucre, huile, blanc d'œuf, sel) et \cle{aliments composés} (apportant plusieurs nutriments : riz, gari, lait, arachide, avocat). Une question pratique se pose alors : \emph{comment identifier le ou les nutriments majeurs présents dans un aliment, simple ou composé ?} Au marché de Mokolo comme au laboratoire, l'aspect visuel ne suffit pas : une poudre blanche peut être de l'amidon, du sucre ou de la farine de manioc ; un liquide jaunâtre peut être de l'huile ou un jus de fruit sucré. La réponse scientifique repose sur des \cle{tests caractéristiques} : des réactions chimiques ciblées qui révèlent la présence d'un nutriment par un changement observable (coloration, précipité, tache). C'est l'objet de cette section : maîtriser les propriétés physicochimiques de chaque groupe de nutriments et la conduite rigoureuse de leur identification, savoir-faire exigible au TP et au Probatoire.

\begin{definition}[Test caractéristique]
Un \cle{test caractéristique} est une réaction chimique spécifique permettant d'identifier un nutriment (ou une substance) grâce à un changement observable : apparition d'une \cle{coloration}, formation d'un \cle{précipité}, ou toute autre modification visible à l'œil nu. Un test est dit \emph{positif} lorsque le changement attendu se produit (le nutriment est présent) et \emph{négatif} dans le cas contraire. La spécificité du test garantit qu'aucun autre nutriment courant ne donne le même résultat dans les mêmes conditions.
\end{definition}

\courssub{Les glucides : des « hydrates de carbone » aux propriétés variées}

\textbf{Situation concrète.} Coupe une pomme de terre crue et verse une goutte d'eau iodée dessus : la surface devient immédiatement bleu nuit. Fais la même chose sur un morceau de sucre en poudre : rien ne se passe, la solution reste jaune-brun. Pourtant, la pomme de terre et le sucre appartiennent tous deux à la famille des \cle{glucides}. Pourquoi des comportements si différents ?

\begin{definition}[Glucides]
Les \cle{glucides} (ou hydrates de carbone) sont des composés organiques de formule brute générale $\ce{C}_{n}\ce{(H2O)}_{m}$ : ils ne contiennent que du carbone, de l'hydrogène et de l'oxygène, ce dernier deux étant dans les mêmes proportions que dans l'eau. On distingue deux grandes classes selon la taille des molécules :
\begin{itemize}
\item les \cle{sucres simples} (oses ou monosaccharides), comme le \cle{glucose} $\ce{C6H12O6}$, le fructose ou le galactose : petites molécules, généralement \emph{solubles dans l'eau}, de saveur sucrée, et \cle{réducteurs} (propriété chimique mise en évidence par le test de Fehling) ;
\item les \cle{sucres complexes} (polyosides ou polysaccharides), comme l'\cle{amidon}, le glycogène ou la cellulose : macromolécules formées de très nombreuses unités d'oses liées entre elles, \emph{insolubles dans l'eau froide} (l'amidon forme une pâte ou un gel à chaud), sans saveur sucrée et \emph{non réducteurs} à l'état polymère.
\end{itemize}
\end{definition}

La \cle{solubilité} et le \cle{pouvoir réducteur} sont donc des propriétés physicochimiques discriminantes : le glucose se dissout facilement dans l'eau (c'est pourquoi le jus de goyave est sucré et limpide), tandis que l'amidon est insoluble à froid. Cette différence de structure (petite molécule libre \emph{vs} longue chaîne polymérique) explique pourquoi les tests d'identification sont différents.

\begin{propriete}[Tests d'identification des glucides]
\begin{itemize}
\item \cle{Test de l'amidon (glucide complexe)} : quelques gouttes d'\cle{eau iodée} (solution de diiode $\ce{I2}$ dans l'iodure de potassium $\ce{KI}$, dite solution de Lugol, jaune-brun) sur l'échantillon $\rightarrow$ coloration \cle{bleu-violet} (dite \emph{bleu nuit}) instantanée, due à la formation d'un complexe d'inclusion entre l'hélice de l'amylose et les ions triiodure $\ce{I3-}$. Le test se fait \emph{à froid}.
\item \cle{Test des sucres réducteurs} (glucose, fructose, maltose, lactose\ldots) : ajout de \cle{liqueur de Fehling} (solution alcaline de tartrate de cuivre(II) $\ce{Cu^2+}$, bleue) puis \cle{chauffage} au bain-marie (environ \SI{80}{\celsius}) $\rightarrow$ formation d'un \cle{précipité rouge brique} d'oxyde de cuivre(I) $\ce{Cu2O}$, insoluble, témoignant de la réduction du $\ce{Cu^2+}$ en $\ce{Cu+}$ par le groupe aldéhyde ou cétone libre du sucre. \textbf{Le chauffage est indispensable.}
\end{itemize}
\end{propriete}

\begin{experience}[Mise en évidence de l'amidon dans la pomme de terre (témoin positif) et absence dans le sucre (témoin négatif)]
\textbf{Matériel} : pomme de terre crue, sucre en poudre, couteau, eau iodée (Lugol), pipettes, lames de verre, deux tubes à essai.\par
\textbf{Protocole} :
\begin{enumerate}
\item Déposer une rondelle fine de pomme de terre sur une lame ; ajouter 2 gouttes d'eau iodée ; observer.
\item Mettre une pincée de sucre dans un tube, ajouter \SI{2}{mL} d'eau, agiter pour dissoudre, puis ajouter 2 gouttes d'eau iodée ; observer.
\end{enumerate}
\textbf{Observations} :
\begin{itemize}
\item Pomme de terre : coloration bleu nuit intense, localisée dans les cellules (grains d'amidon).
\item Sucre : la solution reste jaune-brun (couleur du réactif), pas de coloration bleu nuit.
\end{itemize}
\textbf{Interprétation} : La coloration bleu nuit est spécifique de l'amidon (polysaccharide de réserve). Le saccharose (sucre de table) est un disaccharide non réducteur et ne contient pas d'amidon ; il ne réagit pas à l'eau iodée.
\textbf{Conclusion} : La pomme de terre est un aliment riche en glucides complexes (amidon) ; le sucre est un aliment simple glucidique à base de sucres simples (saccharose), non détectable par ce test.
\end{experience}

\begin{attention}[Le chauffage, condition \emph{sine qua non} du test de Fehling]
Le test à la liqueur de Fehling \emph{exige un chauffage} au bain-marie (environ \SI{80}{\celsius} pendant 2 à 5 min). Sans chauffage, la réduction du $\ce{Cu^2+}$ en $\ce{Cu2O}$ ne se produit pas (cinétique nulle) : le précipité rouge brique n'apparaît \textbf{pas}, même si l'échantillon contient du glucose, et l'élève conclut à tort à un test négatif. C'est l'erreur la plus fréquente au TP et à l'examen. Autre piège : la liqueur de Fehling est bleue ; une solution qui reste bleue ou se trouble légèrement sans dépôt rouge brique après chauffage signifie un test \emph{négatif} (absence de sucres réducteurs), pas un test positif !
\end{attention}

\courssub{Les lipides : des corps gras hydrophobes}

\textbf{Situation concrète.} Verse \SI{5}{mL} d'huile de palme dans un verre contenant \SI{100}{mL} d'eau et agite vigoureusement pendant 30 secondes : de fines gouttelettes se dispersent, le mélange devient trouble (blanchâtre), mais dès que tu cesses d'agiter, l'huile remonte rapidement et reforme une couche distincte en surface. L'huile et l'eau ne se mélangent pas durablement. Cette observation banale de la cuisine révèle la propriété physique fondamentale des lipides.

\begin{definition}[Lipides]
Les \cle{lipides} (corps gras) sont des composés organiques formés essentiellement de carbone, d'hydrogène et d'oxygène (en proportion d'oxygène plus faible que dans les glucides, d'où un pouvoir calorifique plus élevé). Leur propriété physique caractéristique est leur \cle{insolubilité dans l'eau} : on dit qu'ils sont \cle{hydrophobes} (ou lipophiles). En revanche, ils sont \emph{solubles dans les solvants organiques apolaires} : éther, chloroforme, benzène, alcool chaud, acétone.
\end{definition}

Lorsqu'on agite fortement un mélange lipides + eau, on obtient une dispersion temporaire de gouttelettes appelée \cle{émulsion}. Le lait est une émulsion naturelle de globules de matière grasse dans l'eau, stabilisée par des phospholipides et des protéines (caséines).

\begin{propriete}[Tests d'identification des lipides]
\begin{itemize}
\item \cle{Test de la tache translucide (test au papier)} : écraser ou frotter l'échantillon sur du papier filtre, du papier journal ou une feuille de bloc-notes $\rightarrow$ une \cle{tache translucide} (qui laisse passer la lumière comme du papier huilé) \emph{persistante} après séchage (au moins 10 min à l'air libre) indique la présence de lipides. Une tache d'eau pure, elle, s'évapore complètement et le papier redevient opaque.
\item \cle{Test au Sudan III} : le \cle{Sudan III} est un colorant azoïque \cle{liposoluble} (soluble dans les graisses, insoluble dans l'eau) ; ajouté à l'échantillon (ou à l'émulsion), il se dissout sélectivement dans la phase lipidique et lui donne une coloration \cle{rouge-orangé} intense. C'est un test de confirmation très sensible.
\item \cle{Test d'émulsion} : agiter vigoureusement l'échantillon avec de l'eau dans un tube à essai $\rightarrow$ formation d'un mélange trouble laiteux (émulsion) qui se décolore lentement en présence de lipides. Moins spécifique que les deux précédents (certains détergents font de l'émulsion), il est utilisé en complément.
\end{itemize}
\end{propriete}

\begin{exemplebox}[L'arachide, un aliment riche en lipides]
Écrase une cacahuète grillée (arachide) entre deux feuilles de papier journal : une tache grasse translucide apparaît et ne disparaît pas au séchage (test au papier positif). Mise en présence de Sudan III, la pâte d'arachide se colore en rouge-orangé (test au Sudan III positif). Conclusion : l'arachide est riche en lipides — c'est d'ailleurs pour cela qu'on en extrait l'huile d'arachide, largement consommée au Cameroun pour la cuisson et la friture.
\end{exemplebox}

\courssub{Les protéines : les molécules azotées du vivant}

\textbf{Situation concrète.} Casse un œuf dans une poêle chaude : le blanc d'œuf, d'abord liquide, translucide et visqueux, devient blanc, opaque et solide en quelques secondes. Ce changement irréversible, que tout cuisinier connaît, est la \cle{coagulation} (ou dénaturation thermique) des protéines.

\begin{definition}[Protéines]
Les \cle{protéines} (ou protides) sont des macromolécules organiques contenant du carbone, de l'hydrogène, de l'oxygène, de l'\cle{azote} (N) et parfois du \cle{soufre} (S). La présence d'azote est leur signature chimique exclusive parmi les macronutriments énergétiques : on les appelle les \emph{matières azotées} de l'alimentation. Elles sont constituées de longues chaînes d'\cle{acides aminés} liés par des liaisons peptidiques.
\end{definition}

\begin{propriete}[Tests d'identification des protéines]
\begin{itemize}
\item \cle{Test au réactif de Biuret} : ajouter à l'échantillon aqueux (solution ou extrait) un excès de \cle{soude} (hydroxyde de sodium $\ce{NaOH}$, solution concentrée) puis quelques gouttes de \cle{sulfate de cuivre(II)} $\ce{CuSO4}$ (solution bleue) $\rightarrow$ apparition d'une coloration \cle{violette} (mauve) due à la formation d'un complexe de coordination entre les ions $\ce{Cu^2+}$ et les groupes peptidiques ($-\ce{CO-NH}-$) en milieu alcalin. Ce test se fait \emph{à froid}, sans chauffage. \textbf{Ordre impératif : soude d'abord, puis sulfate de cuivre goutte à goutte.}
\item \cle{Coagulation à la chaleur} : chauffées en milieu aqueux (au-dessus de \SI{60}{\celsius} généralement), les protéines se dénaturent (perte de leur structure tridimensionnelle native) et coagulent (agrégation en un réseau insoluble) : passage de l'état liquide/soluble à l'état solide/insoluble (blanc d'œuf, lait qui « tourne », cuisson de la viande).
\end{itemize}
\end{propriete}

\begin{attention}[Ne pas confondre les conditions des tests : le piège classique du Probatoire]
\begin{itemize}
\item \textbf{Fehling (sucres réducteurs)} : \textbf{Chauffage obligatoire} au bain-marie. Précipité rouge brique $\ce{Cu2O}$.
\item \textbf{Biuret (protéines)} : \textbf{À froid}. Coloration violette (complexe $\ce{Cu^2+}$-peptides). Chauffer le mélange soude + sulfate de cuivre est inutile, dégrade le réactif et peut fausser l'observation (précipitation de l'hydroxyde de cuivre $\ce{Cu(OH)2}$ bleu pâle).
\item \textbf{Eau iodée (amidon)} : \textbf{À froid}. Coloration bleu nuit instantanée.
\end{itemize}
Retiens aussi l'ordre d'ajout du réactif de Biuret : d'abord la soude (milieu alcalin nécessaire), ensuite le sulfate de cuivre goutte à goutte. L'inverse peut donner un précipité bleu d'hydroxyde de cuivre masquant la couleur violette.
\end{attention}

\courssub{Conduite pratique d'un test d'identification}

\begin{methode}[Réaliser et interpréter un test caractéristique]
\begin{enumerate}
\item \textbf{Préparer l'échantillon} : placer une petite quantité de l'aliment (solution claire, suspension, ou broyat fin) dans un tube à essai propre et sec. Pour un aliment solide (farine, viande, fruit), réaliser un extrait : broyer avec un peu d'eau distillée, filtrer si nécessaire, récupérer le filtrat.
\item \textbf{Ajouter le réactif} adapté au nutriment recherché (eau iodée, liqueur de Fehling, réactif de Biuret, Sudan III\ldots) en respectant les volumes usuels (\SI{1}{mL} à \SI{2}{mL} de réactif pour \SI{2}{mL} d'échantillon).
\item \textbf{Respecter les conditions} du test : chauffage au bain-marie pour la liqueur de Fehling ; simple agitation ou repos pour l'eau iodée, le Biuret, le Sudan III.
\item \textbf{Observer} attentivement le changement éventuel : coloration de la phase aqueuse ou de la phase grasse, apparition d'un précipité, aspect d'une tache sur papier.
\item \textbf{Conclure} rigoureusement : « Test \emph{positif} pour [nom du nutriment] » ou « Test \emph{négatif} pour [nom du nutriment] », en citant explicitement l'observation qui justifie la conclusion (ex : « coloration bleu nuit $\rightarrow$ présence d'amidon »).
\end{enumerate}
\end{methode}

\begin{exemplebox}[Un jus de goyave testé au laboratoire]
On place \SI{2}{mL} de jus de goyave filtré dans un tube à essai, on ajoute \SI{1}{mL} de liqueur de Fehling et on chauffe au bain-marie pendant 3 minutes. \textbf{Observation} : la solution bleue se trouble, puis un précipité rouge brique se dépose au fond du tube, le surnageant s'éclaircit. \textbf{Conclusion} : le test de Fehling est \emph{positif} pour les sucres réducteurs. Le jus de goyave contient donc des \textbf{glucides simples} (glucose, fructose), source d'énergie rapidement utilisable par l'organisme.
\end{exemplebox}

\begin{popfigure}
\begin{tikzpicture}
% Tube 1 : eau iodée
\draw[thick] (0,0) -- (0,2.6) (1,0) -- (1,2.6);
\draw[thick] (0,0) arc (180:360:0.5);
\fill[popblue!80] (0.02,0.02) -- (0.02,1.6) -- (0.98,1.6) -- (0.98,0.02) arc (180:360:0.48);
\draw (0.5,2.9) node[font=\small\bfseries,popdark]{Eau iodée (Lugol)};
\draw (0.5,-0.8) node[font=\small,popdark,align=center]{Amidon :\\ coloration bleu nuit};
% Tube 2 : Fehling
\draw[thick] (2.5,0) -- (2.5,2.6) (3.5,0) -- (3.5,2.6);
\draw[thick] (2.5,0) arc (180:360:0.5);
\fill[popblue!30] (2.52,0.02) -- (2.52,1.6) -- (3.48,1.6) -- (3.48,0.02) arc (180:360:0.48);
\fill[poporange!90] (2.52,0.02) -- (2.52,0.55) -- (3.48,0.55) -- (3.48,0.02) arc (180:360:0.48);
\foreach \x/\y in {2.7/1.0, 3.0/0.85, 3.3/1.25, 2.85/1.15, 3.15/0.95} \fill[poporange!90] (\x,\y) circle (0.04);
\draw (3,2.9) node[font=\small\bfseries,popdark]{Fehling + chauffage};
\draw (3,-0.8) node[font=\small,popdark,align=center]{Sucres réducteurs :\\ précipité rouge brique};
% Tube 3 : Biuret
\draw[thick] (5,0) -- (5,2.6) (6,0) -- (6,2.6);
\draw[thick] (5,0) arc (180:360:0.5);
\fill[poppurple!60] (5.02,0.02) -- (5.02,1.6) -- (5.98,1.6) -- (5.98,0.02) arc (180:360:0.48);
\draw (5.5,2.9) node[font=\small\bfseries,popdark]{Biuret (à froid)};
\draw (5.5,-0.8) node[font=\small,popdark,align=center]{Protéines :\\ coloration violette};
% Tube 4 : Sudan III
\draw[thick] (7.5,0) -- (7.5,2.6) (8.5,0) -- (8.5,2.6);
\draw[thick] (7.5,0) arc (180:360:0.5);
\fill[poporange!60] (7.52,0.02) -- (7.52,1.6) -- (8.48,1.6) -- (8.48,0.02) arc (180:360:0.48);
\draw (8,2.9) node[font=\small\bfseries,popdark]{Sudan III};
\draw (8,-0.8) node[font=\small,popdark,align=center]{Lipides :\\ coloration rouge-orangé};
\end{tikzpicture}
\legende{Résultats positifs des quatre tests caractéristiques fondamentaux : (1) Eau iodée $\rightarrow$ bleu nuit (amidon) ; (2) Fehling chauffée $\rightarrow$ précipité rouge brique (sucres réducteurs) ; (3) Biuret à froid $\rightarrow$ violet (protéines) ; (4) Sudan III $\rightarrow$ rouge-orangé (lipides).}
\end{popfigure}

\courssub{Identifier un aliment simple inconnu : démarche complète}

Face à un échantillon inconnu (par exemple une poudre blanche « X » ou un liquide « Y » au TP), le biologiste procède méthodiquement, comme un détective qui teste les hypothèses une à une sur des aliquotes séparées.

\begin{methode}[Identifier un aliment simple inconnu]
\begin{enumerate}
\item Répartir l'échantillon en \textbf{plusieurs tubes à essai} (un par test) : on ne mélange \emph{jamais} les réactifs entre eux dans le même tube (réactions croisées, confusion des couleurs).
\item Réaliser le \textbf{test à l'eau iodée} sur le premier aliquot (recherche amidon).
\item Réaliser le \textbf{test à la liqueur de Fehling avec chauffage} sur le deuxième aliquot (recherche sucres réducteurs).
\item Réaliser le \textbf{test au réactif de Biuret} sur le troisième aliquot (recherche protéines).
\item Réaliser le \textbf{test de la tache translucide} (sur papier) et/ou \textbf{Sudan III} sur le quatrième aliquot (recherche lipides).
\item Conclure d'après les résultats : un seul test positif $\rightarrow$ aliment simple (glucidique, lipidique ou protidique) ; plusieurs tests positifs $\rightarrow$ aliment composé.
\end{enumerate}
\end{methode}

\begin{popfigure}
\begin{tikzpicture}[
  grow=down,
  level distance=1.8cm,
  sibling distance=3.5cm,
  level 1/.style={sibling distance=4.5cm},
  level 2/.style={sibling distance=3.0cm},
  level 3/.style={sibling distance=2.5cm},
  every node/.style={font=\small,align=center,inner sep=3pt},
  test/.style={rectangle,rounded corners,draw=popblue,fill=popblueL,thick},
  res/.style={rectangle,rounded corners,draw=popgreen,fill=popgreenL,thick},
  neg/.style={rectangle,rounded corners,draw=popmuted,fill=popcream},
  edge from parent/.style={draw=popdark, thick, ->},
  lbl/.style={font=\scriptsize\itshape,popdark,midway,sloped}
]
\node[test]{Eau iodée ?}
  child { node[res, align=center]{Amidon\\ (bleu nuit)} edge from parent node[lbl,left]{positif} }
  child { node[test]{Fehling + chauffage ?} edge from parent node[lbl,right]{négatif}
    child { node[res, align=center]{Sucres réducteurs\\ (rouge brique)} edge from parent node[lbl,left]{positif} }
    child { node[test]{Biuret ?} edge from parent node[lbl,right]{négatif}
      child { node[res, align=center]{Protéines\\ (violet)} edge from parent node[lbl,left]{positif} }
      child { node[test]{Sudan III / Papier ?} edge from parent node[lbl,right]{négatif}
        child { node[res, align=center]{Lipides\\ (rouge-orangé / tache)} edge from parent node[lbl,left]{positif} }
        child { node[neg, align=center]{Aucun nutriment\\ détecté par ces tests} edge from parent node[lbl,right]{négatif} }
      }
    }
  };
\end{tikzpicture}
\legende{Arbre de décision pour l'identification d'un aliment simple inconnu. Les tests sont réalisés successivement sur des aliquotes séparées. Un résultat positif permet d'identifier le groupe nutritionnel dominant ; on poursuit les tests pour vérifier l'absence d'autres nutriments (aliment simple pur) ou leur présence (aliment composé).}
\end{popfigure}

\begin{exoresolu}[Identifier le nutriment dominant dans trois échantillons]
Au laboratoire du lycée, on dispose de trois échantillons anonymés : \textbf{A} (jus d'orange), \textbf{B} (blanc d'œuf dilué), \textbf{C} (bouillie de maïs). On obtient les résultats suivants :
\begin{itemize}
\item A + liqueur de Fehling chauffée $\rightarrow$ précipité rouge brique.
\item B + soude puis sulfate de cuivre $\rightarrow$ coloration violette.
\item C + eau iodée $\rightarrow$ coloration bleu nuit.
\end{itemize}
Identifier le nutriment dominant de chaque échantillon.
\tcblower
\textbf{Échantillon A (jus d'orange)} : Le précipité rouge brique obtenu avec la liqueur de Fehling \emph{chauffée} est le résultat positif spécifique du test des sucres réducteurs. Le jus d'orange contient donc des \textbf{glucides simples} (principalement glucose, fructose, saccharose hydrolysé). C'est un aliment simple glucidique.

\textbf{Échantillon B (blanc d'œuf)} : La coloration violette avec le réactif de Biuret (soude + sulfate de cuivre, à froid) est caractéristique des \textbf{protéines}. Le blanc d'œuf est effectivement très riche en ovalbumine, une protéine globulaire. C'est un aliment simple protidique.

\textbf{Échantillon C (bouillie de maïs)} : La coloration bleu nuit avec l'eau iodée révèle la présence d'\textbf{amidon}, un glucide complexe (polysaccharide de réserve), substance abondante dans l'endosperme du grain de maïs. C'est un aliment simple glucidique (amylacé).

\textbf{Conclusion} : A est riche en glucides simples (sucres réducteurs), B en protéines, C en glucides complexes (amidon).
\end{exoresolu}

\courssub{Tableau récapitulatif des tests caractéristiques (Mémento Probatoire)}

\begin{center}
\rowcolors{1}{popcream}{white}
\begin{tabularx}{\linewidth}{|l|X|X|X|}
\hline
\rowcolor{popblueL} \textbf{Nutriment recherché} & \textbf{Réactif(s)} & \textbf{Conditions opératoires} & \textbf{Résultat positif (Observation)} \\
\hline
Amidon (glucide complexe) & Eau iodée (Lugol) & À froid, sans chauffage & Coloration \textbf{bleu-violet (bleu nuit)} instantanée \\
\hline
Sucres réducteurs (glucose, fructose, maltose, lactose) & Liqueur de Fehling (tartrate de $\ce{Cu^2+}$ alcalin) & \textbf{Chauffage au bain-marie} (\SI{80}{\celsius}, 2--5 min) & Précipité \textbf{rouge brique} ($\ce{Cu2O}$) au fond du tube \\
\hline
Protéines (liaisons peptidiques) & Réactif de Biuret ($\ce{NaOH}$ + $\ce{CuSO4}$) & \textbf{À froid}, soude \emph{puis} sulfate de cuivre gtte à gtte & Coloration \textbf{violette (mauve)} de la solution \\
\hline
Lipides (corps gras) & Papier filtre / journal ; Sudan III (colorant liposoluble) & À froid, séchage pour le papier & \textbf{Tache translucide persistante} (papier) ; \textbf{Coloration rouge-orangé} (Sudan III) \\
\hline
\end{tabularx}
\end{center}

\begin{aretenir}[L'essentiel de la section]
\begin{itemize}
\item Les glucides ($\ce{C}_{n}\ce{(H2O)}_{m}$) se divisent en sucres simples (solubles, sucrés, réducteurs) et sucres complexes comme l'amidon (insoluble à froid, non réducteur, substance de réserve).
\item Les lipides sont hydrophobes (insolubles dans l'eau, solvants organiques), forment des émulsions instables avec l'eau.
\item Les protéines sont les macromolécules azotées (C, H, O, N, S) ; elles coagulent (dénaturation) par la chaleur.
\item \textbf{Quatre tests de référence} : Eau iodée $\xrightarrow{\text{à froid}}$ bleu nuit (amidon) ; Fehling $\xrightarrow{\text{chauffage}}$ rouge brique (sucres réducteurs) ; Biuret $\xrightarrow{\text{à froid}}$ violet (protéines) ; Papier/Sudan III $\xrightarrow{\text{à froid}}$ tache translucide / rouge-orangé (lipides).
\item \textbf{Rigueur expérimentale} : un test par tube, respect strict des conditions (surtout chauffage Fehling), conclusion explicite citant le nutriment et l'observation.
\end{itemize}
\end{aretenir}

\begin{savaistu}[L'eau iodée, gardienne de la qualité des aliments]
La coloration bleu nuit de l'amidon avec l'iode est si spécifique qu'elle sert d'outil de contrôle de la fraude alimentaire. Un miel ou un sirop de sucre frauduleusement « coupé » avec de la fécule de maïs ou de la farine (adultération pour augmenter le volume) se trahit immédiatement : quelques gouttes d'eau iodée le font virer au bleu nuit, alors qu'un produit pur reste jaune-brun. De même, les services de contrôle de qualité (MINCOMMERCE, laboratoires d'analyse) utilisent ce test pour vérifier la pureté des amidons alimentaires et pharmaceutiques. Une simple réaction de couleur, découverte au XIXe siècle, protège encore aujourd'hui le consommateur camerounais !
\end{savaistu}

\begin{exercice}[À toi de jouer]
On veut comparer la composition de deux aliments locaux : une graine de \textbf{niébé} (haricot) et une tranche d'\textbf{avocat} mûr. Propose, pour chacun, les tests à réaliser et les résultats attendus, sachant que le niébé est réputé riche en protéines et en amidon, et l'avocat riche en lipides. Rédige ta réponse sous forme d'un tableau à trois colonnes : Aliment / Test(s) réalisé(s) / Résultat(s) attendu(s) et Conclusion.
\end{exercice}

\begin{corrige}[Exercice : Niébé vs Avocat]
\begin{center}
\begin{tabularx}{\linewidth}{|l|X|X|}
\hline
\rowcolor{popblueL} \textbf{Aliment} & \textbf{Test(s) réalisé(s)} & \textbf{Résultat(s) attendu(s) et Conclusion} \\
\hline
\textbf{Niébé (graines)} & 1. Eau iodée (sur broyat) & Coloration \textbf{bleu nuit} $\rightarrow$ \textbf{Amidon présent} (réserve du cotylédon). \\
& 2. Réactif de Biuret (sur extrait aqueux) & Coloration \textbf{violette} $\rightarrow$ \textbf{Protéines présentes} (protéines de réserve). \\
& 3. Fehling (sur extrait) & Négatif (bleu) $\rightarrow$ Pas de sucres réducteurs libres. \\
& 4. Sudan III / Papier & Négatif ou trace $\rightarrow$ Peu de lipides. \\
\hline
\textbf{Avocat (pulpe)} & 1. Papier (écrasement) / Sudan III & \textbf{Tache translucide persistante} / \textbf{Rouge-orangé} $\rightarrow$ \textbf{Lipides abondants} (15--20\%). \\
& 2. Eau iodée & Négatif (jaune-brun) $\rightarrow$ Pas d'amidon. \\
& 3. Biuret & Faible ou négatif $\rightarrow$ Peu de protéines (environ 2\%). \\
& 4. Fehling & Négatif $\rightarrow$ Pas de sucres réducteurs. \\
\hline
\end{tabularx}
\end{center}
\textbf{Note} : Le niébé est un \emph{aliment composé} (protidique + glucidique amylacé) ; l'avocat est un \emph{aliment simple lipidique} (fruit gras). Chaque test doit être réalisé sur un aliquot séparé, et la conclusion doit citer explicitement le nutriment identifié.
\end{corrige}

\coursec{Rôle des aliments et alimentation équilibrée}

Dans la section précédente, tu as appris à reconnaître les grands groupes d'aliments simples grâce à leurs caractéristiques physicochimiques et à des tests d'identification (liqueur de Fehling pour les sucres réducteurs, réactif de Biuret pour les protéines, teinture d'iode pour l'amidon, émulsion pour les lipides). Une question surgit alors naturellement : \emph{à quoi servent tous ces nutriments dans notre corps ?} Pourquoi un enfant nourri uniquement de bouillie de manioc, qui mange pourtant « à sa faim », présente-t-il un ventre ballonné et des cheveux décolorés ? Pourquoi une femme enceinte doit-elle manger différemment d'un adolescent ? C'est à ces questions que répond cette section : comprendre les \cle{rôles des aliments}, puis définir et construire une \cle{alimentation équilibrée}.

\courssub{Les trois rôles des aliments}

\textbf{Situation de départ : une observation au village.}\par

Dans un village de la région du Centre, deux enfants du même âge ont des allures très différentes. Le premier, actif et vigoureux, mange chaque jour du ndolé avec du miondo, des arachides, des fruits et parfois du poisson. Le second, apathique, avec les cheveux roux et un abdomen gonflé, ne mange presque que de la bouillie de maïs. Pourtant, les deux mangent en quantité suffisante. La différence ne vient donc pas de la \emph{quantité} mais de la \emph{qualité} : les aliments jouent des rôles différents, et aucun ne peut les assurer tous seul.

\begin{experience}[Observation et hypothèse : pourquoi varier les aliments ?]
\textbf{Observation :} les enfants nourris avec un seul aliment énergétique (bouillie de manioc ou de maïs) présentent des signes de malnutrition malgré un apport calorique suffisant.

\textbf{Hypothèse :} les aliments n'ont pas tous le même rôle ; certains construisent le corps, d'autres fournissent l'énergie, d'autres encore protègent l'organisme.

\textbf{Vérification :} on analyse la composition des aliments consommés. La bouillie de manioc contient presque exclusivement des glucides (amidon mis en évidence par la teinture d'iode), mais pratiquement pas de protéines (test de Biuret négatif). Or les protéines sont les matériaux de construction des tissus.

\textbf{Conclusion :} il faut classer les aliments selon leur rôle physiologique. On distingue trois grands rôles : \cle{énergétique}, \cle{plastique} et \cle{fonctionnel}.
\end{experience}

\textbf{Le rôle énergétique.}\par

Notre corps fonctionne comme un moteur : il lui faut du « carburant » en permanence, même au repos (battements du cœur, respiration, maintien de la température à 37 °C). Ce carburant est fourni principalement par les \cle{glucides} et les \cle{lipides}. Dans la cellule, la respiration cellulaire libère l'énergie contenue dans le glucose sous forme d'ATP :

\[ \ce{C6H12O6 + 6O2 -> 6CO2 + 6H2O + \text{énergie (ATP)}} \]

Les glucides (riz, maïs, manioc, igname, plantain, pain) sont la source d'énergie la plus rapidement utilisable ; les lipides (huile de palme, huile d'arachide, beurre, noix de palme) constituent une réserve énergétique concentrée : 1 g de lipides fournit environ 9 kcal, contre 4 kcal pour 1 g de glucides (ou de protéines).

\textbf{Le rôle plastique (de construction et de réparation).}\par

Le corps se construit et se répare sans cesse : croissance de l'enfant, renouvellement de la peau, des globules rouges, cicatrisation des plaies. Les matériaux de construction sont les \cle{protéines} (viande, poisson, œufs, haricots niébé, arachides, soja), l'\cle{eau} (environ 65 \% de la masse corporelle, constituant structural des cellules) et certains \cle{sels minéraux} : le \cle{calcium} et le \cle{phosphore} pour la minéralisation des os et des dents, le \cle{fer} pour l'hémoglobine des globules rouges.

\textbf{Le rôle fonctionnel (protecteur et régulateur).}\par

Pour que la « machine » tourne bien, il faut des molécules régulatrices, nécessaires en très petites quantités mais absolument indispensables : les \cle{vitamines} (A, B, C, D, E, K...), les \cle{sels minéraux} (iode, fer, zinc, fluor...), les \cle{fibres} alimentaires (qui facilitent le transit intestinal et nourrissent le microbiote) et l'\cle{eau} (milieu de toutes les réactions chimiques, transport des nutriments, élimination des déchets par l'urine, thermorégulation). Ces éléments protègent contre les maladies : la vitamine C contre le scorbut, l'iode contre le goitre, le fer contre l'anémie, la vitamine A contre la cécité nocturne.

\begin{popfigure}
\begin{center}
\begin{tikzpicture}[grow=right, level distance=4.2cm,
  every node/.style={font=\small},
  level 1/.style={sibling distance=3.2cm},
  edge from parent/.style={draw=popmuted,thick}]
\node[draw=popdark, fill=popcream, rounded corners, thick, align=center, font=\bfseries] {Rôles des\\ aliments}
  child { node[draw=popgreen, fill=popgreenL, rounded corners, align=center, text width=4.6cm]
    {\textbf{Rôle fonctionnel}\\ Vitamines, sels minéraux,\\ fibres, eau\\ \emph{protègent et régulent}}
    edge from parent[draw=popgreen, thick] }
  child { node[draw=popblue, fill=popblueL, rounded corners, align=center, text width=4.6cm]
    {\textbf{Rôle plastique}\\ Protéines, eau,\\ calcium, phosphore\\ \emph{construisent et réparent}}
    edge from parent[draw=popblue, thick] }
  child { node[draw=poporange, fill=poporangeL, rounded corners, align=center, text width=4.6cm]
    {\textbf{Rôle énergétique}\\ Glucides, lipides\\ \emph{fournissent l'énergie}}
    edge from parent[draw=poporange, thick] };
\end{tikzpicture}
\end{center}
\legende{Les trois rôles des aliments et les nutriments qui les assurent. Un même aliment peut avoir plusieurs rôles : le poisson est plastique (protéines) mais aussi fonctionnel (iode, vitamine D, oméga-3).}
\end{popfigure}

\begin{attention}[Un aliment, plusieurs rôles]
Ne classe pas un aliment dans un seul rôle de façon rigide. Le lait, par exemple, est à la fois plastique (protéines, calcium), énergétique (lactose, lipides) et fonctionnel (vitamines A, D, B12, eau). Au Probatoire, précise toujours le \emph{nutriment} qui justifie le rôle : « les haricots niébé ont un rôle plastique \emph{grâce à leurs protéines} ».
\end{attention}

\courssub{L'alimentation équilibrée : définition et conditions}

\begin{definition}[Alimentation équilibrée]
Une \cle{alimentation équilibrée} est une alimentation \textbf{variée} qui apporte \textbf{tous les nutriments} (glucides, lipides, protéines, vitamines, sels minéraux, fibres, eau) en \textbf{quantités adaptées} aux besoins de l'individu, \textbf{sans excès ni carence}. Elle couvre les besoins de l'organisme en \emph{quantité} (assez d'énergie) et en \emph{qualité} (tous les nutriments représentés).
\end{definition}

\begin{propriete}[Principe de la variété alimentaire (admis)]
Aucun aliment pris isolément ne contient tous les nutriments nécessaires à l'organisme dans des proportions adéquates. Seule la \textbf{variété} des aliments dans les repas assure l'équilibre alimentaire. C'est pourquoi un repas équilibré associe typiquement un aliment énergétique (céréale ou tubercule), un aliment plastique (viande, poisson, légumineuse) et des aliments fonctionnels (légumes, fruits), accompagnés d'eau.
\end{propriete}

\textbf{Les conditions de l'équilibre alimentaire.}\par

Les besoins ne sont pas les mêmes pour tous : l'équilibre alimentaire est \emph{relatif à chaque individu}. Il dépend de plusieurs facteurs que le programme regroupe en \textbf{cinq déterminants} :

\begin{itemize}
  \item \textbf{L'âge :} le nourrisson se nourrit de lait (maternel idéalement jusqu'à 6 mois) ; l'enfant et l'adolescent ont des besoins élevés en protéines, calcium, phosphore et fer pour la croissance ; l'adulte a des besoins d'entretien ; le vieillard a des besoins énergétiques réduits mais doit maintenir des apports en protéines (lutte contre la sarcopénie) et en calcium (prévention de l'ostéoporose).
  \item \textbf{L'état physiologique :} la grossesse et l'allaitement augmentent fortement les besoins en protéines, fer, acide folique (vitamine B9), calcium et iode.
  \item \textbf{La masse corporelle :} plus le corps est lourd, plus le métabolisme de base et les dépenses de déplacement sont importants.
  \item \textbf{L'activité physique :} un bûcheron, un paysan ou un footballeur dépense bien plus d'énergie qu'un employé de bureau (activité sédentaire).
  \item \textbf{L'état de santé :} la maladie modifie les besoins (un diabétique doit limiter les sucres rapides et privilégier l'index glycémique bas ; un hypertendu doit limiter le chlorure de sodium ; une insuffisance rénale impose de limiter les protéines et le potassium).
\end{itemize}

\begin{exemplebox}[La femme enceinte : des besoins accrus en qualité]
Une femme enceinte « mange pour deux » en qualité plus qu'en quantité. Ses besoins augmentent en \textbf{fer} (expansion du volume sanguin maternel, constitution des réserves fœtales, prévention de l'anémie), en \textbf{calcium} et \textbf{phosphore} (minéralisation du squelette fœtal), en \textbf{protéines} (croissance des tissus fœtaux et placentaires) et en \textbf{acide folique} (prévention des anomalies de fermeture du tube neural). C'est pourquoi les centres de santé du Cameroun distribuent systématiquement des suppléments de fer et d'acide folique aux femmes enceintes lors des consultations prénatales (CPN). Un menu adapté associe par exemple : ignames pilées (énergie), sauce feuilles au poisson fumé (protéines, fer hémique, calcium, iode) et un fruit riche en vitamine C (mangue, orange, goyave) qui améliore l'absorption du fer non hémique.
\end{exemplebox}

\courssub{Les conséquences d'une alimentation déséquilibrée}

\begin{savaistu}[Le kwashiorkor, maladie du sevrage précoce]
Le mot \emph{kwashiorkor} vient de la langue Ga (Ghana) et signifie « la maladie de l'enfant délaissé » (lorsque le suivant naît et que l'aîné est sevré brutalement). Il frappe les enfants sevrés nourris presque exclusivement de bouillie de maïs ou de manioc : les signes cliniques sont la \textbf{dépigmentation des cheveux} (cheveux roux, fins), le \textbf{ventre ballonné} (ascite par hypoalbuminémie : le foie ne synthétise plus assez d'albumine, la pression oncotique chute, l'eau passe dans la cavité péritonéale), un visage bouffi (œdèmes sous-cutanés), une dermatose en « peinture écaillée » et une grande fatigue avec fonte musculaire. C'est une carence en \textbf{protéines} (et acides aminés essentiels) malgré un apport énergétique souvent suffisant. Le \textbf{marasme}, lui, est une carence globale (énergie \emph{et} protéines) : l'enfant est extrêmement maigre, « peau et os », sans œdèmes. Ces deux formes relèvent de la \textbf{malnutrition aiguë sévère}, encore présente dans certaines zones du Cameroun (Extrême-Nord, Nord, Adamaoua) et prise en charge par les programmes de nutrition thérapeutique (aliments thérapeutiques prêts à l'emploi -- ATPE).
\end{savaistu}

On peut classer les déséquilibres nutritionnels en deux grandes catégories :

\begin{itemize}
  \item \textbf{Les carences (apports insuffisants par rapport aux besoins) :}
  \begin{itemize}
    \item \textbf{Malnutrition protéino-énergétique :} kwashiorkor, marasme.
    \item \textbf{Carences en micronutriments :} anémie ferriprive (manque de fer -- pâleur, fatigue, essoufflement), goitre (manque d'iode -- hypertrophie de la thyroïde, fréquent jadis dans les zones montagneuses de l'Ouest et de l'Adamaoua ; d'où l'iodation obligatoire du sel de cuisine au Cameroun depuis 1995), rachitisme (manque de vitamine D et/ou calcium -- déformations osseuses), cécité nocturne et xérophtalmie (manque de vitamine A), scorbut (manque de vitamine C).
  \end{itemize}
  \item \textbf{Les excès (apports trop importants par rapport aux dépenses) :} obésité (excès de glucides et lipides, facteur de risque majeur), diabète de type 2 (résistance à l'insuline), hypertension artérielle (excès de sel/sodium, excès pondéral), dyslipidémies et maladies cardiovasculaires (excès d'acides gras saturés et trans), caries dentaires (excès de sucres fermentescibles, grignotages).
\end{itemize}

\begin{attention}[« Manger beaucoup » n'est pas « bien manger »]
Piège classique au Probatoire ! Un régime uniquement à base de manioc, de foufou de maïs ou de riz blanc \textbf{couvre les besoins énergétiques} mais provoque des \textbf{carences en protéines} (kwashiorkor chez l'enfant), en vitamines et en sels minéraux. L'équilibre alimentaire concerne la \emph{qualité} (variété, densité nutritionnelle) autant que la \emph{quantité}. Inversement, manger trop, même d'aliments variés, conduit à l'obésité et aux maladies chroniques non transmissibles.
\end{attention}

\courssub{La pyramide alimentaire : un guide visuel}

Pour aider chacun à composer ses repas, les nutritionnistes représentent l'équilibre alimentaire par une \cle{pyramide alimentaire} : à la base, les aliments à consommer en grande quantité et à chaque repas ; au sommet, ceux à limiter fortement.

\begin{popfigure}
\begin{center}
\begin{tikzpicture}[font=\small]
% Couches de la pyramide (de bas en haut)
\fill[poporangeL]  (-3.5,0) -- (3.5,0) -- (2.5,1.4) -- (-2.5,1.4) -- cycle;
\fill[popgreenL]   (-2.5,1.4) -- (2.5,1.4) -- (1.6,2.7) -- (-1.6,2.7) -- cycle;
\fill[popblueL]    (-1.6,2.7) -- (1.6,2.7) -- (0.8,3.9) -- (-0.8,3.9) -- cycle;
\fill[poppinkL]    (-0.8,3.9) -- (0.8,3.9) -- (0,5.1) -- cycle;
% Contour
\draw[popdark, thick] (-3.5,0) -- (0,5.1) -- (3.5,0) -- cycle;
% Étiquettes internes
\node[align=center] at (0,0.7) {\textbf{Céréales, tubercules} (riz, maïs, manioc, igname, plantain)};
\node[align=center] at (0,2.05) {\textbf{Fruits et légumes} (ndolé, folong, gombo, mangue, papaye, banane)};
\node[align=center] at (0,3.3) {\textbf{Protéines et laitages} (poisson, viande, œufs, légumineuses, lait)};
\node[align=center] at (0,4.35) {\footnotesize \textbf{Graisses ajoutées, sucres, sel}};
% Annotations latérales droites
\draw[popmuted, thin] (3.5,0.7) -- (4.8,0.7) node[right, align=left, text width=4.5cm] {\footnotesize \textbf{Base} : à chaque repas, en grande quantité. \textbf{Eau} : 1,5--2 L/jour (davantage si chaleur).};
\draw[popmuted, thin] (2.5,2.05) -- (4.8,2.05) node[right, align=left, text width=4.5cm] {\footnotesize \textbf{2ème étage} : plusieurs portions par jour (5 portions fruits/légumes). Rôle fonctionnel.};
\draw[popmuted, thin] (1.6,3.3) -- (4.8,3.3) node[right, align=left, text width=4.5cm] {\footnotesize \textbf{3ème étage} : 1 à 2 portions par jour. Rôle plastique. Privilégier poisson, légumineuses.};
\draw[popmuted, thin] (0.8,4.5) -- (4.8,4.5) node[right, align=left, text width=4.5cm] {\footnotesize \textbf{Sommet} : à \textbf{limiter} (huiles, beurre, sucreries, sodas, sel, cubes).};
% Annotations latérales gauches (rôles)
\draw[popmuted, thin] (-3.5,0.7) -- (-4.8,0.7) node[left, align=right] {\footnotesize Rôle énergétique};
\draw[popmuted, thin] (-2.5,2.05) -- (-4.8,2.05) node[left, align=right] {\footnotesize Rôle fonctionnel};
\draw[popmuted, thin] (-1.6,3.3) -- (-4.8,3.3) node[left, align=right] {\footnotesize Rôle plastique};
\end{tikzpicture}
\end{center}
\legende{La pyramide alimentaire adaptée au contexte camerounais. La base (céréales/tubercules + eau) fournit l'énergie ; le 2ème étage (fruits/légumes) protège ; le 3ème étage (protéines/laitages) construit ; le sommet (graisses ajoutées, sucres, sel) doit être limité.}
\end{popfigure}

La notion de \cle{portion} est un repère pratique : une portion de féculents = un poing fermé (environ 100--150 g cuit) ; une portion de protéines = la paume de la main (environ 100 g poisson/viande ou 2 œufs ou 150 g légumineuses cuites) ; une portion de fruit = une mangue moyenne, une orange, une banane plantain ; une portion de légumes = deux mains jointes. Un adulte devrait boire au moins 1,5 à 2 litres d'eau par jour, davantage sous climat chaud ou lors d'effort physique.

\begin{exoresolu}[Analyser un menu et le corriger]
\textbf{Énoncé.} Un élève de Première déjeune chaque midi d'un grand plat de bâton de manioc (miondo) avec de l'huile de palme rouge et du sel. 1) Ce repas est-il équilibré ? Justifier en identifiant les rôles couverts et ceux manquants. 2) Quel trouble nutritionnel risque-t-il à long terme ? 3) Proposer un menu corrigé adapté à un adolescent (besoins élevés : croissance, activité scolaire), en utilisant des aliments locaux accessibles.

\tcblower

\textbf{Solution.}

\textbf{1) Analyse du repas.}
\begin{itemize}
  \item \textbf{Bâton de manioc :} source majeure de \textbf{glucides} (amidon) $\rightarrow$ \textbf{Rôle énergétique} couvert.
  \item \textbf{Huile de palme rouge :} source de \textbf{lipides} $\rightarrow$ \textbf{Rôle énergétique} (renforcé) + apport en \textbf{vitamine A} (caroténoïdes) $\rightarrow$ début de \textbf{Rôle fonctionnel}.
  \item \textbf{Sel :} apport en \textbf{sodium} et \textbf{iode} (si sel iodé) $\rightarrow$ \textbf{Rôle fonctionnel} (partiel).
\end{itemize}
\textbf{Manquants :} \textbf{Protéines} (pas de rôle plastique), \textbf{Fibres} (transit), \textbf{Vitamines C, B, E, K} (pas de légumes/fruits frais), \textbf{Calcium, Fer} (biodisponibles). Le repas n'est \textbf{pas équilibré} : il est déséquilibré qualitativement (monotonie, absence de groupe plastique et fonctionnel complet).

\textbf{2) Risque à long terme.}
Chez un adolescent en pleine croissance : \textbf{retard de croissance staturo-pondéral}, \textbf{anémie ferriprive} (fatigue, baisse des performances scolaires), \textbf{carence en vitamine A} (cécité crépusculaire, infections respiratoires/digestives répétées), risque de \textbf{kwashiorkor} si l'apport protéique reste très bas. L'excès d'huile de palme (acides gras saturés) sans activité physique intense favorise aussi le surpoids futur.

\textbf{3) Menu corrigé pour un adolescent (équilibré, varié, local) :}
\begin{itemize}
  \item \textbf{Énergétique (Base) :} Bâton de manioc (miondo) \emph{ou} riz blanc \emph{ou} igname pilée \emph{ou} plantain bouilli.
  \item \textbf{Plastique (Construction) :} Poisson braisé (carpe, capitaine) \emph{ou} sauce d'arachides avec morceaux de viande/abats \emph{ou} haricots niébé (haricots blancs) + œuf dur. \textit{Justification : protéines de haute valeur biologique, fer hémique (poisson/viande), calcium (arêtes de poisson mangées, légumineuses).}
  \item \textbf{Fonctionnel (Protection) :} Ndolé (feuilles de vernonia) \emph{ou} folong (feuilles de gnetum) \emph{ou} sauce gombo \emph{ou} salade de tomates/oignons/concombre. \textit{Justification : fibres, vitamines (A, C, B9), fer non hémique, calcium, antioxydants.}
  \item \textbf{Dessert :} Un fruit de saison (mangue, papaye, ananas, goyave, banane douce). \textit{Justification : vitamine C (absorption fer), fibres, eau.}
  \item \textbf{Boisson :} Eau potable à volonté (pas de soda, pas de jus sucré industriel).
\end{itemize}
Ce menu associe explicitement les trois rôles des aliments : il est varié, dense en micronutriments, adapté aux goûts locaux et aux besoins de l'adolescent.
\end{exoresolu}

\begin{exercice}[À toi de jouer]
Une grand-mère de 70 ans, peu active (retraitée, vie sédentaire), souffre d'hypertension artérielle (HTA) traitée et d'ostéoporose (fragilité osseuse). Indique \textbf{deux} conseils alimentaires précis, adaptés à son état de santé et à son âge, en justifiant \textbf{chacun} par le rôle des aliments ou nutriments concernés.
\end{exercice}

\begin{corrige}[Exercice : la grand-mère hypertendue et ostéoporotique]
\begin{itemize}
  \item \textbf{Conseil 1 : Limiter strictement le sel ajouté et les cubes d'assaisonnement industriels (type Maggi, Jumbo), riches en chlorure de sodium et glutamate.}
  \textit{Justification :} Le sodium (rôle fonctionnel : pression osmotique, influx nerveux) devient délétère en excès : il favorise la rétention d'eau et l'augmentation de la pression artérielle, aggravant l'HTA et le risque d'AVC. L'équilibre, c'est « sans excès ni carence ». Remplacer le sel par les épices locales (poivre de Penja, gingembre, ail, oignon, persil, céleri) pour le goût.
  \item \textbf{Conseil 2 : Maintenir, voire augmenter, les apports en calcium et en protéines de bonne qualité, malgré l'appétit réduit et les faibles dépenses énergétiques.}
  \textit{Justification :} Le calcium et le phosphore (rôle plastique : minéralisation osseuse) et les protéines (rôle plastique : matrice osseuse, masse musculaire) sont indispensables pour freiner l'ostéoporose et la sarcopénie (fonte musculaire du vieillard). Sources locales : lait caillé (kindirmo), yaourt, fromage local (wara), sardines à l'huile (avec arêtes), légumes à feuilles sombres (ndolé, folong), haricots niébé. Réduire en revanche les graisses saturées (huile de palme, beurre) et les sucres rapides (sommet de la pyramide) pour éviter la prise de poids inutile.
\end{itemize}
\end{corrige}

\begin{aretenir}[L'essentiel de la section]
\begin{itemize}
  \item Les aliments ont trois rôles physiologiques : \textbf{énergétique} (glucides, lipides $\rightarrow$ ATP), \textbf{plastique} (protéines, eau, Ca, P, Fe $\rightarrow$ structure, croissance, réparation) et \textbf{fonctionnel} (vitamines, sels minéraux, fibres, eau $\rightarrow$ régulation, protection, transit).
  \item Une \textbf{alimentation équilibrée} = variée + quantités adaptées + sans excès ni carence. Aucun aliment seul ne suffit (principe de variété).
  \item Les besoins (donc l'équilibre) sont \textbf{individuels} : dépendent de l'\textbf{âge}, de l'\textbf{état physiologique} (grossesse, allaitement), de la \textbf{masse corporelle}, de l'\textbf{activité physique} et de l'\textbf{état de santé}.
  \item Les déséquilibres provoquent des \textbf{carences} (kwashiorkor, marasme, anémie, goitre, rachitisme, cécité vitamine A) ou des \textbf{excès} (obésité, diabète type 2, HTA, maladies cardiovasculaires).
  \item La \textbf{pyramide alimentaire} guide les choix quotidiens : base = céréales/tubercules + eau ; étage 1 = fruits/légumes ; étage 2 = protéines/laitages ; sommet = graisses ajoutées/sucres/sel (à limiter).
\end{itemize}
\end{aretenir}

Maintenant que tu sais \emph{pourquoi} varier les aliments (les rôles) et \emph{comment} composer une assiette équilibrée (la pyramide), il reste à savoir \emph{combien} en consommer précisément : c'est l'objet de la section suivante, où nous apprendrons à calculer la \cle{valeur énergétique} d'un repas et à élaborer des \cle{rations alimentaires} chiffrées adaptées à chaque situation.

\coursec{Valeur énergétique des aliments et calcul d'un repas}

Dans la section précédente, nous avons vu qu'une alimentation équilibrée doit couvrir les besoins de construction (protéines), de protection (vitamines, sels minéraux) et \emph{d'énergie} (glucides, lipides). Mais comment savoir si un repas apporte \emph{assez} d'énergie ? Comment comparer un plat de ndolè, une assiette de riz ou une cuillère d'huile de palme ? Il faut pour cela apprendre à \cle{mesurer l'énergie} contenue dans les aliments. C'est l'objet de cette section, très souvent évaluée au Probatoire.

\courssub{La notion de valeur énergétique}

\textbf{Une observation de la vie courante.} Un élève qui déjeune le matin d'un beignet et d'un verre de jus tient difficilement jusqu'à midi ; un autre qui mange du foufou avec de la sauce garde sa concentration toute la matinée. De même, les cultivateurs de la région de l'Ouest savent qu'une cuillère d'huile de palme « nourrit » davantage qu'une cuillère de bouillon. Tous les aliments ne fournissent pas la même quantité d'énergie : il faut une grandeur pour le quantifier.

\textbf{Hypothèse.} Les aliments contiennent des nutriments (glucides, protéines, lipides) dont l'oxydation dans nos cellules (respiration cellulaire) libère de l'énergie utilisable sous forme d'ATP, avec dégagement de chaleur, exactement comme la combustion du bois libère de la chaleur. On peut donc mesurer cette énergie.

\textbf{Données expérimentales.} Historiquement, on a utilisé la \cle{bombe calorimétrique} (appareil dans lequel on brûle complètement une masse connue d'aliment sec dans l'oxygène pur et où l'on mesure la chaleur dégagée) pour déterminer l'énergie brute des nutriments. Cependant, dans l'organisme, l'oxydation des protéines est incomplète (production d'urée) et la digestibilité n'est pas totale. Les valeurs de référence utilisées en nutrition humaine sont les \cle{facteurs d'Atwater}, issus de mesures calorimétriques corrigées de la digestibilité et des pertes azotées urinaires. Ces valeurs, constantes et admises internationalement, sont les suivantes :

\begin{center}
\begin{tabularx}{\linewidth}{|l|X|X|}
\hline
\rowcolor{popblueL} \textbf{Nutriment} & \textbf{Énergie libérée par gramme} & \textbf{Équivalent en kcal} \\
\hline
Glucides & 17 kJ/g & $\approx$ 4 kcal/g \\
\hline
Protéines & 17 kJ/g & $\approx$ 4 kcal/g \\
\hline
Lipides & 38 kJ/g & $\approx$ 9 kcal/g \\
\hline
\end{tabularx}
\end{center}

\begin{definition}[Valeur énergétique d'un aliment]
La \cle{valeur énergétique} d'un aliment est la quantité d'énergie métabolisable libérée par l'oxydation complète de ses nutriments dans l'organisme. Elle s'exprime en \cle{kilojoules} (kJ) ou en \cle{kilocalories} (kcal), le plus souvent pour 100 g d'aliment tel que consommé. L'unité internationale d'énergie est le \cle{joule} (J) ; on utilise ses multiples : $1~\text{kJ} = 1000~\text{J}$. La relation entre les deux unités usuelles est :
\[ 1~\text{kcal} = 4{,}18~\text{kJ} \qquad \text{donc} \qquad 1~\text{kJ} \approx 0{,}24~\text{kcal}. \]
\end{definition}

\begin{propriete}[Valeurs énergétiques de référence -- Facteurs d'Atwater (admis)]
L'oxydation métabolisable d'un gramme de nutriment libère :
\begin{itemize}
\item 1 g de \textbf{glucides} $\rightarrow$ 17 kJ ($\approx$ 4 kcal) ;
\item 1 g de \textbf{protéines} $\rightarrow$ 17 kJ ($\approx$ 4 kcal) ;
\item 1 g de \textbf{lipides} $\rightarrow$ 38 kJ ($\approx$ 9 kcal).
\end{itemize}
Les \cle{lipides} sont donc les nutriments les plus énergétiques : à masse égale, ils fournissent \emph{plus du double} de l'énergie des glucides ou des protéines. Ces valeurs sont admises et doivent être connues par cœur pour le Probatoire.
\end{propriete}

\begin{popfigure}
\begin{center}
\begin{tikzpicture}
\begin{axis}[
    ybar, bar width=1.2cm,
    width=0.85\linewidth, height=6.5cm,
    ymin=0, ymax=45,
    ylabel={Énergie libérée (kJ/g)},
    symbolic x coords={Glucides, Protéines, Lipides},
    xtick=data,
    nodes near coords,
    every node near coord/.append style={font=\small\bfseries, color=popdark},
    ytick={0,10,17,20,30,38,40},
    axis lines=left,
    ymajorgrids=true, grid style={popline!40},
]
\addplot[fill=popblue, draw=popdark] coordinates {(Glucides,17) (Protéines,17) (Lipides,38)};
\end{axis}
\end{tikzpicture}
\end{center}
\legende{Énergie métabolisable libérée par l'oxydation d'un gramme de chaque nutriment (facteurs d'Atwater). Les lipides (38 kJ/g) sont plus de deux fois plus énergétiques que les glucides et les protéines (17 kJ/g).}
\end{popfigure}

\begin{exemplebox}[Pourquoi l'huile de palme est-elle si « riche » ?]
100 g d'huile de palme contiennent environ 100 g de lipides. L'énergie fournie est donc :
\[ E = 38 \times 100 = 3800~\text{kJ} \approx 909~\text{kcal}. \]
Pour obtenir la même énergie avec du riz cuit (environ 28 g de glucides pour 100 g), il faudrait :
\[ m = \frac{3800}{17} \approx 224~\text{g de glucides}, \quad \text{soit environ } \frac{224}{0{,}28} \approx 800~\text{g de riz cuit}. \]
Une simple cuillère d'huile (10 g $\approx$ 380 kJ) « pèse » donc énergétiquement autant qu'une grande portion de riz (environ 80 g). C'est pourquoi les régimes trop riches en huile favorisent l'obésité, problème de santé croissant dans les villes camerounaises.
\end{exemplebox}

\courssub{La formule de calcul de la valeur énergétique}

Un aliment ou un repas contient en général les trois types de nutriments. L'énergie totale est la \emph{somme} des énergies fournies par chacun.

\begin{aretenir}[Formule fondamentale]
Pour un aliment ou un repas contenant $m_G$ grammes de glucides, $m_P$ grammes de protéines et $m_L$ grammes de lipides, la valeur énergétique totale est :
\[ E = 17 \times m_G \;+\; 17 \times m_P \;+\; 38 \times m_L \qquad (E \text{ en kJ, masses en g}). \]
Pour convertir en kilocalories : $E(\text{kcal}) = \dfrac{E(\text{kJ})}{4{,}18}$.
\end{aretenir}

\begin{methode}[Calculer la valeur énergétique d'un repas]
\begin{enumerate}
\item \textbf{Identifier} les aliments du repas et relever, dans une table de composition des aliments, les masses de glucides, protéines et lipides qu'ils contiennent (attention : les tables donnent les valeurs pour 100 g d'aliment tel que consommé ; il faut faire une règle de trois si la portion diffère).
\item \textbf{Additionner} les masses de chaque nutriment apportées par tous les aliments : on obtient $m_G$, $m_P$, $m_L$ totales.
\item \textbf{Appliquer la formule} : multiplier chaque masse par son coefficient (17, 17, 38) puis additionner.
\item \textbf{Convertir} éventuellement le résultat en kcal en divisant par 4,18.
\item \textbf{Comparer} le résultat aux besoins de l'individu (voir section suivante sur les rations alimentaires).
\end{enumerate}
\end{methode}

\begin{exoresolu}[Valeur énergétique d'un repas]
Énoncé. Le déjeuner d'un élève de Première apporte au total 80 g de glucides, 20 g de protéines et 15 g de lipides. Calculer la valeur énergétique de ce repas en kJ, puis en kcal.
\tcblower
Solution détaillée.
\begin{enumerate}
\item On identifie les masses : $m_G = 80$ g ; $m_P = 20$ g ; $m_L = 15$ g.
\item On applique la formule :
\begin{align*}
E &= 17 \times m_G + 17 \times m_P + 38 \times m_L \\
E &= 17 \times 80 + 17 \times 20 + 38 \times 15 \\
E &= 1360 + 340 + 570 \\
E &= 2270~\text{kJ}.
\end{align*}
\item Conversion en kilocalories :
\[ E = \frac{2270}{4{,}18} \approx 543~\text{kcal}. \]
\end{enumerate}
Conclusion : ce repas fournit 2270 kJ, soit environ 543 kcal. Pour un adolescent dont les besoins quotidiens sont de l'ordre de 10 000 à 12 000 kJ, ce déjeuner couvre environ un cinquième des besoins de la journée : c'est cohérent pour un repas.
\end{exoresolu}

\courssub{Lecture et utilisation d'une table de composition des aliments}

Les diététiciens et les services de nutrition (par exemple ceux des hôpitaux de Yaoundé ou de Douala) utilisent des \cle{tables de composition des aliments} qui indiquent, pour 100 g d'aliment tel que consommé, les masses de nutriments et la valeur énergétique. Voici un extrait simplifié portant sur des aliments courants au Cameroun :

\begin{center}
\begin{tabularx}{\linewidth}{|l|X|X|X|X|}
\hline
\rowcolor{popblueL} \textbf{Aliment (pour 100 g)} & \textbf{Glucides (g)} & \textbf{Protéines (g)} & \textbf{Lipides (g)} & \textbf{Énergie (kJ)} \\
\hline
Riz blanc cuit & 28 & 2{,}7 & 0{,}3 & $\approx$ 538 \\
\hline
Arachides grillées & 16 & 26 & 49 & $\approx$ 2576 \\
\hline
Huile de palme & 0 & 0 & 100 & $\approx$ 3800 \\
\hline
Poisson fumé (capitaine) & 0 & 22 & 5 & $\approx$ 564 \\
\hline
\end{tabularx}
\end{center}

\begin{exemplebox}[Vérification avec la table]
Vérifions la valeur énergétique des arachides grillées :
\[ E = 17 \times 16 + 17 \times 26 + 38 \times 49 = 272 + 442 + 1862 = 2576~\text{kJ} \approx 616~\text{kcal}. \]
On retrouve bien la valeur de la table. Remarquons que le poisson fumé, très riche en protéines mais pauvre en lipides, est peu énergétique : c'est un excellent aliment de construction sans excès d'énergie.
\end{exemplebox}

\begin{attention}[Les pièges classiques du calcul énergétique]
\begin{itemize}
\item \textbf{Ne jamais additionner des kJ et des kcal.} Si un énoncé mélange les deux unités, convertis d'abord tout dans la même unité ($1~\text{kcal} = 4{,}18~\text{kJ}$). C'est l'erreur la plus fréquente au Probatoire.
\item \textbf{Les coefficients s'appliquent aux masses de nutriments, pas à la masse d'aliment.} 100 g de riz cuit ne contiennent que 28 g de glucides : le reste est de l'eau, des vitamines, des sels minéraux, des fibres, qui ne fournissent pas d'énergie significative.
\item \textbf{Respecte la règle de trois} quand la portion n'est pas de 100 g : pour 250 g de riz cuit, les glucides valent $28 \times 2{,}5 = 70$ g.
\item \textbf{Présente le calcul en étapes} (énergie des glucides, des protéines, des lipides, puis somme) : une réponse sans détail perd des points même si le résultat est juste.
\end{itemize}
\end{attention}

\begin{savaistu}[Étiquettes, santé publique et facteurs d'Atwater]
Les étiquettes des produits alimentaires vendus dans les supermarchés de Douala ou Yaoundé indiquent la valeur énergétique en kJ \emph{et} en kcal pour 100 g (ou par portion). Savoir lire ces étiquettes est une compétence de citoyen : elle permet de limiter les excès énergétiques responsables de l'obésité et du diabète de type 2, maladies en forte progression au Cameroun. À l'inverse, dans les zones où sévit la malnutrition infantile, les centres de santé utilisent des pâtes thérapeutiques prêtes à l'emploi (comme le « Plumpy'nut »\textregistered), à base d'arachide, de lait en poudre, d'huile et de sucre, enrichies en vitamines et minéraux. Ces pâtes concentrent plus de 2000 kJ dans 100 g grâce à leur forte teneur en lipides (facteur 38 kJ/g) : c'est l'application directe de la propriété selon laquelle les lipides sont les nutriments les plus énergétiques. Notons que les facteurs d'Atwater (17, 17, 38) sont des moyennes ; la valeur réelle peut varier légèrement selon la nature précise des acides gras ou des glucides (fibres vs amidon), mais ces valeurs de référence sont celles exigées aux examens.
\end{savaistu}

\begin{exercice}[À toi de jouer]
Un repas est composé de 200 g de riz blanc cuit, 50 g d'arachides grillées et 10 g d'huile de palme. En utilisant la table de composition ci-dessus, calculer la valeur énergétique totale de ce repas en kJ puis en kcal.
\end{exercice}

\begin{corrige}[Exercice : valeur énergétique d'un repas composé]
Étape 1 : masses de nutriments par aliment (règle de trois à partir des valeurs pour 100 g).
\begin{itemize}
\item Riz (200 g) : glucides $28 \times 2 = 56$ g ; protéines $2{,}7 \times 2 = 5{,}4$ g ; lipides $0{,}3 \times 2 = 0{,}6$ g.
\item Arachides (50 g) : glucides $16 \times 0{,}5 = 8$ g ; protéines $26 \times 0{,}5 = 13$ g ; lipides $49 \times 0{,}5 = 24{,}5$ g.
\item Huile de palme (10 g) : lipides $100 \times 0{,}1 = 10$ g.
\end{itemize}
Étape 2 : totaux : $m_G = 56 + 8 = 64$ g ; $m_P = 5{,}4 + 13 = 18{,}4$ g ; $m_L = 0{,}6 + 24{,}5 + 10 = 35{,}1$ g.

Étape 3 : application de la formule :
\begin{align*}
E &= 17 \times 64 + 17 \times 18{,}4 + 38 \times 35{,}1 \\
E &= 1088 + 312{,}8 + 1333{,}8 \\
E &\approx 2735~\text{kJ} \approx \frac{2735}{4{,}18} \approx 654~\text{kcal}.
\end{align*}
Conclusion : ce repas fournit environ 2735 kJ (654 kcal). On remarque que plus de la moitié de l'énergie provient des lipides (1334 kJ sur 2735 kJ), alors qu'ils ne représentent que 35 g sur 117,5 g de nutriments totaux : cela illustre à nouveau leur forte densité énergétique.
\end{corrige}

\begin{aretenir}[L'essentiel de la section]
\begin{itemize}
\item La valeur énergétique d'un aliment est l'énergie métabolisable libérée par l'oxydation de ses nutriments ; elle s'exprime en kJ ou en kcal ($1~\text{kcal} = 4{,}18~\text{kJ}$).
\item Coefficients à connaître (facteurs d'Atwater) : 17 kJ/g pour les glucides, 17 kJ/g pour les protéines, 38 kJ/g pour les lipides.
\item Formule : $E = 17\,m_G + 17\,m_P + 38\,m_L$ (masses en g, $E$ en kJ).
\item Les tables de composition donnent les nutriments pour 100 g : pense à la règle de trois pour toute autre portion.
\end{itemize}
\end{aretenir}

\coursec{Rations alimentaires et élaboration de menus équilibrés}

Après avoir appris à calculer la valeur énergétique d'un repas, nous devons maintenant répondre à la question centrale : \emph{quelle quantité totale d'aliments un individu doit-il consommer par jour pour être en bonne santé ?} C'est l'objet de la notion de \cle{ration alimentaire}. Au Cameroun, comme ailleurs, la malnutrition (sous-nutrition ou surcharge pondérale) touche encore de nombreux foyers, souvent par méconnaissance des besoins réels adaptés à chaque profil. Construisons ensemble les outils pour élaborer des menus équilibrés, variés et ancrés dans notre terroir.

\courssub{Définition et détermination de la ration alimentaire}

\begin{definition}[Ration alimentaire]
La \textbf{ration alimentaire} est la quantité totale d'aliments (et de boissons) consommée par un individu au cours d'une période de 24 heures. Elle doit couvrir l'ensemble des besoins : énergétiques (glucides, lipides, protéines), plastiques (protéines pour la construction et la réparation), fonctionnels (vitamines, minéraux, fibres, eau) et de plaisir.
\end{definition}

Les besoins énergétiques ne sont pas identiques pour tous. Ils dépendent de paramètres précis que tout bon planificateur nutritionnel doit identifier.

\begin{propriete}[Facteurs de variation des besoins énergétiques (admis)]
Les besoins énergétiques quotidiens (BEJ) varient principalement selon :
\begin{itemize}
    \item \textbf{L'âge :} forts besoins chez l'enfant et l'adolescent (croissance) ; baisse progressive chez le sujet âgé.
    \item \textbf{Le sexe :} à masse corporelle et activité égales, l'homme a des besoins supérieurs de 10 à 15 \% à ceux de la femme (masse musculaire plus importante).
    \item \textbf{La masse corporelle et la taille :} plus le corps est grand et lourd, plus le métabolisme de base est élevé.
    \item \textbf{L'activité physique :} c'est le facteur le plus variable (sédentaire, modéré, intense, sportif de haut niveau).
    \item \textbf{L'état physiologique :} grossesse (besoins majorés au 2\textsuperscript{e} et 3\textsuperscript{e} trimestres), allaitement, convalescence, croissance pubertaire.
\end{itemize}
\end{propriete}

\begin{popfigure}
\begin{tikzpicture}
\begin{axis}[
    xbar, width=\linewidth, height=8cm,
    xlabel={Besoins énergétiques quotidiens (kJ/jour)},
    ytick={1,2,3,4,5},
    yticklabels={Enfant (6--10 ans), Adolescent (15--18 ans), Adulte actif, Femme enceinte/allaitante, Sportif de haut niveau},
    xmin=0, xmax=17000,
    bar width=12pt,
    nodes near coords, nodes near coords align={horizontal},
    point meta=explicit symbolic,
    enlarge y limits=0.3,
    axis lines*=left,
    tick label style={font=\small},
    label style={font=\small},
    x tick label style={font=\small, /pgf/number format/use comma, /pgf/number format/1000 sep={\,}},
]
\addplot[popblue, fill=popblue!60] coordinates {
    (7000,1) [7000 kJ]
    (11500,2) [11500 kJ]
    (11000,3) [11000 kJ]
    (12500,4) [12500 kJ]
    (16000,5) [16000 kJ]
};
\end{axis}
\end{tikzpicture}
\legende{Comparaison des besoins énergétiques quotidiens indicatifs selon le profil (valeurs moyennes arrondies). Source : ANC/FAO/OMS adaptées.}
\end{popfigure}

\begin{aretenir}[Ordres de grandeur des besoins énergétiques quotidiens (kJ/jour)]
\begin{center}
\begin{tabularx}{\linewidth}{|l|X|}\hline
\textbf{Profil} & \textbf{Fourchette indicative (kJ/jour)} \\ \hline
Enfant (6--10 ans) & 6000 -- 8000 \\ \hline
Adolescent (15--18 ans) & 10000 -- 13000 \\ \hline
Adulte actif (homme) & 10000 -- 12000 \\ \hline
Adulte active (femme) & 9000 -- 11000 \\ \hline
Femme enceinte (3\textsuperscript{e} trimestre) / Allaitante & + 1000 à 2000 kJ par rapport à la base \\ \hline
Sportif (entraînement intense) & 13000 -- 16000 et plus \\ \hline
\end{tabularx}
\end{center}
\end{aretenir}

\courssub{Répartition qualitative de la ration : les repères nutritionnels}

Une ration équilibrée ne se résume pas à un total de kilojoules. La \cle{qualité} des apports est primordiale. Les recommandations officielles (PNNS, OMS, FAO) fixent des pourcentages cibles pour la répartition de l'énergie totale entre les trois macronutriments.

\begin{propriete}[Répartition énergétique de référence pour l'adulte et l'adolescent]
Pour une alimentation équilibrée, l'apport énergétique total (AET) doit provenir approximativement de :
\begin{itemize}
    \item \textbf{Glucides :} 50 à 55 \% (dont sucres ajoutés < 10 \%) ;
    \item \textbf{Lipides :} 30 à 35 \% (dont acides gras saturés < 12 \%) ;
    \item \textbf{Protéines :} 12 à 15 \% (animales + végétales, ratio 1/1 souhaitable).
\end{itemize}
Les fibres (25--30 g/j), l'eau (1,5--2 L/j) et la densité en micronutriments complètent ce tableau.
\end{propriete}

\begin{popfigure}
\begin{tikzpicture}
    % Camembert manuel avec TikZ standard (sans pgf-pie)
    \colorlet{gluc}{popgreen!80}
    \colorlet{lip}{poporange!80}
    \colorlet{prot}{popblue!80}
    % Glucides 55% = 198°
    \draw[fill=gluc, thick] (0,0) -- (3,0) arc (0:198:3) -- cycle;
    % Lipides 30% = 108° (198 à 306)
    \draw[fill=lip, thick] (0,0) -- (198:3) arc (198:306:3) -- cycle;
    % Protéines 15% = 54° (306 à 360)
    \draw[fill=prot, thick] (0,0) -- (306:3) arc (306:360:3) -- cycle;
    % Légende
    \begin{scope}[xshift=4cm, yshift=1cm, font=\small]
        \fill[gluc] (0,0) rectangle (0.5,0.3) node[right, xshift=0.2cm] {Glucides (55\%)};
        \fill[lip] (0,-0.6) rectangle (0.5,-0.3) node[right, xshift=0.2cm] {Lipides (30\%)};
        \fill[prot] (0,-1.2) rectangle (0.5,-0.9) node[right, xshift=0.2cm] {Protéines (15\%)};
    \end{scope}
\end{tikzpicture}
\legende{Répartition idéale de l'apport énergétique total (AET) entre les macronutriments pour un adolescent ou un adulte.}
\end{popfigure}

\courssub{Répartition chronologique : le rythme des repas}

L'organisme humain fonctionne selon des rythmes (nycthéméraux). Fractionner la ration sur la journée optimise la digestion, l'assimilation et évite les pics d'hypoglycémie ou de surcharge digestive.

\begin{propriete}[Répartition conseillée de la ration journalière]
\begin{center}
\begin{tabularx}{\linewidth}{|l|X|X|}\hline
\textbf{Prise alimentaire} & \textbf{Part de l'AET} & \textbf{Composition type} \\ \hline
Petit-déjeuner & 20 -- 25 \% & Céréale + produit laitier + fruit + boisson (réhydratation) \\ \hline
Déjeuner & 35 -- 40 \% & Entrée crue + Plat protidique + Féculent/Légume + Fruit/Laitage \\ \hline
Collation (goûter) & 5 -- 10 \% (optionnel) & Fruit + produit céréalier ou laitier (évite le grignotage) \\ \hline
Dîner & 25 -- 30 \% & Plus léger que le déjeuner : Potage/Légumes + Protéine maigre + Féculent modéré + Fruit \\ \hline
\end{tabularx}
\end{center}
\end{propriete}

\begin{attention}[Pièges fréquents]
\begin{itemize}
    \item \textbf{Sauter le petit-déjeuner :} Entraîne une hypoglycémie matinale, baisse de vigilance, grignotage sucré/gras en matinée et surcharge au déjeuner.
    \item \textbf{Dîner trop copieux :} Perturbe le sommeil, favorise le stockage lipidique (insuline élevée le soir) et la prise de poids.
    \item \textbf{Confondre "équilibré" et "calorique" :} Un menu de 12000 kJ fait uniquement de riz-huile-sucre apporte l'énergie mais zéro protéine, zéro vitamine, zéro fibre $\rightarrow$ \textbf{malnutrition qualitative}.
\end{itemize}
\end{attention}

\courssub{Méthode d'élaboration d'un menu équilibré}

\begin{methode}[Construire un menu équilibré adapté à un profil -- 5 étapes]
\begin{enumerate}
    \item \textbf{Profilier l'individu :} Âge, sexe, poids, taille, activité physique (PAL), état physiologique (grossesse, pathologie).
    \item \textbf{Fixer la ration énergétique cible (BEJ) :} Utiliser les tables de référence (voir tableau ci-dessus) ou formules de calcul (ex: Black et al. pour le métabolisme de base $\times$ PAL).
    \item \textbf{Répartir l'énergie entre les macronutriments :} Appliquer les pourcentages 55/30/15. Convertir en grammes (1 g Glucides = 17 kJ ; 1 g Lipides = 37 kJ ; 1 g Protéines = 17 kJ).
    \item \textbf{Choisir des aliments locaux, variés, de saison :} Respecter la structure des repas (entrée, plat, dessert). Intégrer les 3 groupes d'aliments simples (énergétiques, constructeurs, protecteurs) + eau + fibres à chaque repas principal.
    \item \textbf{Vérifier par le calcul :} Calculer l'énergie réelle du menu proposé. Comparer à la cible ($\pm$ 10 \%). Vérifier la présence de fruits/légumes crus et cuits, source de calcium, source de fer hémique.
\end{enumerate}
\end{methode}

\courssub{Application : Menu équilibré camerounais pour un élève de Première}

Mettons la méthode en œuvre pour un cas concret : \textbf{Kévin, 17 ans, élève en 1ère A, activité modérée (marche 30 min/jour, cours assis), 60 kg, 1,70 m}.

\begin{exoresolu}[Élaboration et vérification d'un menu pour Kévin (besoin $\approx$ 12000 kJ)]
\textbf{Étape 1 -- Profil :} Adolescent, garçon, activité modérée $\rightarrow$ BEJ visé = \textbf{12 000 kJ}.

\textbf{Étape 2 -- Répartition macronutriments cibles :}
\begin{itemize}
    \item Glucides (55 \%) : $0,55 \times 12000 = 6600$ kJ $\rightarrow$ $6600 / 17 \approx \textbf{388 g}$.
    \item Lipides (30 \%) : $0,30 \times 12000 = 3600$ kJ $\rightarrow$ $3600 / 37 \approx \textbf{97 g}$.
    \item Protéines (15 \%) : $0,15 \times 12000 = 1800$ kJ $\rightarrow$ $1800 / 17 \approx \textbf{106 g}$.
\end{itemize}

\textbf{Étape 3 -- Proposition de menu (aliments locaux) :}

\begin{center}
\begin{tabularx}{\linewidth}{|l|X|}\hline
\textbf{Repas} & \textbf{Menu détaillé} \\ \hline
\textbf{Petit-déj. (25 \% $\approx$ 3000 kJ)} & Bouillie de maïs fermenté (300 g) enrichie en lait en poudre (20 g) et arachides grillées broyées (15 g) + 1 banane plantain mûre (100 g) + Eau/Thé léger. \\ \hline
\textbf{Déjeuner (40 \% $\approx$ 4800 kJ)} & Riz blanc cuit (200 g cru) + Ndolé (feuilles de vernonia 150 g + arachide 30 g + crevettes séchées 20 g + huile rouge 15 ml) + 1 orange (150 g) + Eau. \\ \hline
\textbf{Collation (5 \% $\approx$ 600 kJ)} & 1 mangue (150 g) + Pain de maïs (50 g) ou beignet de haricot (accara) petit format. \\ \hline
\textbf{Dîner (30 \% $\approx$ 3600 kJ)} & Ignames pilées (300 g cuites) + Sauce graine (pâte de noix de palme 30 g + poisson fumé 50 g + légumes-feuilles 100 g + huile 10 ml) + Papaye (100 g) + Eau. \\ \hline
\end{tabularx}
\end{center}

\textbf{Étape 4 -- Vérification numérique (calcul simplifié sur bases moyennes) :}

\begin{center}
\begin{tabular}{|l|cccc|c|}\hline
\textbf{Aliment (portion)} & \textbf{Glucides (g)} & \textbf{Lipides (g)} & \textbf{Protéines (g)} & \textbf{Energie (kJ)} \\ \hline
Bouillie maïs/lait/arachide (300+20+15g) & 55 & 12 & 14 & 1750 \\
Banane plantain (100g) & 22 & 0,2 & 1,2 & 400 \\
Riz cru (200g) & 156 & 1 & 14 & 3000 \\
Ndolé (feuilles+arachide+crevettes+huile) & 15 & 25 & 18 & 1350 \\
Orange (150g) & 15 & 0,2 & 1,5 & 280 \\
Mangue (150g) & 20 & 0,3 & 1 & 380 \\
Pain maïs (50g) & 25 & 1,5 & 3 & 580 \\
Ignames cuites (300g) & 65 & 0,5 & 4 & 1200 \\
Sauce graine (noix palme+poisson+légumes+huile) & 10 & 35 & 25 & 1850 \\
Papaye (100g) & 8 & 0,1 & 0,5 & 150 \\ \hline
\textbf{TOTAL JOUR} & \textbf{391 g} & \textbf{75,8 g} & \textbf{82,2 g} & \textbf{10940 kJ} \\ \hline
\end{tabular}
\end{center}

\textbf{Interprétation :}
\begin{itemize}
    \item Énergie totale : \textbf{10 940 kJ} (légèrement sous la cible de 12 000 kJ, $\approx$ -9 \%, acceptable). On peut augmenter légèrement les portions de riz/igname ou ajouter une collation plus consistante.
    \item Glucides : 391 g $\times$ 17 = 6647 kJ (\textbf{60,7 \%}) $\rightarrow$ Dans la cible (50-55 \%), légèrement haut mais acceptable pour un adolescent actif.
    \item Lipides : 75,8 g $\times$ 37 = 2805 kJ (\textbf{25,6 \%}) $\rightarrow$ Légèrement bas (cible 30-35 \%). L'huile rouge et l'arachide apportent des AGS et AGMI/AGPI. On peut ajouter 5-10 ml d'huile.
    \item Protéines : 82,2 g $\times$ 17 = 1397 kJ (\textbf{12,8 \%}) $\rightarrow$ Dans la cible (12-15 \%). Ratio Protéines Animales/Végétales $\approx$ 40/60 (poisson, crevettes, lait vs arachide, maïs, légumineuses). Correct.
    \item \textbf{Qualité :} Présence de fruits/légumes à chaque repas (Vit C, fibres, caroténoïdes), sources de Calcium (lait, feuilles vertes, poisson avec arêtes), Fer hémique (poisson, crevettes) + Fer non hémique + Vit C (absorption favorisée). Menu \textbf{équilibré}.
\end{itemize}
\end{exoresolu}

\begin{exemplebox}[Exemple chiffré : Besoin en glucides d'un adolescent]
Un adolescent ayant un besoin énergétique de 12 000 kJ/jour doit tirer 55 \% de son énergie des glucides.
\[ E_{\text{gluc}} = 0,55 \times 12000 = 6600 \text{ kJ} \]
Puisque 1 g de glucides libère 17 kJ :
\[ m_{\text{gluc}} = \frac{6600}{17} \approx 388 \text{ g de glucides par jour}. \]
Cela représente environ 70 g au petit-déjeuner, 150 g au déjeuner, 40 g au goûter, 130 g au dîner. Une bouillie de maïs (100 g farine) + pain + riz + igname + fruits couvrent aisément ce besoin.
\end{exemplebox}

\courssub{Exercices d'application}

\begin{exercice}[Calcul de ration pour une femme allaitante]
Marie, 25 ans, 60 kg, allaite son enfant de 3 mois (allaitement exclusif). Son activité est modérée (ménage, marche). Son besoin de base (avant grossesse) est de 9500 kJ/j.
\begin{enumerate}
    \item Calculer son besoin énergétique journalier actuel (majoration allaitement = + 2000 kJ/j).
    \item Déterminer les apports cibles en g de glucides, lipides, protéines (répartition 55/30/15).
    \item Proposer un menu type camerounais pour une journée (Petit-déj, Déj, Dîner) en justifiant les choix (disponibilité locale, densité nutritionnelle).
\end{enumerate}
\end{exercice}

\begin{corrige}[Exercice n°1]
\begin{enumerate}
    \item BEJ = 9500 + 2000 = \textbf{11 500 kJ/j}.
    \item Cibles : Glucides : $0,55 \times 11500 / 17 = 372$ g. Lipides : $0,30 \times 11500 / 37 = 93$ g. Protéines : $0,15 \times 11500 / 17 = 101$ g.
    \item \textbf{Menu proposé :} Pdj : Bouillie mil/maïs + lait + arachide + papaye. Déj : Fufu maïs/manioc (300g) + Sauce okro/feuilles (poisson sec + huile + légumes) + Ananas. Dîner : Pâte de maïs (200g) + Sauce pistache (arachide/courge + viande/poisson + feuilles) + Banane. Riche en calcium (lait, feuilles, poisson entier), fer, folates, protéines, énergie dense.
\end{enumerate}
\end{corrige}

\begin{exercice}[Analyse critique de menu]
Un élève propose ce menu pour un camarade de 17 ans (besoin 12000 kJ) :
\textit{Pdj : Pain beurré + Chocolat chaud sucré. Déj : Frites (300g) + Steak haché (100g) + Ketchup + Soda (33cl). Dîner : Pâtes (200g cru) + Sauce tomate industrielle + Fromage râpé + Glace.}
\begin{enumerate}
    \item Ce menu est-il énergétiquement suffisant ? (Estimer grossièrement).
    \item Est-il \textbf{équilibré} ? Justifier en citant au moins 4 écarts majeurs par rapport aux recommandations (répartition macronutriments, fibres, micronutriments, rythme, qualité lipides/glucides).
    \item Proposer une version corrigée en gardant l'esprit "repas rapide/aimé des jeunes" mais équilibrée.
\end{enumerate}
\end{exercice}

\begin{corrige}[Exercice n°2]
\begin{enumerate}
    \item Estimation : Pain (600) + Beurre (300) + Choc chaud (800) $\approx$ 1700. Frites (1500) + Steak (800) + Soda (600) $\approx$ 2900. Pâtes (1400) + Sauce/Fromage (500) + Glace (800) $\approx$ 2700. Total $\approx$ 7300 kJ. \textbf{Insuffisant} (déficit $\approx$ 40 \%).
    \item \textbf{Non équilibré. Écarts :} 1) Lipides $\gg$ 35 \% (frites, beurre, fromage, steak gras, glace) -- excès AGS. 2) Glucides raffinés/sucres ajoutés $\gg$ 10 \% (soda, ketchup, chocolat, pain blanc, glace) -- index glycémique élevé. 3) Fibres quasi absentes (pas de légumes crus, peu de complets, pas de fruits frais). 4) Micronutriments faibles (Vit C, B9, Mg, K). 5) Répartition journalière déséquilibrée (dîner trop gras/lourd). 6) Apport hydrique insuffisant (soda $\neq$ eau).
    \item \textbf{Version corrigée :} Pdj : Pain complet + beurre léger + yaourt + fruit frais + eau. Déj : Sandwich pain complet (poulet grillé + crudités + léger fromage) + fruit + eau. Dîner : Pâtes complètes + sauce tomate maison (légumes, huile olive) + poisson/viande maigre + salade verte + fruit/laitage.
\end{enumerate}
\end{corrige}

\begin{savaistu}[Le sportif de haut niveau : une machine énergétique]
Un footballeur professionnel (Lion Indomptable) ou un marathonien camerounais en période de compétition peut dépenser \textbf{16 000 à 20 000 kJ par jour} (soit 2 fois un adulte sédentaire). Leur ration doit être hyper-énergétique mais \emph{digeste} : charge glucidique massive (pâtes, riz, bouillies, patates douces) avant l'effort, protéines de haute qualité (poisson, œufs, lait) pour la récupération, hydratation continue (eau, boissons d'effort). L'erreur classique : manger "lourd" (beaucoup d'huile, viandes grasses, piments forts) avant le match $\rightarrow$ troubles digestifs, baisse de performance. La nutrition sportive est une science à part entière !
\end{savaistu}

\begin{aretenir}[Bilan : Les 3 piliers d'une ration réussie]
\begin{enumerate}
    \item \textbf{Quantité adaptée :} Énergie = Besoins (âge, sexe, activité, état physiologique). Ni trop (surpoids), ni trop peu (maigreur, carences).
    \item \textbf{Qualité répartie :} 55 \% Glucides (complexes > simples), 30 \% Lipides (variés, oméga-3), 15 \% Protéines (animales + végétales). Fibres, eau, micronutriments couverts par la diversité.
    \item \textbf{Rythme respecté :} 3 à 4 prises/jour. Petit-déjeuner obligatoire. Dîner léger. Collation si besoin réel.
\end{enumerate}
\end{aretenir}

\newpage

\coursec{Activité d'intégration}

\courssub{Objectifs de l'activité}

À l'issue de cette activité d'intégration, tu seras capable de :

\begin{itemize}
\item mobiliser la \cle{classification des aliments} en groupes nutritionnels (énergétiques glucidiques, énergétiques lipidiques, bâtisseurs, protecteurs) pour composer des repas variés ;
\item appliquer la formule de la \cle{valeur énergétique} $E = 17G + 17P + 38L$ (énergie en kJ, masses en g) pour calculer l'énergie fournie par un aliment puis par un repas complet ;
\item élaborer un \cle{menu équilibré} adapté à une situation concrète (âge, masse corporelle, activité physique intense) en respectant la répartition énergétique de référence 55/30/15 ;
\item vérifier quantitativement la pertinence d'un menu par le calcul d'un \cle{écart relatif} et justifier ses choix devant la classe ;
\item développer des savoir-être : travail en équipe, esprit critique, argumentation scientifique rigoureuse.
\end{itemize}

\courssub{Situation}

\textbf{Diététicien d'un jour : un menu pour le champion de lutte du quartier.}

Dans le quartier Nkolndongo à Yaoundé, le jeune \cle{Ekotto}, 17 ans, est le champion de lutte traditionnelle de son arrondissement. La grande compétition annuelle approche et son entraîneur s'inquiète : Ekotto s'affaiblit, il s'essouffle rapidement et a perdu deux combats amicaux. Sa mère, dépassée, déclare : «~Mon fils mange pourtant beaucoup de foutou chaque jour !~» L'infirmier du centre de santé, consulté, explique que le problème n'est pas la \emph{quantité} mais la \emph{qualité} et la \emph{répartition} de l'alimentation : un sportif de haut niveau a des besoins énergétiques considérables et doit équilibrer glucides, lipides et protides.

Le médecin du sport a établi la fiche suivante :

\begin{center}
\begin{tabularx}{\linewidth}{|l|X|}
\hline
\rowcolor{popblueL} \textbf{Caractéristique} & \textbf{Donnée} \\
\hline
Âge / sexe & 17 ans, masculin (adolescent en croissance) \\
\hline
Masse corporelle & 70 kg \\
\hline
Activité & entraînement intensif de lutte, 3 h par jour \\
\hline
Besoin énergétique estimé & \SI{14000}{kJ} par jour \\
\hline
Consigne du médecin & répartition : 55 \% de l'énergie sous forme de glucides, 30 \% sous forme de lipides, 15 \% sous forme de protides \\
\hline
\end{tabularx}
\end{center}

Ton groupe de SVTEEHB est sollicité comme équipe de «~diététiciens d'un jour~». Vous disposez de la table de composition simplifiée suivante (masses en grammes pour 100 g d'aliment) :

\begin{center}
\begin{tabularx}{\linewidth}{|X|c|c|c|}
\hline
\rowcolor{popgreenL} \textbf{Aliment (100 g)} & \textbf{Glucides (g)} & \textbf{Protides (g)} & \textbf{Lipides (g)} \\
\hline
Riz cuit & 28 & 2,7 & 0,3 \\
\hline
Foutou de manioc & 40 & 1 & 0,2 \\
\hline
Ndolé (feuilles + arachide + crevettes) & 5 & 8 & 12 \\
\hline
Arachides grillées & 10 & 26 & 49 \\
\hline
Huile de palme & 0 & 0 & 100 \\
\hline
Poisson fumé & 0 & 25 & 5 \\
\hline
Haricots cuits & 20 & 9 & 0,5 \\
\hline
Banane plantain cuite & 32 & 1 & 0,3 \\
\hline
Avocat & 2 & 2 & 15 \\
\hline
Lait & 5 & 3,3 & 3,5 \\
\hline
Sucre & 100 & 0 & 0 \\
\hline
Feuilles de manioc (pondu) & 4 & 7 & 1 \\
\hline
\end{tabularx}
\end{center}

\courssub{Consignes}

Par groupes de 4 élèves, vous traiterez les tâches suivantes, dans l'ordre :

\begin{enumerate}
\item \textbf{Classer.} Répartir les 12 aliments du tableau dans les quatre groupes nutritionnels : \emph{aliments énergétiques glucidiques}, \emph{aliments énergétiques lipidiques}, \emph{aliments bâtisseurs (riches en protides)} et \emph{aliments protecteurs (riches en vitamines et sels minéraux)}. Justifier chaque classement à partir des données du tableau (teneur majoritaire en macronutriments) et de la nature de l'aliment (simple ou composé).
\item \textbf{Composer.} Proposer un menu complet de la journée (petit-déjeuner, déjeuner, collation, dîner) en choisissant des aliments du tableau et en fixant la masse de chaque portion. Le menu doit être varié, réaliste (aliments disponibles au marché local) et respecter la répartition énergétique 55/30/15.
\item \textbf{Calculer.} Pour chaque aliment, calculer l'énergie fournie par la portion à l'aide de la formule $E = 17G + 17P + 38L$, puis calculer la valeur énergétique totale du menu.
\item \textbf{Vérifier.} Calculer l'écart relatif entre la valeur énergétique du menu et le besoin de \SI{14000}{kJ}. Le menu est accepté si l'écart est inférieur à 10 \%. Vérifier également la répartition effective entre glucides, lipides et protides.
\item \textbf{Présenter et justifier.} Exposer le menu devant la classe en 5 minutes : choix des aliments, équilibre entre les groupes, présence d'aliments protecteurs (vitamines, sels minéraux, fibres), adaptation à l'effort du lutteur. Le barème collectif attribue des points à la justesse des calculs, à la variété du menu et à la qualité de la justification.
\end{enumerate}

\courssub{Ressources à mobiliser}

\begin{definition}[Constituants de la matière vivante et classification des aliments]
La matière vivante est constituée d'eau (\SI{\approx 60}{\%} de la masse corporelle), de matières minérales (sels minéraux) et de matières organiques. Les matières organiques se divisent en : \textbf{glucides} (sucres, amidon, fibres), \textbf{lipides} (graisses, huiles), \textbf{protides} (protéines), \textbf{acides nucléiques} et \textbf{vitamines}.
\begin{itemize}
\item \textbf{Aliment simple :} ne fournit qu'un seul type de nutriment énergétique ou structurel majeur (ex: sucre = glucides ; huile = lipides ; blanc d'œuf = protides).
\item \textbf{Aliment composé :} associe plusieurs nutriments (ex: ndolé, foutou, lait, haricots). La plupart de nos aliments quotidiens sont des aliments composés.
\item \textbf{Groupes nutritionnels (classification fonctionnelle) :}
\begin{itemize}
\item \textit{Énergétiques glucidiques :} apportent surtout du glucose (riz, foutou, banane plantain, sucre, haricots, pain, maïs, igname, manioc).
\item \textit{Énergétiques lipidiques :} apportent surtout des acides gras (huiles, beurre, arachides, avocat, lait entier, fromage, jaune d'œuf).
\item \textit{Bâtisseurs (protides) :} apportent des acides aminés essentiels (viande, poisson, œufs, lait, légumineuses : haricots, soja, arachides ; ndolé grâce aux arachides/crevettes).
\item \textit{Protecteurs / Régulateurs :} apportent vitamines, sels minéraux, fibres, eau (légumes-feuilles : pondu, ndolé, épinards ; fruits : banane, avocat, mangue, papaye ; lait).
\end{itemize}
\end{itemize}
\end{definition}

\begin{methode}[Tests caractéristiques des constituants organiques (TP de mise en évidence)]
Pour identifier un aliment simple ou caractériser un aliment composé, on réalise des tests colorés sur un extrait aqueux (broyat + eau, filtré) :
\begin{enumerate}
\item \textbf{Glucides réducteurs (glucose, fructose, maltose, lactose) :} Test de \textbf{Fehling} (ou Benedict). \textit{Protocole :} 2 mL extrait + 2 mL Fehling A+B $\rightarrow$ bain-marie 5 min. \textit{Résultat + :} précipité rouge brique (oxyde de cuivre(I) $\ce{Cu2O}$).
\item \textbf{Amidon :} Test à l'\textbf{iode} (solution de Lugol). \textit{Protocole :} 1 goutte Lugol sur extrait ou aliment solide. \textit{Résultat + :} coloration bleu-noir/violet intense (complexe amidon-iode).
\item \textbf{Protides :} Test de \textbf{Biuret}. \textit{Protocole :} 2 mL extrait + 2 mL soude $\ce{NaOH}$ 10\% + quelques gouttes sulfate de cuivre $\ce{CuSO4}$ 1\%. \textit{Résultat + :} coloration violet/rose (complexe cuivre-peptides).
\item \textbf{Lipides :} Test au \textbf{Rouge Soudan III} (ou IV). \textit{Protocole :} 1 mL extrait + 1 goutte colorant rouge soudan III (dans l'alcool) + eau. \textit{Résultat + :} coloration rouge des gouttelettes lipidiques (flottent en surface). \textit{Alternative :} tache grasse persistante sur papier buvard.
\item \textbf{Eau :} Test au \textbf{sulfate de cuivre anhydre} (blanc $\rightarrow$ bleu) ou chlorure de cobalt (bleu $\rightarrow$ rose).
\end{enumerate}
\emph{Attention :} Ces tests sont qualitatifs (présence/absence). Pour le dosage quantitatif, on utilise des méthodes instrumentales (spectrophotométrie, chromatographie).
\end{methode}

\begin{aretenir}[Boîte à outils de l'activité]
\begin{itemize}
\item \textbf{Valeurs énergétiques de référence (facteurs d'Atwater) :} 1 g de glucides $\rightarrow$ \SI{17}{kJ} ; 1 g de protides $\rightarrow$ \SI{17}{kJ} ; 1 g de lipides $\rightarrow$ \SI{38}{kJ}.
\item \textbf{Valeur énergétique d'un aliment (portion de masse $m$ en g) :} \[ E = 17 \times G + 17 \times P + 38 \times L \] où $G$, $P$, $L$ sont les masses (en g) de glucides, protides, lipides \textbf{dans la portion consommée} (calculées par règle de trois à partir du tableau pour 100 g).
\item \textbf{Alimentation équilibrée :} ration adaptée à l'âge, à la masse corporelle, à l'état physiologique (croissance, grossesse, allaitement, vieillissement) et à l'activité physique ; répartition conseillée de l'énergie apportée : \textbf{55 \% glucides, 30 \% lipides, 15 \% protides}. Couverture des besoins en vitamines, sels minéraux, fibres et eau.
\item \textbf{Écart relatif (validation du menu) :} $\text{écart} = \dfrac{|E_{\text{menu}} - E_{\text{besoin}}|}{E_{\text{besoin}}} \times 100$ (en \%). Menu valide si écart $< 10\,\%$.
\end{itemize}
\end{aretenir}

\courssub{Démarche et corrigé commenté}

\begin{exoresolu}[Un menu possible pour Ekotto (corrigé type)]
\textbf{Étape 1 — Classification des aliments du tableau.}

On classe selon le nutriment \emph{majoritaire} ou la \emph{fonction nutritionnelle principale}, en notant qu'un aliment composé peut appartenir à deux groupes.

\begin{center}
\begin{tabularx}{\linewidth}{|l|X|}
\hline
\rowcolor{poporangeL} \textbf{Groupe nutritionnel} & \textbf{Aliments du tableau (justification)} \\
\hline
Énergétiques glucidiques & \textbf{Riz cuit} (G=28 g), \textbf{Foutou} (G=40 g), \textbf{Banane plantain} (G=32 g), \textbf{Sucre} (G=100 g), \textbf{Haricots} (G=20 g, aussi bâtisseur). \\
\hline
Énergétiques lipidiques & \textbf{Huile de palme} (L=100 g), \textbf{Arachides grillées} (L=49 g, aussi bâtisseur), \textbf{Avocat} (L=15 g, aussi protecteur), \textbf{Ndolé} (L=12 g, aussi bâtisseur/protecteur). \\
\hline
Bâtisseurs (Protides) & \textbf{Poisson fumé} (P=25 g), \textbf{Arachides grillées} (P=26 g), \textbf{Haricots} (P=9 g), \textbf{Ndolé} (P=8 g), \textbf{Feuilles de manioc} (P=7 g), \textbf{Lait} (P=3,3 g). \\
\hline
Protecteurs (Vitamines, Sels minéraux, Fibres) & \textbf{Feuilles de manioc (pondu)}, \textbf{Ndolé (feuilles)}, \textbf{Avocat}, \textbf{Banane plantain}, \textbf{Lait}. \\
\hline
\end{tabularx}
\end{center}

\textbf{Étape 2 — Menu proposé (portions estimées pour un adolescent sportif de 70 kg).}

\begin{itemize}
\item \textbf{Petit-déjeuner (7 h) :} Banane plantain bouillie 150 g, Arachides grillées 40 g, Lait 250 g sucré avec 30 g de sucre. \textit{(Apport glucidique rapide et soutenu, protides/lipides arachides, calcium lait).}
\item \textbf{Déjeuner (13 h) :} Riz cuit 350 g, Ndolé 200 g, Huile de palme 20 g (dans la sauce), Avocat 150 g. \textit{(Base glucidique riz, protides/lipides ndolé, lipides/vitamines avocat).}
\item \textbf{Collation post-entraînement (17 h) :} Banane plantain 150 g, Arachides grillées 20 g. \textit{(Recharge glycogénique rapide + protides récupération).}
\item \textbf{Dîner (20 h) :} Foutou de manioc 300 g, Feuilles de manioc (pondu) 150 g, Poisson fumé 100 g, Huile de palme 10 g, Haricots cuits 200 g. \textit{(Glucides lents foutou/haricots, protides nobles poisson/haricots, protecteurs pondu, lipides modérés).}
\end{itemize}

\textbf{Étape 3 — Calculs énergétiques détaillés.}

On calcule l'énergie pour 100 g ($E_{100}$), puis pour la portion ($E_{\text{portion}} = E_{100} \times \frac{m}{100}$).

\begin{center}
\small
\begin{tabularx}{\linewidth}{|X|c|c|c|}
\hline
\rowcolor{poporangeL} \textbf{Aliment (Portion)} & \textbf{$E_{100}$ (kJ)} & \textbf{Masse (g)} & \textbf{$E_{\text{portion}}$ (kJ)} \\
\hline
Banane plantain (150 g) & $17\times32 + 17\times1 + 38\times0,3 = 572,4$ & 150 & $572,4 \times 1,5 = 859$ \\
Arachides (40 g) & $17\times10 + 17\times26 + 38\times49 = 2474$ & 40 & $2474 \times 0,4 = 990$ \\
Lait (250 g) & $17\times5 + 17\times3,3 + 38\times3,5 = 274,1$ & 250 & $274,1 \times 2,5 = 685$ \\
Sucre (30 g) & $17\times100 = 1700$ & 30 & $1700 \times 0,3 = 510$ \\
\hline
Riz (350 g) & $17\times28 + 17\times2,7 + 38\times0,3 = 533,3$ & 350 & $533,3 \times 3,5 = 1867$ \\
Ndolé (200 g) & $17\times5 + 17\times8 + 38\times12 = 677$ & 200 & $677 \times 2 = 1354$ \\
Huile palme (20 g) & $38\times100 = 3800$ & 20 & $3800 \times 0,2 = 760$ \\
Avocat (150 g) & $17\times2 + 17\times2 + 38\times15 = 638$ & 150 & $638 \times 1,5 = 957$ \\
\hline
Banane plantain (150 g) & 572,4 & 150 & 859 \\
Arachides (20 g) & 2474 & 20 & 495 \\
\hline
Foutou (300 g) & $17\times40 + 17\times1 + 38\times0,2 = 704,6$ & 300 & $704,6 \times 3 = 2114$ \\
Feuilles manioc (150 g) & $17\times4 + 17\times7 + 38\times1 = 225$ & 150 & $225 \times 1,5 = 338$ \\
Poisson fumé (100 g) & $17\times0 + 17\times25 + 38\times5 = 615$ & 100 & 615 \\
Huile palme (10 g) & 3800 & 10 & 380 \\
Haricots (200 g) & $17\times20 + 17\times9 + 38\times0,5 = 512$ & 200 & $512 \times 2 = 1024$ \\
\hline
\textbf{TOTAL JOURNÉE} & & & \textbf{13 807} \\
\hline
\end{tabularx}
\end{center}
\emph{Note :} Les légères différences avec le tableau de l'énoncé (13 805 kJ) viennent de l'arrondi des $E_{100}$ (ici conservés à 1 décimale). Les deux valeurs sont acceptables.

\textbf{Étape 4 — Vérification quantitative.}

\begin{itemize}
\item \textbf{Écart relatif sur l'énergie totale :}
\[ \text{écart} = \frac{|13807 - 14000|}{14000} \times 100 = \frac{193}{14000} \times 100 \approx 1,4\,\% < 10\,\% \]
$\rightarrow$ \textbf{Menu validé} sur le plan énergétique global.
\item \textbf{Vérification de la répartition (calcul approché des macronutriments totaux) :}
\begin{align*}
\text{Total Glucides } G_{\text{tot}} &\approx 1,5(32) + 0,4(10) + 2,5(5) + 0,3(100) + 3,5(28) + 2(5) + 0,2(0) + 1,5(2) + 1,5(32) + 0,2(10) + 3(40) + 1,5(4) + 1(0) + 0,1(0) + 2(20) \\
&\approx 48 + 4 + 12,5 + 30 + 98 + 10 + 0 + 3 + 48 + 2 + 120 + 6 + 0 + 0 + 40 = \textbf{421,5 g} \\
E_G &= 421,5 \times 17 = \textbf{7166 kJ} \quad (51,9\,\%) \\[4pt]
\text{Total Lipides } L_{\text{tot}} &\approx 1,5(0,3) + 0,4(49) + 2,5(3,5) + 0 + 3,5(0,3) + 2(12) + 0,2(100) + 1,5(15) + 1,5(0,3) + 0,2(49) + 3(0,2) + 1,5(1) + 1(5) + 0,1(100) + 2(0,5) \\
&\approx 0,45 + 19,6 + 8,75 + 0 + 1,05 + 24 + 20 + 22,5 + 0,45 + 9,8 + 0,6 + 1,5 + 5 + 10 + 1 = \textbf{124,7 g} \\
E_L &= 124,7 \times 38 = \textbf{4739 kJ} \quad (34,3\,\%) \\[4pt]
\text{Total Protides } P_{\text{tot}} &\approx 1,5(1) + 0,4(26) + 2,5(3,3) + 0 + 3,5(2,7) + 2(8) + 0 + 1,5(2) + 1,5(1) + 0,2(26) + 3(1) + 1,5(7) + 1(25) + 0 + 2(9) \\
&\approx 1,5 + 10,4 + 8,25 + 0 + 9,45 + 16 + 0 + 3 + 1,5 + 5,2 + 3 + 10,5 + 25 + 0 + 18 = \textbf{111,8 g} \\
E_P &= 111,8 \times 17 = \textbf{1901 kJ} \quad (13,8\,\%)
\end{align*}
\textbf{Bilan répartition :} 52 \% / 34 \% / 14 \%. Proche du 55/30/15. L'écart sur les lipides (+4 \%) s'explique par l'huile de palme et l'avocat/arachides. Pour se rapprocher de 30 \%, on pourrait réduire l'huile de palme à 15 g (déjeuner) et 5 g (dîner) et augmenter le riz à 400 g.
\end{itemize}

\textbf{Étape 5 — Justification nutritionnelle et contextuelle.}

\begin{itemize}
\item \textbf{Couverture des besoins :} Énergie (13 807 kJ $\approx$ 14 000 kJ), Protéines (112 g $\approx$ 1,6 g/kg/j, optimal pour sportif en croissance), Glucides (stocks glycogéniques pour 3 h entraînement), Lipides (acides gras essentiels via arachides/avocat/poisson).
\item \textbf{Équilibre et variété :} 10 aliments différents, 4 groupes nutritionnels représentés à chaque repas principal. Présence de \textbf{protides d'origine animale} (poisson, lait) et \textbf{végétale} (légumineuses, arachides) $\rightarrow$ complémentarité acides aminés.
\item \textbf{Aliments protecteurs :} Ppondu et Ndolé (vit A, C, B9, Fe, Ca, fibres), Avocat (vit E, K, K, Mg), Banane (K, vit B6, C), Lait (Ca, vit D, B12). Couverture large des micronutriments.
\item \textbf{Adaptation à l'effort :} Collation glucidique/protidique immédiate post-entraînement (fenêtre métabolique). Dîner riche en glucides complexes (foutou, haricots) pour recharge glycogénique nocturne.
\item \textbf{Acceptabilité culturelle :} Tous aliments de base au Cameroun (marché Nkolndongo), modes de cuisson habituels (bouilli, sauce, grillé). Coût maîtrisé (produits locaux, peu de viande coûteuse).
\end{itemize}
\end{exoresolu}

\begin{attention}[Pièges fréquents au Probatoire et en TP]
\begin{itemize}
\item \textbf{Conversion portions / 100 g :} Le tableau donne $G, P, L$ pour \textbf{100 g}. Pour une portion de $m$ g, les masses réelles sont $G \times \frac{m}{100}$, etc. \textbf{Ne jamais} utiliser directement les valeurs du tableau pour une portion $\neq$ 100 g.
\item \textbf{Unités de la formule :} $E = 17G + 17P + 38L$ donne l'énergie en \textbf{kJ} si $G, P, L$ sont en \textbf{grammes}. Ne pas confondre avec kcal (facteurs 4, 4, 9).
\item \textbf{Écart relatif :} Le dénominateur est \textbf{toujours la valeur de référence (besoin)}, pas la valeur trouvée. Formule : $\frac{|V_{\text{trouvé}} - V_{\text{ref}}|}{V_{\text{ref}}} \times 100$.
\item \textbf{Équilibre $\neq$ Énergie seule :} Un menu à 14 000 kJ fait seulement de foutou + huile + sucre est \textbf{non équilibré} (carence protides, vitamines, fibres, excès lipides/glucides simples). La répartition 55/30/15 et la diversité sont obligatoires.
\item \textbf{Aliments «~purs~» vs composés :} L'huile de palme et le sucre sont des aliments \emph{simples} (100\% lipides / 100\% glucides). Le ndolé, le foutou, le lait, les haricots sont \emph{composés}. Les classer demande d'identifier le(s) nutriment(s) dominant(s).
\item \textbf{Tests biochimiques :} Fehling $\rightarrow$ glucides \emph{réducteurs} (pas l'amidon direct). Biuret $\rightarrow$ liaisons peptidiques (protides). Rouge Soudan $\rightarrow$ lipides (coloration des gouttelettes, pas de l'eau). Iode $\rightarrow$ amidon spécifiquement.
\end{itemize}
\end{attention}

\courssub{Bilan de l'activité}

Cette activité t'a permis de mettre en évidence plusieurs idées essentielles de la leçon «~Amélioration de la santé par la nutrition~» :

\begin{itemize}
\item \textbf{La quantité ne suffit pas} : Ekotto mangeait «~beaucoup~» (énergie peut-être couverte) mais mal. Une alimentation équilibrée exige une \emph{répartition} correcte entre macronutriments (55/30/15), la présence d'aliments protecteurs (vitamines, sels minéraux, fibres, eau) et une adaptation aux rythmes (collation post-effort).
\item \textbf{Le besoin énergétique est individuel} : Il dépend de l'âge (croissance), du sexe, de la masse corporelle, de l'état physiologique et surtout du \textbf{niveau d'activité physique} (NAP). Un lutteur de 17 ans (NAP élevé) a un besoin (\SI{14000}{kJ}) bien supérieur à un élève sédentaire (\SI{\approx 10000}{kJ}) ou une personne âgée.
\item \textbf{Le calcul est un outil de décision objectif} : La formule $E = 17G + 17P + 38L$ et le calcul d'écart relatif transforment une impression subjective («~il mange beaucoup~») en un diagnostic scientifique validé (écart 1,4 \%). C'est la démarche du diététicien/nutritionniste.
\item \textbf{Les aliments locaux suffisent pour l'équilibre} : Avec les produits du terroir camerounais (tubercules : manioc, plantain, igname ; céréales : riz, maïs ; légumineuses : haricots, arachides, soja ; légumes-feuilles : pondu, ndolé, folong ; protéines animales accessibles : poisson fumé, œufs, lait ; fruits : banane, avocat, papaye, mangue), on peut composer des rations parfaitement équilibrées, à condition de \textbf{les associer judicieusement} (céréales + légumineuses = protides complets ; vitamine C + fer végétal = absorption).
\item \textbf{La démarche scientifique en nutrition :} Observer (fatigue, contre-performances) $\rightarrow$ Hypothèse (déséquilibre alimentaire) $\rightarrow$ Expérimenter/Calculer (analyse ration, calculs énergétiques) $\rightarrow$ Conclure (déficit protidique, excès lipides, manque protecteurs) $\rightarrow$ Proposer/Valider (menu corrigé, vérification 55/30/15, écart $<10\,\%$).
\end{itemize}

Tu es maintenant prêt à appliquer cette démarche rigoureuse à d'autres situations du programme : \textbf{ration d'une femme enceinte/allaitante} (besoins accrus en fer, acide folique, Ca, protides, énergie), \textbf{menu d'un enfant en bas âge} (sevrage, densité énergétique, allergies), \textbf{alimentation d'un malade en convalescence} (hyperprotidique, fractionnée, digestible) ou \textbf{gestion d'une pathologie métabolique} (diabète de type 2 : index glycémique, charge glycémique, répartition glucidique).

\newpage

\coursec{Méthodes et exercices résolus}

\coursec{Amélioration de la santé de la nutrition}

\courssub{Les constituants de la matière vivante}

\begin{definition}[Constituants de la matière vivante]
Tout être vivant, animal ou végétal, est constitué de trois catégories de substances fondamentales :
\begin{itemize}
\item l'\cle{eau} (solvant universel du métabolisme, $60$ à $90\,\%$ de la masse fraîche) ;
\item la \cle{matière organique} (composés carbonés : glucides, lipides, protéines, acides nucléiques, vitamines) ;
\item la \cle{matière minérale} (sels minéraux : calcium, phosphore, potassium, fer, magnésium...), mise en évidence sous forme de \cle{cendres} après calcination.
\end{itemize}
\end{definition}

\begin{experience}[Mise en évidence des constituants par calcination (TP type)]
\textbf{Matériel :} tube à essai, bec Bunsen, pince à tube, balance de précision (\SI{0,01}{g}), échantillon frais (graine de maïs, feuille, muscle, lait), sulfate de cuivre anhydre (blanc) pour tester l'eau.
\textbf{Protocole :}
\begin{enumerate}
\item Peser l'échantillon frais : masse $m_1$.
\item Chauffer doucement : la buée sur les parois supérieures condense ; approcher le sulfate de cuivre anhydre $\rightarrow$ devient bleu (preuve de l'eau).
\item Dessécher à l'étuve à \SI{100}{\celsius} jusqu'à masse constante $m_2$ (masse sèche).
\item Calciner fortement (flamme vive) jusqu'à disparition des fumées et odeurs : la matière organique s'oxyde (noircissement puis combustion).
\item Laisser refroidir et peser le résidu : masse $m_3$ (cendres = matière minérale).
\end{enumerate}
\textbf{Interprétation :}
\begin{itemize}
\item Eau $= m_1 - m_2$ ;
\item Matière organique $= m_2 - m_3$ ;
\item Matière minérale $= m_3$.
\end{itemize}
\end{experience}
\begin{popfigure}
\begin{tikzpicture}[scale=0.9, every node/.style={font=\small}]
% Support
\draw[thick] (0,0) -- (6,0);
% Tube à essai
\draw[thick] (3,0.2) -- (3,3.5) -- (4,3.5) -- (4,0.2);
\draw[thick] (3,0.2) .. controls (3.5,0) and (3.5,0) .. (4,0.2); % fond arrondi
% Échantillon
\fill[poporange!50] (3.1,0.4) rectangle (3.9,1.2);
\node at (3.5,0.8) {Échantillon frais};
% Buée
\draw[popblue!50, decorate, decoration={snake, amplitude=1pt, segment length=4pt}] (3.1,2.8) -- (3.9,2.8);
\node[popblue] at (3.5,3.1) {Buée (H$_2$O)};
% Bec Bunsen
\draw[thick] (2.5,-0.5) -- (2.5,0);
\draw[thick] (3.5,-0.5) -- (3.5,0);
\draw[poporange, decorate, decoration={zigzag, amplitude=2pt, segment length=3pt}] (3,-0.5) -- (3,0);
\node at (3,-0.8) {Flamme};
% Sulfate de cuivre (tube test eau)
\draw[thick] (5,2.5) -- (5,3.5) -- (5.5,3.5) -- (5.5,2.5);
\draw[thick] (5,2.5) .. controls (5.25,2.3) and (5.25,2.3) .. (5.5,2.5);
\fill[white] (5.05,2.6) rectangle (5.45,3.4);
\node[rotate=90] at (5.25,3.0) {CuSO$_4$ blanc $\rightarrow$ bleu};
% Cendres final
\fill[popmuted] (3.1,0.4) rectangle (3.9,0.8);
\node at (3.5,0.6) {Cendres (minéraux)};
\end{tikzpicture}
\legende{Mise en évidence de l'eau (buée, sulfate de cuivre) et séparation matière organique / matière minérale par calcination.}
\end{popfigure}

\begin{methode}[Analyser quantitativement les constituants d'un échantillon biologique]
\begin{enumerate}
\item \textbf{Peser} la masse fraîche $m_1$.
\item \textbf{Dessécher} à l'étuve (\SI{100}{\celsius}) jusqu'à masse constante $m_2$ $\rightarrow$ eau $= m_1 - m_2$.
\item \textbf{Calciner} le résidu sec jusqu'à masse constante $m_3$ (cendres) $\rightarrow$ matière minérale $= m_3$ ; matière organique $= m_2 - m_3$.
\item \textbf{Calculer} les pourcentages par rapport à la masse fraîche $m_1$ : $\% = \dfrac{\text{masse constituant}}{m_1} \times 100$.
\item \textbf{Vérifier} : $\%\,\text{eau} + \%\,\text{organique} + \%\,\text{minéral} = 100\,\%$.
\end{enumerate}
\textbf{Réflexes :} toujours travailler sur masse fraîche pour les pourcentages ; noter l'odeur caractéristique (brûlé, protéines) lors de la calcination ; le sulfate de cuivre anhydre est un test spécifique de l'eau (blanc $\rightarrow$ bleu).
\end{methode}

\begin{exoresolu}[Analyse quantitative d'une graine de maïs]
On pèse une graine de maïs fraîche : $m_1 = \SI{2,0}{g}$. Après séchage à l'étuve (\SI{100}{\celsius}), masse $m_2 = \SI{1,7}{g}$. Après calcination, masse des cendres $m_3 = \SI{0,02}{g}$.
\begin{enumerate}
\item Calculer le pourcentage d'eau, de matière organique et de matière minérale.
\item Conclure sur la composition de la matière vivante.
\end{enumerate}
\tcblower
\textbf{1. Calculs :}
\begin{itemize}
\item \textbf{Eau} : $m_{\text{eau}} = m_1 - m_2 = 2,0 - 1,7 = \SI{0,3}{g}$ ; $\%\,\text{eau} = \dfrac{0,3}{2,0} \times 100 = 15\,\%$.
\item \textbf{Matière minérale (cendres)} : $m_{\text{min}} = m_3 = \SI{0,02}{g}$ ; $\%\,\text{min} = \dfrac{0,02}{2,0} \times 100 = 1\,\%$.
\item \textbf{Matière organique} : $m_{\text{org}} = m_2 - m_3 = 1,7 - 0,02 = \SI{1,68}{g}$ ; $\%\,\text{org} = \dfrac{1,68}{2,0} \times 100 = 84\,\%$.
\end{itemize}
Vérification : $15 + 84 + 1 = 100\,\%$.

\textbf{2. Conclusion :} La graine de maïs, comme toute matière vivante, associe \textbf{eau} ($15\,\%$), \textbf{matière organique} ($84\,\%$, majoritairement amidon = glucides) et \textbf{matière minérale} ($1\,\%$). Les proportions varient selon l'organe et l'espèce (ex. : feuille $\approx 80\,\%$ eau, graine oléagineuse $\approx 40\,\%$ lipides).
\end{exoresolu}

\courssub{Aliments simples et aliments composés}

\begin{definition}[Aliment simple et aliment composé]
Un \cle{aliment simple} est un aliment naturel non transformé (ou peu transformé) qui apporte principalement un type de nutriment dominant (ex. : riz, huile, œuf, orange). Un \cle{aliment composé} résulte de la transformation culinaire ou industrielle associant plusieurs aliments simples (ex. : ndolé préparé, pain, beignet, sauce arachide).
\end{definition}

\begin{propriete}[Classification fonctionnelle des aliments simples (Groupes alimentaires)]
Les aliments simples se répartissent en quatre groupes selon leur rôle physiologique dominant :
\begin{center}
\begin{tabularx}{\linewidth}{|l|X|X|}
\hline
\rowcolor{popblueL} \textbf{Groupe} & \textbf{Nutriment dominant} & \textbf{Exemples locaux (Cameroun)} \\
\hline
\textbf{Énergétiques (Glucides)} & Glucides complexes (amidon) & Maïs, riz, manioc, igname, banane plantain, taro, mil, sorgho, pain \\
\hline
\textbf{Énergétiques (Lipides)} & Lipides (acides gras) & Huile de palme, huile d'arachide, huile de coton, beurre, margarine \\
\hline
\textbf{Constructeurs (Bâtisseurs)} & Protéines (animales ou végétales) & Viande, poisson (frais/fumé), œufs, lait, fromage, arachide, niébé, soja, sésame \\
\hline
\textbf{Protecteurs (Fonctionnels)} & Vitamines, sels minéraux, fibres & Légumes-feuilles (ndolé, folong, éru, kelen), fruits (mangue, orange, goyave, ananas, papaye), carotte, tomate \\
\hline
\textbf{Eau} & H$_2$O & Eau de source, eau traitée, jus naturels non sucrés \\
\hline
\end{tabularx}
\end{center}
\textbf{Remarque :} Certains aliments simples ont un \textbf{double rôle} (ex. : arachide = constructeur \textbf{et} énergétique lipidique ; lait = constructeur \textbf{et} protecteur calcique ; niébé = constructeur \textbf{et} énergétique glucidique).
\end{propriete}

\begin{methode}[Classer un aliment simple dans son groupe fonctionnel]
\begin{enumerate}
\item Identifier l'aliment \textbf{simple} (ne jamais classer un plat composé).
\item Déterminer son \textbf{nutriment majoritaire} (composition moyenne connue ou étiquette).
\item L'affecter au groupe correspondant : \textbf{Énergétique} (Glucides ou Lipides), \textbf{Constructeur}, \textbf{Protecteur}, \textbf{Eau}.
\item Préciser le cas échéant le \textbf{double rôle}.
\end{enumerate}
\end{methode}

\begin{exoresolu}[Classification d'aliments du marché de Mokolo]
Liste : huile de palme, niébé, banane plantain, éru, œuf, sucre, eau, arachide.
\tcblower
\begin{center}
\begin{tabularx}{\linewidth}{|l|X|}
\hline
\rowcolor{popblueL} \textbf{Groupe fonctionnel} & \textbf{Aliments} \\
\hline
Énergétiques (Glucides) & Banane plantain, sucre \\
\hline
Énergétiques (Lipides) & Huile de palme \\
\hline
Constructeurs (Protéines) & Niébé, œuf, arachide \\
\hline
Protecteurs (Vitamines, Sels minéraux, Fibres) & Éru (légume-feuille) \\
\hline
Eau & Eau \\
\hline
\end{tabularx}
\end{center}
\textbf{Précision :} L'arachide est riche en protéines (\textbf{constructeur}) \textbf{et} en lipides (\textbf{énergétique lipide}). Le niébé apporte protéines et glucides. Le sucre est un glucide simple (énergie rapide) ; la banane plantain apporte de l'amidon (énergie lente) et du potassium.
\end{exoresolu}

\courssub{Caractéristiques physicochimiques et tests d'identification}

\begin{propriete}[Tests caractéristiques des groupes d'aliments simples]
\begin{center}
\begin{tabularx}{\linewidth}{|l|l|X|l|}
\hline
\rowcolor{popblueL} \textbf{Groupe ciblé} & \textbf{Réactif} & \textbf{Condition / Observation positive} & \textbf{Conclusion} \\
\hline
Amidon (Glucide complexe) & Eau iodée (jaune-orangé) & À froid : coloration \textbf{bleue-noire} (violet intense) & Présence d'amidon \\
\hline
Glucose / Sucres réducteurs & Liqueur de Fehling (bleue, Cu$^{2+}$) & \textbf{Chauffage} : précipité \textbf{rouge brique} (Cu$_2$O) & Présence de sucres réducteurs \\
\hline
Lipides & Papier acétate (ou filtre) & Tache \textbf{translucide}, persistante, inextensible à la chaleur & Présence de lipides \\
\hline
Protéines & Réactif de Biuret (NaOH + CuSO$_4$) & À froid : coloration \textbf{violette} / pourpre & Présence de liaisons peptidiques \\
\hline
Eau & Sulfate de cuivre anhydre (blanc) & Devient \textbf{bleu} (hydratation) & Présence d'eau \\
\hline
\end{tabularx}
\end{center}
\textbf{Attention :} La liqueur de Fehling \textbf{doit être chauffée} (ébullition) pour réduire le Cu$^{2+}$ en Cu$_2$O. L'eau iodée ne réagit \textbf{pas} avec le glucose (seule l'amidon donne le bleu-noir).
\end{propriete}

\begin{methode}[Réaliser et interpréter une série de tests d'identification]
\begin{enumerate}
\item Préparer une \textbf{solution test} (écraser l'aliment dans l'eau, filtrer si nécessaire).
\item Effectuer \textbf{chaque test} sur une aliquote distincte : noter \textbf{réactif}, \textbf{condition} (chauffage ?), \textbf{observation} (couleur, précipité, tache).
\item \textbf{Conclure} pour chaque test : « Test X positif $\rightarrow$ présence de [nutriment] » ou « Test X négatif $\rightarrow$ absence de [nutriment] ».
\item \textbf{Synthétiser} : lister les aliments simples mis en évidence.
\end{enumerate}
\textbf{Rédaction type :} « L'ajout d'eau iodée sur la solution test a entraîné une coloration bleue-noire, signe de la présence d'amidon. Le chauffage avec la liqueur de Fehling n'a pas donné de précipité rouge brique : absence de sucres réducteurs. »
\end{methode}

\begin{exoresolu}[Analyse d'aliments locaux : pâte d'arachide et bouillie de maïs sucrée]
Résultats obtenus au laboratoire :
\begin{center}
\begin{tabularx}{\linewidth}{|l|X|X|}
\hline
\rowcolor{popblueL} \textbf{Test} & \textbf{Pâte d'arachide} & \textbf{Bouillie de maïs sucrée} \\
\hline
Eau iodée & Coloration bleue-noire & Coloration bleue-noire \\
\hline
Papier acétate & Tache translucide persistante & Pas de tache \\
\hline
Réactif Biuret & Coloration violette & Coloration violette faible \\
\hline
Liqueur de Fehling + chauffage & Pas de précipité & Précipité rouge brique \\
\hline
\end{tabularx}
\end{center}
\tcblower
\textbf{Pâte d'arachide :}
\begin{itemize}
\item Amidon $+$ (eau iodée) ; Lipides $+$ (papier acétate) ; Protéines $+$ (Biuret) ; Glucose $-$ (Fehling).
\item \textbf{Conclusion :} Aliment \textbf{énergétique} (amidon + lipides) \textbf{et constructeur} (protéines). Très complet.
\end{itemize}
\textbf{Bouillie de maïs sucrée :}
\begin{itemize}
\item Amidon $+$ (eau iodée) ; Lipides $-$ ; Protéines $+$ (faible) ; Glucose $+$ (Fehling : hydrolyse du saccharose ajouté par la chaleur/salive).
\item \textbf{Conclusion :} Aliment essentiellement \textbf{énergétique} (glucides complexes + simples), pauvre en lipides et protéines.
\end{itemize}
\end{exoresolu}

\courssub{Rôle des aliments et alimentation équilibrée}

\begin{definition}[Alimentation équilibrée]
Une \cle{alimentation équilibrée} est une alimentation qui couvre \textbf{quantitativement} (énergie, macronutriments) et \textbf{qualitativement} (micronutriments, eau, fibres) les besoins de l'organisme, en tenant compte de l'\textbf{âge}, du \textbf{sexe}, de l'\textbf{état physiologique} (grossesse, allaitement, croissance), de l'\textbf{activité physique} et de l'\textbf{état de santé} (pathologies).
\end{definition}

\begin{propriete}[Rôles physiologiques des groupes d'aliments]
\begin{itemize}
\item \textbf{Aliments énergétiques (Glucides/Lipides)} : Fournissent l'ATP nécessaire au travail musculaire, à la thermorégulation (\SI{37}{\celsius}), à la biosynthèse et au fonctionnement cérébral. $1\,\text{g}$ glucides $= \SI{17}{kJ}$ ; $1\,\text{g}$ lipides $= \SI{38}{kJ}$.
\item \textbf{Aliments constructeurs (Protéines)} : Apportent les acides aminés indispensables pour la synthèse des protéines structurales (muscles, collagène, kératine), enzymatiques, hormonales (insuline), de transport (hémoglobine) et de défense (anticorps). $1\,\text{g}$ protéines $= \SI{17}{kJ}$.
\item \textbf{Aliments protecteurs (Vitamines, Sels minéraux, Fibres)} : Cofacteurs enzymatiques (vitamines B), antioxydants (vit C, E, $\beta$-carotène), régulation osseuse (Ca, P, vit D), transport O$_2$ (Fe), équilibre hydrique (Na, K), transit intestinal (fibres).
\item \textbf{Eau} : Milieu réactionnel (hydrolyse, synthèse), transport sanguin/lymphatique, élimination rénale (urée), thermorégulation (sueur).
\end{itemize}
\end{propriete}

\begin{savaistu}[Le "Kwashiorkor" et le "Marasme" : visages de la malnutrition au Cameroun]
Le \textbf{kwashiorkor} (carence protéique avec apport calorique suffisant) se manifeste par un œdème, une dermatose, une hépatomégalie et une chevelure décolorée ("cheveux en drapeau"). Le \textbf{marasme} (carence globale énergétique et protéique) entraîne un amaigrissement extrême, une atrophie musculaire et une immunodépression sévère. Ces deux formes de malnutrition aiguë sévère (MAS) touchent encore des enfants de moins de 5 ans dans les régions de l'Extrême-Nord et du Nord (insécurité alimentaire, conflits, mauvaises pratiques d'allaitement). La prise en charge en Unité de Nutrition Thérapeutique Ambulatoire (UNTA) ou en hôpital (CRENI) utilise des aliments thérapeutiques prêts à l'emploi (ATPE/Plumpy'Nut\textregistered) riches en protéines, lipides, vitamines et minéraux.
\end{savaistu}

\begin{methode}[Adapter l'alimentation à un état physiologique ou pathologique]
\begin{enumerate}
\item \textbf{Identifier le facteur} : croissance (adolescent), grossesse/allaitement, vieillissement, sport intense, pathologie (diabète, HTA, anémie, obésité, VIH/Sida).
\item \textbf{Modifier les apports ciblés} :
\begin{itemize}
\item \textbf{Croissance / Grossesse / Allaitement} : $\uparrow$ Protéines (construction fœtus/lait), $\uparrow$ Fer (hémoglobine, réserves fœtales), $\uparrow$ Calcium (ossification), $\uparrow$ Acide folique (vit B9, fermeture tube neural), $\uparrow$ Énergie (modérée).
\item \textbf{Sportif / Travailleur de force} : $\uparrow\uparrow$ Glucides (glycogène musculaire/hépatique), $\uparrow$ Eau (pertes sudorales), $\uparrow$ Protéines (réparation micro-lésions).
\item \textbf{Diabète type 2} : $\downarrow$ Sucres rapides (index glycémique bas), $\uparrow$ Fibres (légumineuses, céréales complètes), répartition glucidique sur 3 repas + collations.
\item \textbf{Hypertension artérielle (HTA)} : $\downarrow\downarrow$ Sel (NaCl $< \SI{5}{g/j}$), $\uparrow$ K (fruits/légumes), $\downarrow$ Graisses saturées/alcool.
\item \textbf{Anémie ferriprive} : $\uparrow$ Fer héminique (viande rouge, foie, sang, poisson) + Vit C (absorption) ; éviter thé/café au repas (tanins bloquent Fe).
\item \textbf{Obésité} : $\downarrow$ Densité énergétique (lipides, sucres), $\uparrow$ Volume/satiété (légumes, fibres), activité physique.
\end{itemize}
\item \textbf{Justifier} chaque changement par le mécanisme physiologique.
\end{enumerate}
\end{methode}

\begin{exoresolu}[Alimentation d'une femme enceinte anémique à Bafoussam]
Femme, 6 mois de grossesse, anémie (Hb $< \SI{11}{g/dL}$).
\begin{enumerate}
\item Expliquer l'augmentation des besoins en protéines, fer, calcium.
\item Proposer 3 aliments locaux à privilégier et 1 à limiter, justifier.
\end{enumerate}
\tcblower
\textbf{1. Besoins augmentés :}
\begin{itemize}
\item \textbf{Protéines} : Synthèse des tissus fœtaux (muscles, organes), placenta, utérus, seins, expansion du volume sanguin maternel.
\item \textbf{Fer} : Hémoglobine pour le volume sanguin maternel accru ($+40\,\%$) + constitution des réserves hépatiques fœtales (6 premiers mois de vie). Risque majeur d'anémie dilutionnelle + carentielle.
\item \textbf{Calcium} : Minéralisation squelette fœtal ($\approx \SI{30}{g}$ Ca total). Si apport insuffisant $\rightarrow$ déminéralisation osseuse maternelle (ostéomalacie).
\end{itemize}
\textbf{2. Aliments locaux :}
\begin{itemize}
\item \textbf{Privilégier :} \textbf{Folong / Ndolé} (Fe non héminique + vit C + folates) ; \textbf{Viande rouge / Foie / Poisson fumé} (Fe héminique bien absorbé + protéines) ; \textbf{Lait caillé "Kindirmou" / Yaourt} (Ca biodisponible + protéines). \textit{Astuce :} ajouter jus de goyave/orange (vit C) au repas pour tripler l'absorption du fer végétal.
\item \textbf{Limiter :} \textbf{Boissons sucrées / Bonbons / Pain blanc} (calories "vides", pic glycémique, risque diabète gestationnel, pas de Fe/Ca).
\end{itemize}
\end{exoresolu}

\courssub{Valeur énergétique des aliments et calcul d'un repas}

\begin{propriete}[Coefficients énergétiques d'Atwater (valeurs de référence Probatoire)]
\begin{center}
\begin{tabular}{|c|c|c|}
\hline
\rowcolor{popblueL} \textbf{Nutriment} & \textbf{kJ/g} & \textbf{kcal/g} \\
\hline
Glucides & 17 & 4 \\
\hline
Lipides & 38 & 9 \\
\hline
Protéines & 17 & 4 \\
\hline
\end{tabular}
\end{center}
Formule : $E_{\text{total}}(\text{kJ}) = 17 \times m_G + 38 \times m_L + 17 \times m_P$ (masses $m$ en grammes de nutriments purs).
Conversion : $1\,\text{kcal} = \num{4,18}\,\text{kJ}$.
\end{propriete}

\begin{methode}[Calculer la valeur énergétique d'un repas]
\begin{enumerate}
\item \textbf{Lister} les aliments simples composant le repas et leur \textbf{masse brute} (g).
\item \textbf{Convertir} en masses de nutriments purs (g) : $m_{\text{nutriment}} = m_{\text{aliment}} \times \frac{\%\,\text{nutriment}}{100}$ (données de composition moyenne ou étiquette).
\item \textbf{Appliquer} les coefficients : $E_G = 17 m_G$, $E_L = 38 m_L$, $E_P = 17 m_P$.
\item \textbf{Sommer} : $E_{\text{repas}} = E_G + E_L + E_P$ (en kJ).
\item \textbf{Convertir} en kcal si demandé : $E_{\text{kcal}} = E_{\text{kJ}} / 4,18$.
\item \textbf{Comparer} aux besoins du repas (ex. : $\frac{1}{3}$ ou $40\,\%$ des besoins journaliers) et \textbf{conclure}.
\end{enumerate}
\end{methode}

\begin{exoresolu}[Repas de cantine scolaire à Yaoundé]
Menu : \SI{150}{g} riz cuit ($25\,\%$ glucides), \SI{50}{g} bœuf ($20\,\%$ prot, $10\,\%$ lip), \SI{10}{g} huile (lipides purs). Besoin repas $\approx \SI{3500}{kJ}$.
\tcblower
\textbf{Masses de nutriments :}
\begin{itemize}
\item Glucides : $150 \times 0,25 = \SI{37,5}{g}$ (riz).
\item Protéines : $50 \times 0,20 = \SI{10}{g}$ (bœuf).
\item Lipides : $50 \times 0,10 = \SI{5}{g}$ (bœuf) $+ \SI{10}{g}$ (huile) $= \SI{15}{g}$.
\end{itemize}
\textbf{Énergie :}
\begin{align*}
E_G &= 37,5 \times 17 = \SI{637,5}{kJ} \\
E_P &= 10 \times 17 = \SI{170}{kJ} \\
E_L &= 15 \times 38 = \SI{570}{kJ} \\
E_{\text{total}} &= \SI{1377,5}{kJ} \quad (\approx \SI{330}{kcal})
\end{align*}
\textbf{Analyse :} $\SI{1377,5}{kJ} < \SI{3500}{kJ}$ (couverture $39\,\%$ seulement).
\textbf{Conclusion :} Repas \textbf{insuffisant}. Il faut augmenter la ration glucidique (ex. : \SI{300}{g} riz ou ajout igname/banane) et/ou ajouter une source lipidique/protéique (arachide dans la sauce).
\end{exoresolu}

\courssub{Rations alimentaires et élaboration de menus équilibrés}

\begin{definition}[Ration alimentaire journalière]
Ensemble des aliments consommés en 24 h. Elle doit être \textbf{équilibrée} (3 groupes + eau à chaque repas), \textbf{variée} (différents aliments dans chaque groupe), \textbf{adaptée} (besoins spécifiques) et \textbf{répartie} sur 3 repas principaux (+ collation si besoin).
\end{definition}

\begin{propriete}[Repartition énergétique journalière type (adulte/adolescent actif)]
\begin{itemize}
\item \textbf{Petit-déjeuner} : $20$ à $25\,\%$ (réveil métabolique, glycémie).
\item \textbf{Déjeuner} : $35$ à $40\,\%$ (repas principal, activité après-midi).
\item \textbf{Dîner} : $30$ à $35\,\%$ (léger, $> 2$ h avant coucher).
\end{itemize}
Besoins journaliers moyens : Adolescent/Adulte actif $\approx \SIrange{10000}{12000}{kJ}$ ($\approx \SIrange{2400}{2900}{kcal}$). Femme enceinte/allaitante / Sportif : $+ \SIrange{500}{1000}{kJ}$. Personne âgée / Sédentaire : $- \SIrange{1000}{2000}{kJ}$.
\end{propriete}

\begin{methode}[Élaborer un menu équilibré pour une situation donnée]
\begin{enumerate}
\item \textbf{Analyser le sujet} : Âge, sexe, poids, taille, activité, pathologie, contexte (internat, village, ville, budget).
\item \textbf{Estimer les besoins journaliers} (kJ) et répartir par repas ($25/40/35\,\%$).
\item \textbf{Construire chaque repas} (assiette idéale) :
\begin{itemize}
\item $1/2$ assiette : \textbf{Légumes / Fruits} (Protecteurs) ;
\item $1/4$ assiette : \textbf{Féculents / Céréales} (Énergétiques Glucides) ;
\item $1/4$ assiette : \textbf{Viande / Poisson / Œuf / Légumineuses} (Constructeurs) ;
\item \textbf{Matière grasse} (cuisson/assaisonnement) modérée ;
\item \textbf{Eau} à volonté.
\end{itemize}
\item \textbf{Varier} les sources (animaux/végétaux, crus/cuits, couleurs) sur la journée et la semaine.
\item \textbf{Vérifier} la couverture des micronutriments (Fe, Ca, Vit A, C, D, B9, Iode, Zn).
\end{enumerate}
\end{methode}

\begin{exoresolu}[Menu pour une élève de 16 ans en internat à Yaoundé (Besoins $\SI{10000}{kJ/j}$)]
\tcblower
\textbf{Menu proposé :}
\begin{center}
\begin{tabularx}{\linewidth}{|l|X|}
\hline
\rowcolor{popblueL} \textbf{Repas} & \textbf{Composition (Aliments locaux)} \\
\hline
\textbf{Petit-déj.} ($\approx \SI{2500}{kJ}$) & Bouillie de maïs/soja enrichie (lait + arachide moulue), beignet niébé, fruit de saison (banane/papaye), eau \\
\hline
\textbf{Déjeuner} ($\approx \SI{4000}{kJ}$) & Riz blanc, sauce pistache (graines de courge) au poisson sec + ndolé (feuilles), avocat, eau \\
\hline
\textbf{Dîner} ($\approx \SI{3500}{kJ}$) & Igname bouillie, omelette tomate/oignon/poivron, salade carotte/chou, mangue, eau \\
\hline
\end{tabularx}
\end{center}
\textbf{Justification :}
\begin{itemize}
\item \textbf{Énergétiques} : Maïs/soja + beignet (matin), Riz (midi), Igname (soir) + Lipides (arachide, huile sauce, avocat, œuf).
\item \textbf{Constructeurs} : Arachide/soja/lait (matin), Poisson sec (midi), Œuf (soir) $\rightarrow$ protéines variées, couverture besoins croissance.
\item \textbf{Protecteurs} : Ndolé (Fe, Ca, Vit A, folates), Avocat (Vit E, K, K), Carotte/Chou (β-carotène, Vit C), Fruits (Vit C, fibres).
\item \textbf{Eau} : Présente à chaque repas.
\item \textbf{Contexte} : Aliments disponibles au marché local, coût maîtrisé, respect habitudes culinaires (sauces, bouillies).
\end{itemize}
\end{exoresolu}

\begin{exercice}[Entraînement : Calcul de ration et adaptation]
\textbf{Sujet :} Un maçon de 30 ans, \SI{70}{kg}, travaille 8 h/j (port de charges). Besoins estimés $\SI{12500}{kJ/j}$.
\begin{enumerate}
\item Proposer un menu équilibré pour une journée (3 repas) avec aliments camerounais.
\item Calculer approximativement l'apport énergétique du déjeuner proposé (donner les masses estimées de nutriments et le total en kJ).
\item Il est diagnostiqué hypertendu. Quelle modification majeure apporter à ce menu ? Justifier.
\end{enumerate}
\end{exercice}

\begin{corrige}[Exercice : Maçon hypertendu]
\textbf{1. Menu proposé (exemple) :}
\begin{itemize}
\item \textbf{Matin :} Bouillie mil/maïs (lait, arachide), beignet haricot, orange, eau.
\item \textbf{Midi :} Fufu maïs/manioc, sauce gombo/feuilles (viande/poisson + crevettes séchées), mangue, eau.
\item \textbf{Soir :} Riz, sauce tomate/poisson fumé + légumes (carotte/haricot vert), yaourt, eau.
\end{itemize}
\textbf{2. Calcul déjeuner (estimation) :}
Fufu \SI{300}{g} ($30\,\%$ glucides $\rightarrow \SI{90}{g}$ G), Sauce : Poisson \SI{100}{g} ($20\,\%$ P $\rightarrow \SI{20}{g}$ P ; $5\,\%$ L $\rightarrow \SI{5}{g}$ L), Huile palme \SI{20}{g} (L), Gombo/feuilles (négligeable énergie).
$E_G = 90 \times 17 = \SI{1530}{kJ}$ ; $E_P = 20 \times 17 = \SI{340}{kJ}$ ; $E_L = (5+20) \times 38 = \SI{950}{kJ}$.
$E_{\text{déj}} \approx \SI{2820}{kJ}$ ($\approx 22,5\,\%$ du total, un peu faible pour $40\,\%$ visé $\rightarrow$ augmenter fufu ou ajouter banane plantain).
\textbf{3. Modification HTA :} \textbf{Réduire drastiquement le sel (NaCl)} : ne pas saler l'eau de cuisson du fufu/riz, limiter le poisson fumé/crevettes séchées (très salés), ne pas ajouter de cube bouillon (glutamate + sel), privilégier poisson frais/viande fraîche, augmenter légumes/fruits (K). Justification : Le Na$^+$ retient l'eau $\rightarrow$ augmentation volume sanguin $\rightarrow$ élévation PA.
\end{corrige}

\begin{attention}[Les 5 erreurs fatales au Probatoire (Série A)]
\begin{enumerate}
\item \textbf{Tests biochimiques :} Confondre \textbf{Eau iodée} (Amidon $\rightarrow$ \textbf{Bleu-noir} À FROID) et \textbf{Liqueur de Fehling} (Glucose/Sucres réducteurs $\rightarrow$ \textbf{Rouge brique} \textbf{APRÈS CHAUFFAGE}). Oublier "chauffage" pour Fehling = 0 point.
\item \textbf{Calcul énergétique :} Utiliser la \textbf{masse brute de l'aliment} au lieu de la \textbf{masse du nutriment pur} (oublier le $\%$ de composition). Inverser les coefficients : Lipides $38\,\text{kJ/g}$ (pas 17) ; Glucides/Protéines $17\,\text{kJ/g}$ (pas 38).
\item \textbf{Classification :} Classer un \textbf{aliment composé} (ex. : "sauce arachide", "ndolé cuit") au lieu de l'aliment simple (arachide, ndolé feuille). Oublier le \textbf{double rôle} de l'arachide, soja, lait, niébé.
\item \textbf{Menu équilibré :} Proposer un menu "copieux" (beaucoup de riz/huile) sans les \textbf{trois groupes à chaque repas} (ex. pas de légumes/fruits au dîner). Absence de \textbf{justification} liée à la situation (âge, pathologie, activité).
\item \textbf{Adaptation pathologique :} Dire "manger moins" pour un diabétique/obèse au lieu de "réduire sucres rapides / lipides, augmenter fibres". Dire "manger plus de fer" pour anémique sans citer la \textbf{Vitamine C} (absorption) ni distinguer fer héminique/non héminique.
\end{enumerate}
\end{attention}

\newpage

\coursec{Exercices}

\courssub{Je vérifie mes connaissances}

\begin{exercice}[QCM : les constituants de la matière vivante]
Pour chaque question, une seule réponse est correcte. Entoure la lettre correspondante.

\textbf{1.} Les constituants minéraux de la matière vivante sont :
\begin{itemize}
\item[a)] l'eau et les sels minéraux ;
\item[b)] les glucides et les lipides ;
\item[c)] les protéines et les vitamines.
\end{itemize}

\textbf{2.} L'élément chimique le plus abondant dans la matière vivante est :
\begin{itemize}
\item[a)] l'azote ;
\item[b)] l'oxygène ;
\item[c)] le carbone.
\end{itemize}

\textbf{3.} Le réactif qui permet de mettre en évidence les protéines est :
\begin{itemize}
\item[a)] l'eau iodée ;
\item[b)] la liqueur de Fehling ;
\item[c)] le réactif de Biuret.
\end{itemize}

\textbf{4.} Un aliment composé est un aliment qui :
\begin{itemize}
\item[a)] contient un seul type de nutriment ;
\item[b)] contient au moins deux types de nutriments en proportions notables ;
\item[c)] ne contient que des sels minéraux et de l'eau.
\end{itemize}
\end{exercice}

\begin{exercice}[Vrai ou faux — justifier]
Réponds par VRAI ou FAUX, puis justifie ta réponse en une ou deux phrases.

\textbf{1.} L'eau iodée devient bleu-noir en présence de glucose.

\textbf{2.} L'huile de palme rouge est un aliment simple à dominante lipidique.

\textbf{3.} Les besoins énergétiques d'un adolescent et d'une femme enceinte sont identiques.

\textbf{4.} Les protéines assurent un rôle plastique (construction et renouvellement des tissus).

\textbf{5.} Un aliment énergétique ne peut jamais jouer un rôle plastique.
\end{exercice}

\begin{exercice}[Définitions]
Définis brièvement et précisément les termes suivants :
\begin{enumerate}
\item aliment simple ;
\item aliment composé ;
\item alimentation équilibrée ;
\item ration alimentaire ;
\item valeur énergétique d'un aliment.
\end{enumerate}
\end{exercice}

\begin{exercice}[Calculs directs]
On rappelle les valeurs énergétiques : glucides : \SI{17}{kJ/g} ; protéines : \SI{17}{kJ/g} ; lipides : \SI{38}{kJ/g}.

\textbf{1.} Calcule la valeur énergétique de \SI{50}{g} de glucides.

\textbf{2.} Calcule la valeur énergétique de \SI{20}{g} de lipides.

\textbf{3.} Un en-cas contient \SI{30}{g} de glucides et \SI{10}{g} de protéines. Calcule sa valeur énergétique totale en kJ, puis convertis-la en kcal (on rappelle : $\SI{1}{kcal} = \SI{4,18}{kJ}$).
\end{exercice}

\courssub{Je m'entraîne}

\begin{exercice}[QCM d'identification par les tests]
On réalise trois tests sur un jus de fruit enrichi commercialisé à Douala. Résultats obtenus :
\begin{itemize}
\item eau iodée : coloration jaune-brun persistante (test négatif) ;
\item liqueur de Fehling chauffée : précipité rouge brique (test positif) ;
\item réactif de Biuret : coloration violette (test positif).
\end{itemize}

\textbf{1.} Pour chaque test, indique le nutriment recherché et conclus sur sa présence ou son absence dans le jus.

\textbf{2.} Ce jus est-il un aliment simple ou un aliment composé ? Justifie.

\textbf{3.} Propose un quatrième test permettant de rechercher les lipides dans ce jus, en précisant le résultat attendu.
\end{exercice}

\begin{exercice}[Classification d'aliments camerounais]
Voici dix aliments couramment consommés au Cameroun : huile de palme, sucre, blanc d'œuf, ndolé, lait, miondo (bâton de manioc), arachide, avocat, haricots, eru.

\textbf{1.} Recopie et complète le tableau ci-dessous en classant chaque aliment.

\begin{center}
\begin{tabularx}{\linewidth}{|l|X|X|}
\hline
\rowcolor{popblueL} \textbf{Aliment} & \textbf{Simple ou composé} & \textbf{Groupe de nutriments dominant} \\
\hline
Huile de palme & & \\
\hline
Sucre & & \\
\hline
Blanc d'œuf & & \\
\hline
Ndolé & & \\
\hline
Lait & & \\
\hline
Miondo & & \\
\hline
Arachide & & \\
\hline
Avocat & & \\
\hline
Haricots & & \\
\hline
Eru & & \\
\hline
\end{tabularx}
\end{center}

\textbf{2.} Parmi ces aliments, lesquels sont des aliments énergétiques ? Lesquels sont des aliments de construction (rôle plastique) ?
\end{exercice}

\begin{exercice}[Valeur énergétique d'un repas]
Le déjeuner d'un élève de Première au lycée de Bafoussam se compose de :
\begin{itemize}
\item glucides : \SI{120}{g} ;
\item protéines : \SI{30}{g} ;
\item lipides : \SI{25}{g}.
\end{itemize}
On rappelle : glucides : \SI{17}{kJ/g} ; protéines : \SI{17}{kJ/g} ; lipides : \SI{38}{kJ/g} ; $\SI{1}{kcal} = \SI{4,18}{kJ}$.

\textbf{1.} Calcule la valeur énergétique apportée par chaque groupe de nutriments.

\textbf{2.} Calcule la valeur énergétique totale de ce repas en kJ, puis en kcal.

\textbf{3.} Le déjeuner doit théoriquement couvrir 35 \% d'une ration quotidienne de \SI{11000}{kJ}. Ce repas convient-il comme déjeuner ? Justifie par le calcul.
\end{exercice}

\begin{exercice}[Protocoles expérimentaux]
Tu disposes au laboratoire du lycée d'une bouillie de maïs et de lait de soja.

\textbf{1.} Décris le protocole complet (matériel, réactif, conditions expérimentales, résultat attendu) permettant de mettre en évidence la présence d'amidon dans la bouillie de maïs.

\textbf{2.} Décris de même le protocole permettant de mettre en évidence les protéines dans le lait de soja.

\textbf{3.} Pourquoi faut-il réaliser un tube témoin (avec de l'eau distillée) dans chaque test ?
\end{exercice}

\courssub{Je me prépare au Probatoire}

\begin{exercice}[La ration d'une femme allaitante — type Probatoire]
Madame Etoundi, mère allaitante résidant à Maroua, a des besoins énergétiques estimés à \SI{11500}{kJ/jour}. Son menu quotidien lui apporte :
\begin{itemize}
\item glucides : \SI{380}{g} ;
\item protéines : \SI{85}{g} ;
\item lipides : \SI{60}{g}.
\end{itemize}
On rappelle : glucides : \SI{17}{kJ/g} ; protéines : \SI{17}{kJ/g} ; lipides : \SI{38}{kJ/g} ; $\SI{1}{kcal} = \SI{4,18}{kJ}$.

\textbf{1.} Définis les termes : ration alimentaire ; alimentation équilibrée.

\textbf{2.} Calcule la valeur énergétique totale du menu de Madame Etoundi en kJ, puis en kcal.

\textbf{3.} Compare cette valeur à ses besoins et conclus : le menu couvre-t-il les besoins de la journée ?

\textbf{4.} En cas de déficit, calcule sa valeur en kJ, puis propose un aliment local riche en lipides (par exemple l'huile de palme ou l'arachide) et calcule la masse de lipides qu'il faudrait ajouter pour combler exactement ce déficit.

\textbf{5.} Explique pourquoi les besoins d'une femme allaitante sont supérieurs à ceux d'une femme non allaitante, en précisant quels nutriments doivent être particulièrement augmentés et pourquoi.
\end{exercice}

\begin{exercice}[Le kwashiorkor dans un village de l'Extrême-Nord — type Probatoire]
Dans un village de la région de l'Extrême-Nord, un enfant de 3 ans, sevré puis nourri presque exclusivement de foutou (pâte de maïs) et d'eau, présente un ventre ballonné, des œdèmes aux membres inférieurs, des cheveux décolorés et un retard de croissance. Le médecin du centre de santé diagnostique un kwashiorkor.

Analyse du régime de l'enfant : glucides : \SI{250}{g/jour} ; protéines : \SI{12}{g/jour} ; lipides : \SI{8}{g/jour}. Les besoins d'un enfant de cet âge sont estimés à \SI{6000}{kJ/jour}, dont au moins \SI{30}{g} de protéines.

\textbf{1.} Calcule la valeur énergétique totale du régime de cet enfant et montre qu'elle peut paraître proche des besoins énergétiques.

\textbf{2.} Montre par le calcul que l'apport protéique est très inférieur aux besoins.

\textbf{3.} Rappelle le rôle plastique des protéines et explique alors les symptômes observés (retard de croissance, œdèmes, fonte musculaire).

\textbf{4.} Déduis-en la définition d'une alimentation équilibrée : pourquoi ne suffit-il pas de couvrir les besoins énergétiques ?

\textbf{5.} Propose deux aliments locaux riches en protéines, accessibles dans cette région, pour améliorer durablement le régime de l'enfant.
\end{exercice}

\begin{exercice}[Enquête au marché et menus équilibrés — type Probatoire]
Dans le cadre d'une semaine de la santé organisée au lycée de Ngaoundéré, les élèves de Première A doivent concevoir des menus équilibrés. On donne les valeurs énergétiques : glucides : \SI{17}{kJ/g} ; protéines : \SI{17}{kJ/g} ; lipides : \SI{38}{kJ/g}.

\textbf{Partie A — Identification des aliments.} Trois échantillons anonymes A, B et C prélevés au marché sont soumis aux tests suivants :

\begin{center}
\begin{tabularx}{\linewidth}{|l|X|X|X|}
\hline
\rowcolor{popblueL} \textbf{Test} & \textbf{Échantillon A} & \textbf{Échantillon B} & \textbf{Échantillon C} \\
\hline
Eau iodée & bleu-noir & jaune-brun & bleu-noir \\
\hline
Liqueur de Fehling (à chaud) & négatif & rouge brique & négatif \\
\hline
Réactif de Biuret & négatif & violet & violet \\
\hline
\end{tabularx}
\end{center}

\textbf{1.} Identifie les nutriments présents dans chaque échantillon.

\textbf{2.} Précise, pour chaque échantillon, s'il s'agit d'un aliment simple ou composé, et propose un exemple d'aliment camerounais correspondant.

\textbf{Partie B — Élaboration d'un menu.} Un élève de 16 ans a des besoins quotidiens de \SI{10500}{kJ}, répartis ainsi : petit-déjeuner 25 \%, déjeuner 35 \%, dîner 40 \%.

\textbf{3.} Calcule la valeur énergétique que doit apporter chacun des trois repas.

\textbf{4.} Le déjeuner proposé contient \SI{100}{g} de glucides, \SI{25}{g} de protéines et \SI{20}{g} de lipides. Vérifie par le calcul s'il respecte la part attendue, et propose un ajustement si nécessaire.

\textbf{5.} Cite les quatre conditions dont dépend une alimentation équilibrée et illustre chacune par un exemple concret.
\end{exercice}

\newpage

\coursec{Corrigés des exercices}

\coursec{Exercices d'application et corrigés détaillés}

\begin{corrige}[Exercice 1 — QCM : les constituants de la matière vivante]
\textbf{1. Réponse a).} Les constituants \cle{minéraux} de la matière vivante sont l'\cle{eau} (environ \SI{65}{\%} de la masse corporelle) et les \cle{sels minéraux} (calcium, phosphore, fer, sodium, potassium...). Les glucides, lipides, protéines et vitamines sont des constituants \emph{organiques}.

\textbf{2. Réponse b).} L'\cle{oxygène} est l'élément chimique le plus abondant dans la matière vivante (environ \SI{65}{\%} de la masse), principalement parce qu'il entre dans la composition de l'eau, très abondante. Viennent ensuite le carbone, l'hydrogène puis l'azote.

\textbf{3. Réponse c).} Le \cle{réactif de Biuret} donne une coloration violette en présence de protéines. L'\cle{eau iodée} met en évidence l'amidon (bleu-noir) et la \cle{liqueur de Fehling} chauffée met en évidence les sucres réducteurs comme le glucose (précipité rouge brique).

\textbf{4. Réponse b).} Un \cle{aliment composé} contient au moins deux types de nutriments en proportions notables (exemple : le lait contient glucides, lipides et protéines). Un aliment qui ne contient qu'un seul type de nutriment est un aliment \emph{simple}.
\end{corrige}

\begin{corrige}[Exercice 2 — Vrai ou faux — justifier]
\textbf{1. FAUX.} L'eau iodée devient bleu-noir en présence d'\emph{amidon}, pas de glucose. Le glucose est mis en évidence par la liqueur de Fehling chauffée (précipité rouge brique).

\textbf{2. VRAI.} L'huile de palme rouge est constituée presque exclusivement de lipides : c'est un aliment simple du groupe des corps gras (aliment énergétique à dominante lipidique).

\textbf{3. FAUX.} Les besoins énergétiques dépendent de l'âge, du sexe, de la masse corporelle, de l'activité physique et de l'état physiologique. Une femme enceinte a des besoins \emph{augmentés} (construction des tissus du fœtus), différents de ceux d'un adolescent.

\textbf{4. VRAI.} Les protéines fournissent les acides aminés nécessaires à la fabrication des protéines de l'organisme : construction des tissus chez l'enfant, renouvellement et réparation des tissus chez l'adulte. C'est le \cle{rôle plastique}.

\textbf{5. FAUX.} Certains nutriments cumulent les deux rôles : les protéines sont plastiques mais peuvent aussi être oxydées pour fournir de l'énergie (\SI{17}{kJ/g}) ; les \cle{phospholipides} entrent dans la constitution des membranes cellulaires. La classification par rôle n'est pas étanche.
\end{corrige}

\begin{corrige}[Exercice 3 — Définitions]
\begin{enumerate}
\item \textbf{Aliment simple :} aliment qui ne contient qu'\emph{un seul type de nutriment} en proportion notable (exemples : huile, sucre, blanc d'œuf).
\item \textbf{Aliment composé :} aliment qui contient \emph{au moins deux types de nutriments} en proportions notables (exemples : lait, arachide, ndolé).
\item \textbf{Alimentation équilibrée :} alimentation qui couvre, en \emph{quantité} (énergie) et en \emph{qualité} (tous les nutriments : glucides, lipides, protéines, vitamines, sels minéraux, eau), les besoins de l'organisme, lesquels varient selon l'âge, la masse corporelle, l'état physiologique et l'état de santé.
\item \textbf{Ration alimentaire :} quantité totale d'aliments nécessaire à un individu pendant \SI{24}{h} pour couvrir ses besoins énergétiques et plastiques.
\item \textbf{Valeur énergétique d'un aliment :} quantité d'énergie libérée par l'oxydation complète d'une masse donnée de cet aliment ; elle s'exprime en kJ (ou kcal) par gramme ou par \SI{100}{g}.
\end{enumerate}
\end{corrige}

\begin{corrige}[Exercice 4 — Calculs directs]
\textbf{1.} Glucides : $E = 50 \times 17 = \SI{850}{kJ}$.

\textbf{2.} Lipides : $E = 20 \times 38 = \SI{760}{kJ}$.

\textbf{3.} En-cas :
\begin{align*}
E_{\text{glucides}} &= 30 \times 17 = \SI{510}{kJ}\\
E_{\text{protéines}} &= 10 \times 17 = \SI{170}{kJ}\\
E_{\text{totale}} &= 510 + 170 = \SI{680}{kJ}
\end{align*}
Conversion : $E = \dfrac{680}{4{,}18} \approx \SI{162,7}{kcal}$.
\end{corrige}

\begin{corrige}[Exercice 5 — QCM d'identification par les tests]
\textbf{1.} Interprétation des tests :
\begin{itemize}
\item \textbf{Eau iodée} : recherche de l'\emph{amidon}. Coloration jaune-brun persistante $\Rightarrow$ test négatif : \textbf{absence d'amidon}.
\item \textbf{Liqueur de Fehling chauffée} : recherche des \emph{sucres réducteurs} (glucose...). Précipité rouge brique $\Rightarrow$ test positif : \textbf{présence de glucose}.
\item \textbf{Réactif de Biuret} : recherche des \emph{protéines}. Coloration violette $\Rightarrow$ test positif : \textbf{présence de protéines}.
\end{itemize}

\textbf{2.} Le jus contient au moins deux types de nutriments en proportions notables (glucides solubles et protéines) : c'est un \textbf{aliment composé}.

\textbf{3.} Test des lipides : déposer une goutte de jus sur un papier (ou frotter sur papier sulfurisé) : l'apparition d'une \textbf{tache translucide persistante} (qui ne s'évapore pas, contrairement à une tache d'eau) traduit la présence de lipides. On peut aussi émulsionner le jus dans l'eau : un trouble laiteux persistant indique des lipides.
\end{corrige}

\begin{corrige}[Exercice 6 — Classification d'aliments camerounais]
\textbf{1.} Tableau complété :

\begin{center}
\begin{tabularx}{\linewidth}{|l|X|X|}
\hline
\rowcolor{popblueL} \textbf{Aliment} & \textbf{Simple ou composé} & \textbf{Groupe de nutriments dominant} \\
\hline
Huile de palme & Simple & Lipides \\
\hline
Sucre & Simple & Glucides \\
\hline
Blanc d'œuf & Simple & Protéines \\
\hline
Ndolé & Composé & Protéines et lipides (+ glucides) \\
\hline
Lait & Composé & Glucides, lipides et protéines \\
\hline
Miondo & Simple (quasi) & Glucides (amidon) \\
\hline
Arachide & Composé & Lipides et protéines \\
\hline
Avocat & Composé & Lipides dominants (+ eau, fibres) \\
\hline
Haricots & Composé & Protéines et glucides (amidon) \\
\hline
Eru & Composé & Eau, sels minéraux, vitamines, fibres (+ huile) \\
\hline
\end{tabularx}
\end{center}

\textbf{2.} \textbf{Aliments énergétiques} (riches en glucides et/ou lipides) : huile de palme, sucre, miondo, arachide, avocat. \textbf{Aliments de construction (rôle plastique)} : blanc d'œuf, haricots, arachide, lait, ndolé (riches en protéines). Remarque : l'arachide et le lait cumulent les deux rôles, ce qui illustre que les catégories ne sont pas exclusives.
\end{corrige}

\begin{corrige}[Exercice 7 — Valeur énergétique d'un repas]
\textbf{1.} Valeur énergétique par groupe :
\begin{align*}
E_{\text{glucides}} &= 120 \times 17 = \SI{2040}{kJ}\\
E_{\text{protéines}} &= 30 \times 17 = \SI{510}{kJ}\\
E_{\text{lipides}} &= 25 \times 38 = \SI{950}{kJ}
\end{align*}

\textbf{2.} Total :
\[E_{\text{totale}} = 2040 + 510 + 950 = \SI{3500}{kJ}\]
Conversion : $E = \dfrac{3500}{4{,}18} \approx \SI{837,3}{kcal}$.

\textbf{3.} Part attendue du déjeuner : $E_{\text{attendue}} = 0{,}35 \times 11000 = \SI{3850}{kJ}$.

Comparaison : $3500 < 3850$ kJ ; déficit : $3850 - 3500 = \SI{350}{kJ}$.

\textbf{Conclusion :} ce repas est légèrement insuffisant pour un déjeuner (il manque \SI{350}{kJ}, soit environ \SI{9}{\%} de la part attendue). On peut le compléter, par exemple, par un fruit ou une tranche d'avocat.
\end{corrige}

\begin{corrige}[Exercice 8 — Protocoles expérimentaux]
\textbf{1. Mise en évidence de l'amidon dans la bouillie de maïs.}
\begin{itemize}
\item \emph{Matériel :} deux tubes à essai, pipette, eau iodée, eau distillée, portoir.
\item \emph{Protocole :} verser environ \SI{2}{mL} de bouillie de maïs dans le tube 1 et \SI{2}{mL} d'eau distillée dans le tube 2 (témoin). Ajouter 2 à 3 gouttes d'eau iodée dans chaque tube. Agiter et observer \emph{à froid} (sans chauffage).
\item \emph{Résultat attendu :} coloration \textbf{bleu-noir} dans le tube 1 (présence d'amidon) ; coloration jaune-brun persistante dans le tube témoin.
\end{itemize}

\textbf{2. Mise en évidence des protéines dans le lait de soja.}
\begin{itemize}
\item \emph{Matériel :} deux tubes à essai, pipettes, réactif de Biuret (ou soude + sulfate de cuivre dilué), eau distillée.
\item \emph{Protocole :} verser \SI{2}{mL} de lait de soja dans le tube 1 et \SI{2}{mL} d'eau distillée dans le tube 2. Ajouter quelques gouttes de réactif de Biuret dans chaque tube, agiter, observer \emph{à froid}.
\item \emph{Résultat attendu :} coloration \textbf{violette} dans le tube 1 (présence de protéines) ; coloration bleu clair (couleur du réactif) dans le témoin.
\end{itemize}

\textbf{3. Rôle du tube témoin :} il sert de \emph{référence négative} : il montre la couleur du réactif en l'absence du nutriment recherché. La comparaison tube expérience / tube témoin permet d'affirmer sans ambiguïté que le changement de couleur observé est bien dû à la présence du nutriment et non au réactif lui-même.
\end{corrige}

\begin{corrige}[Exercice 9 — La ration d'une femme allaitante — type Probatoire]
\textbf{1.} \textbf{Ration alimentaire :} quantité totale d'aliments nécessaire à un individu pendant \SI{24}{h} pour couvrir ses besoins énergétiques et plastiques. \textbf{Alimentation équilibrée :} alimentation qui couvre les besoins de l'organisme en quantité (énergie) et en qualité (tous les nutriments), en tenant compte de l'âge, de la masse corporelle, de l'état physiologique et de l'état de santé.

\textbf{2.} Valeur énergétique du menu :
\begin{align*}
E_{\text{glucides}} &= 380 \times 17 = \SI{6460}{kJ}\\
E_{\text{protéines}} &= 85 \times 17 = \SI{1445}{kJ}\\
E_{\text{lipides}} &= 60 \times 38 = \SI{2280}{kJ}\\
E_{\text{totale}} &= 6460 + 1445 + 2280 = \SI{10185}{kJ}
\end{align*}
Conversion : $E = \dfrac{10185}{4{,}18} \approx \SI{2436,6}{kcal}$.

\textbf{3.} Comparaison : $10185 < 11500$ kJ. Le menu \textbf{ne couvre pas} les besoins de la journée : il est déficitaire.

\textbf{4.} Déficit : $11500 - 10185 = \SI{1315}{kJ}$.

Masse de lipides à ajouter (aliment proposé : huile de palme ou pâte d'arachide) :
\[m = \dfrac{1315}{38} \approx \SI{34,6}{g}\]
Il faudrait donc ajouter environ \textbf{\SI{35}{g} de lipides} (par exemple 3 cuillères à soupe d'huile de palme réparties dans les plats, ou une portion de pâte d'arachide).

\textbf{5.} La femme allaitante produit chaque jour du lait, aliment composé riche en lipides, protéines, glucides (lactose), calcium et eau. Cette production exige un surplus d'énergie et de matière première. Il faut donc augmenter particulièrement :
\begin{itemize}
\item les \textbf{protéines} (rôle plastique : fabrication des protéines du lait, renouvellement des tissus maternels) ;
\item les \textbf{lipides et glucides} (énergie de lactation) ;
\item l'\textbf{eau} et les \textbf{sels minéraux} (calcium) et les \textbf{vitamines}, exportés dans le lait.
\end{itemize}

\textbf{Barème indicatif (sur 10) :} définitions : 2 pts ; calculs énergétiques détaillés : 3 pts ; comparaison et conclusion : 1 pt ; déficit et masse de lipides : 2 pts ; justification physiologique : 2 pts.

\textbf{Points de vigilance :} ne pas oublier la conversion en kcal quand elle est demandée ; expliciter la formule $E = m \times$ valeur énergétique ; la conclusion doit citer les deux valeurs comparées ; pour la question 4, diviser le déficit par \SI{38}{kJ/g} (lipides) et non par 17.
\end{corrige}

\begin{corrige}[Exercice 10 — Le kwashiorkor dans un village de l'Extrême-Nord — type Probatoire]
\textbf{1.} Valeur énergétique du régime :
\begin{align*}
E_{\text{glucides}} &= 250 \times 17 = \SI{4250}{kJ}\\
E_{\text{protéines}} &= 12 \times 17 = \SI{204}{kJ}\\
E_{\text{lipides}} &= 8 \times 38 = \SI{304}{kJ}\\
E_{\text{totale}} &= 4250 + 204 + 304 = \SI{4758}{kJ}
\end{align*}
Cette valeur (\SI{4758}{kJ}, soit environ \SI{80}{\%} des \SI{6000}{kJ} requis) est du \emph{même ordre de grandeur} que les besoins énergétiques : à première vue, l'enfant ne semble pas « manquer de nourriture ».

\textbf{2.} Apport protéique : \SI{12}{g/jour} contre au moins \SI{30}{g/jour} requis, soit un déficit de $30 - 12 = \SI{18}{g/jour}$ : l'enfant ne reçoit que $\dfrac{12}{30} = 0{,}40$, c'est-à-dire \textbf{40 \%} de ses besoins en protéines.

\textbf{3.} Les protéines assurent le \cle{rôle plastique} : elles fournissent les acides aminés nécessaires à la \emph{construction} des tissus (croissance) et à leur \emph{renouvellement} (muscles, enzymes, protéines du sang comme l'albumine). La carence protéique explique :
\begin{itemize}
\item le \textbf{retard de croissance} : faute d'acides aminés, les tissus ne peuvent pas être construits ;
\item la \textbf{fonte musculaire} : l'organisme puise dans ses propres protéines musculaires ;
\item les \textbf{œdèmes} et le ventre ballonné : la baisse des protéines sanguines (albumine) diminue la \cle{pression oncotique} du sang, l'eau quitte les vaisseaux et s'accumule dans les tissus.
\end{itemize}

\textbf{4.} Une \textbf{alimentation équilibrée} ne se résume pas à couvrir les besoins énergétiques : elle doit apporter, en quantités adaptées, \emph{tous} les nutriments (glucides, lipides, protéines, vitamines, sels minéraux, eau). Ici, l'énergie est presque suffisante mais la carence protéique provoque une maladie grave : la \emph{qualité} de la ration compte autant que la \emph{quantité}.

\textbf{5.} Aliments locaux riches en protéines accessibles dans l'Extrême-Nord : le \textbf{niébé} (haricot local), l'\textbf{arachide} (ou la pâte d'arachide), le \textbf{poisson séché} du Logone, le \textbf{lait caillé} (pédi). (Deux réponses suffisent.)

\textbf{Barème indicatif (sur 10) :} calcul énergétique total : 2 pts ; calcul du déficit protéique : 2 pts ; rôle plastique et explication des symptômes : 3 pts ; définition de l'alimentation équilibrée : 2 pts ; aliments locaux : 1 pt.

\textbf{Points de vigilance :} bien distinguer besoins énergétiques (kJ) et besoins plastiques (g de protéines) ; relier explicitement chaque symptôme à la carence ; ne pas écrire que l'enfant « manque de nourriture » en général : c'est une carence \emph{qualitative}.
\end{corrige}

\begin{corrige}[Exercice 11 — Enquête au marché et menus équilibrés — type Probatoire]
\textbf{Partie A.}

\textbf{1.} Nutriments présents :
\begin{itemize}
\item \textbf{Échantillon A :} eau iodée bleu-noir $\Rightarrow$ \emph{amidon} présent ; Fehling négatif (pas de sucre réducteur) ; Biuret négatif (pas de protéines). A contient uniquement de l'\textbf{amidon}.
\item \textbf{Échantillon B :} Fehling rouge brique $\Rightarrow$ \emph{glucose} présent ; Biuret violet $\Rightarrow$ \emph{protéines} présentes ; pas d'amidon.
\item \textbf{Échantillon C :} eau iodée bleu-noir $\Rightarrow$ \emph{amidon} ; Biuret violet $\Rightarrow$ \emph{protéines} ; pas de sucre réducteur.
\end{itemize}

\textbf{2.} Nature des aliments :
\begin{itemize}
\item \textbf{A : aliment simple} (un seul nutriment) $\Rightarrow$ exemple : miondo, tubercule de manioc ou d'igname, farine de maïs.
\item \textbf{B : aliment composé} (glucides + protéines) $\Rightarrow$ exemple : lait, jus de fruit enrichi.
\item \textbf{C : aliment composé} (amidon + protéines) $\Rightarrow$ exemple : haricots, niébé, graines de courge (egusi).
\end{itemize}

\textbf{Partie B.}

\textbf{3.} Répartition des \SI{10500}{kJ} :
\begin{align*}
E_{\text{petit-déjeuner}} &= 0{,}25 \times 10500 = \SI{2625}{kJ}\\
E_{\text{déjeuner}} &= 0{,}35 \times 10500 = \SI{3675}{kJ}\\
E_{\text{dîner}} &= 0{,}40 \times 10500 = \SI{4200}{kJ}
\end{align*}
Vérification : $2625 + 3675 + 4200 = \SI{10500}{kJ}$. Correct.

\textbf{4.} Valeur énergétique du déjeuner proposé :
\begin{align*}
E_{\text{glucides}} &= 100 \times 17 = \SI{1700}{kJ}\\
E_{\text{protéines}} &= 25 \times 17 = \SI{425}{kJ}\\
E_{\text{lipides}} &= 20 \times 38 = \SI{760}{kJ}\\
E_{\text{totale}} &= 1700 + 425 + 760 = \SI{2885}{kJ}
\end{align*}
Comparaison : $2885 < 3675$ kJ ; déficit : $3675 - 2885 = \SI{790}{kJ}$. Le déjeuner \textbf{ne respecte pas} la part attendue.

\emph{Ajustement possible :} ajouter des glucides : $m = \dfrac{790}{17} \approx \SI{46,5}{g}$ de glucides supplémentaires (par exemple une portion de miondo ou de riz en plus) ; ou ajouter $m = \dfrac{790}{38} \approx \SI{20,8}{g}$ de lipides (une cuillère d'huile de palme, un quart d'avocat). Toute combinaison intermédiaire est acceptable.

\textbf{5.} Les quatre conditions d'une alimentation équilibrée :
\begin{itemize}
\item \textbf{l'âge} : un nourrisson et un adolescent n'ont pas les mêmes besoins (croissance rapide) ;
\item \textbf{l'état physiologique} : une femme enceinte ou allaitante a des besoins augmentés (fœtus, lactation) ;
\item \textbf{la masse corporelle} (et l'activité physique) : un manœuvre ou un footballeur dépense plus qu'un employé de bureau sédentaire ;
\item \textbf{l'état de santé} : un malade ou un convalescent a des besoins modifiés (protéines pour la réparation des tissus).
\end{itemize}

\textbf{Barème indicatif (sur 20) :} identification des nutriments : 4 pts ; simple/composé + exemples : 3 pts ; répartition des trois repas : 3 pts ; calcul du déjeuner et ajustement : 5 pts ; conditions et exemples : 5 pts.

\textbf{Points de vigilance :} dans le tableau de tests, raisonner test par test et conclure explicitement « présent / absent » ; vérifier que la somme des trois repas redonne \SI{10500}{kJ} ; pour l'ajustement, préciser \emph{quel} nutriment on ajoute et justifier la masse par un calcul ; chaque condition doit être illustrée par un exemple concret, pas seulement citée.
\end{corrige}

\newpage

\coursec{Fiche bilan}

\begin{popfigure}
\begin{center}\tcbox[colback=popblue,colframe=popblue,coltext=white,arc=9pt,boxsep=4pt,left=6pt,right=6pt]{\bfseries\small Amélioration de la santé par la nutrition}\end{center}
\vspace{-2pt}
\begin{tcbraster}[raster columns=3,raster equal height=rows,raster column skip=6pt,raster row skip=6pt,enhanced,blankest]
\begin{tcolorbox}[enhanced,colback=white,colframe=popblue,boxrule=0.9pt,arc=3pt,title={Constituants matière vivante},fonttitle=\bfseries\scriptsize,coltitle=white,colbacktitle=popblue,left=4pt,right=4pt,top=2pt,bottom=2pt,boxsep=1pt,toptitle=1.5pt,bottomtitle=1.5pt,halign=left]
\begin{itemize}[nosep,leftmargin=*,itemsep=1pt,font=\scriptsize]
\item Minéraux : eau, sels
\item Organiques : G, L, P, V
\end{itemize}
\end{tcolorbox}
\begin{tcolorbox}[enhanced,colback=white,colframe=popgreen,boxrule=0.9pt,arc=3pt,title={Aliments simples \& composés},fonttitle=\bfseries\scriptsize,coltitle=white,colbacktitle=popgreen,left=4pt,right=4pt,top=2pt,bottom=2pt,boxsep=1pt,toptitle=1.5pt,bottomtitle=1.5pt,halign=left]
\begin{itemize}[nosep,leftmargin=*,itemsep=1pt,font=\scriptsize]
\item Simples : 1 constituant
\item Composés : mélanges
\item Classification 3 groupes
\end{itemize}
\end{tcolorbox}
\begin{tcolorbox}[enhanced,colback=white,colframe=poporange,boxrule=0.9pt,arc=3pt,title={Tests identification},fonttitle=\bfseries\scriptsize,coltitle=white,colbacktitle=poporange,left=4pt,right=4pt,top=2pt,bottom=2pt,boxsep=1pt,toptitle=1.5pt,bottomtitle=1.5pt,halign=left]
\begin{itemize}[nosep,leftmargin=*,itemsep=1pt,font=\scriptsize]
\item Fehling/Benedict : G réducteurs
\item Iode : Amidon
\item Biuret : Protéines
\item Soudan III : Lipides
\end{itemize}
\end{tcolorbox}
\begin{tcolorbox}[enhanced,colback=white,colframe=poppurple,boxrule=0.9pt,arc=3pt,title={Rôles \& Équilibre},fonttitle=\bfseries\scriptsize,coltitle=white,colbacktitle=poppurple,left=4pt,right=4pt,top=2pt,bottom=2pt,boxsep=1pt,toptitle=1.5pt,bottomtitle=1.5pt,halign=left]
\begin{itemize}[nosep,leftmargin=*,itemsep=1pt,font=\scriptsize]
\item Énergétique (G, L)
\item Constructeur (P)
\item Protecteur (V, Sels, Eau)
\item Conditions : âge, état, masse
\end{itemize}
\end{tcolorbox}
\begin{tcolorbox}[enhanced,colback=white,colframe=poppink,boxrule=0.9pt,arc=3pt,title={Valeur énergétique},fonttitle=\bfseries\scriptsize,coltitle=white,colbacktitle=poppink,left=4pt,right=4pt,top=2pt,bottom=2pt,boxsep=1pt,toptitle=1.5pt,bottomtitle=1.5pt,halign=left]
\begin{itemize}[nosep,leftmargin=*,itemsep=1pt,font=\scriptsize]
\item G : 17 kJ/g (4 kcal/g)
\item L : 38 kJ/g (9 kcal/g)
\item P : 17 kJ/g (4 kcal/g)
\item VE repas = $\sum m_i \times VE_i$
\end{itemize}
\end{tcolorbox}
\begin{tcolorbox}[enhanced,colback=white,colframe=popgold,boxrule=0.9pt,arc=3pt,title={Rations \& Menus},fonttitle=\bfseries\scriptsize,coltitle=white,colbacktitle=popgold,left=4pt,right=4pt,top=2pt,bottom=2pt,boxsep=1pt,toptitle=1.5pt,bottomtitle=1.5pt,halign=left]
\begin{itemize}[nosep,leftmargin=*,itemsep=1pt,font=\scriptsize]
\item ANC : Besoins journaliers
\item Répartition : 3-4 repas
\item Diversité groupes aliments
\item Adaptation pathologie
\end{itemize}
\end{tcolorbox}
\end{tcbraster}

\legende{Carte mentale de la leçon NA01 : 6 branches principales couvrant le programme officiel 1ère A.}
\end{popfigure}

\begin{bilan}

\courssub{Définitions clés}

\begin{itemize}
  \item \cle{Matière vivante} : Ensemble des substances constituant les êtres vivants. Deux catégories : \textbf{minéraux} (eau \textasciitilde 60-70\%, sels minéraux) et \textbf{organiques} (glucides, lipides, protides, vitamines).
  \item \cle{Aliment simple} : Aliment ne contenant qu'un seul constituant organique majoritaire (ex: huile = lipides, sucre = glucides, blanc d'œuf = protides).
  \item \cle{Aliment composé} : Aliment contenant plusieurs constituants organiques en proportions variables (ex: lait, œuf entier, farine, viande).
  \item \cle{Alimentation équilibrée} : Alimentation qui couvre quantitativement et qualitativement les besoins de l'organisme (énergie, construction, régulation) sans excès ni carence, adaptée à l'âge, l'état physiologique (grossesse, allaitement, sport), la masse corporelle et l'état de santé.
  \item \cle{Valeur énergétique (VE)} : Énergie libérée par l'oxydation complète de 1 g de nutriment. Unités : kJ/g ou kcal/g (1 kcal = 4,184 kJ).
  \item \cle{Ration alimentaire journalière} : Quantité totale d'aliments ingérés en 24 h, répartie en 3 à 4 prises, satisfaisant les \cle{ANC} (Apports Nutritionnels Conseillés).
\end{itemize}

\courssub{Formules et valeurs de référence à connaître par cœur}

\begin{center}
\begin{tabularx}{\linewidth}{|>{\bfseries}l|X|X|}\hline
\rowcolor{popblueL}
\textbf{Nutriment (groupe)} & \textbf{Valeur énergétique} & \textbf{Rôle principal} \\ \hline
Glucides (Énergétiques) & $17\ \text{kJ/g}\ (4\ \text{kcal/g})$ & Carburant cellulaire rapide (ATP) \\ \hline
Lipides (Énergétiques) & $38\ \text{kJ/g}\ (9\ \text{kcal/g})$ & Stockage énergie, structure membranes, précurseurs hormonaux \\ \hline
Protides (Constructeurs) & $17\ \text{kJ/g}\ (4\ \text{kcal/g})$ & Construction/réparation tissus, enzymes, transport, défense \\ \hline
Vitamines, Sels minéraux, Eau (Protecteurs/Régulateurs) & $0$ & Cofacteurs enzymatiques, équilibre hydro-électrolytique, protection immunitaire \\ \hline
\end{tabularx}
\end{center}

\[
\boxed{\text{VE}_{\text{repas}}\ (\text{kJ}) = \sum_{i} \left( m_i\ (\text{g}) \times \text{VE}_i\ (\text{kJ/g}) \right)}
\]
\textbf{Répartition énergétique de référence (adulte sain) :} Glucides 50-55\% ; Lipides 30-35\% ; Protides 12-15\%.

\courssub{Méthodes express}

\begin{methode}[Tester un aliment simple]
  \begin{enumerate}
    \item \textbf{Glucides réducteurs} : Quelques gouttes de \cle{Fehling} (A+B) ou \cle{Benedict} $\rightarrow$ bain-marie 5 min. \textbf{Positif} : Précipité rouge brique (\ce{Cu2O}).
    \item \textbf{Amidon} : Quelques gouttes de \cle{solution iodée} $\rightarrow$ à froid. \textbf{Positif} : Coloration bleu-noir intense.
    \item \textbf{Protéines} : \cle{Biuret} (NaOH + \ce{CuSO4}) $\rightarrow$ à froid. \textbf{Positif} : Coloration violette/rose. \textit{Alternative} : Ninhydrine (chauffage) $\rightarrow$ violet.
    \item \textbf{Lipides} : \cle{Soudan III} (colorant lipophile) $\rightarrow$ mélange. \textbf{Positif} : Coloration rouge-orangé des gouttelettes. \textit{Alternative} : Tache de graisse sur papier buvard (translucide, indélébile).
  \end{enumerate}
\end{methode}

\begin{methode}[Calculer la VE d'un repas]
  \begin{enumerate}
    \item Identifier les aliments composant le repas et leur masse (pesée ou tableau de composition).
    \item Déterminer pour chaque aliment la teneur en G, L, P (étiquette ou table Ciqual/ANC Cameroun).
    \item Calculer l'apport de chaque macronutriment : $m_{\text{G}} \times 17$, $m_{\text{L}} \times 38$, $m_{\text{P}} \times 17$ (en kJ).
    \item Sommer les apports = VE totale du repas. Vérifier la répartition \% par rapport aux repères.
  \end{enumerate}
\end{methode}

\begin{methode}[Élaborer un menu équilibré (ex: adulte 2200 kcal/j)]
  \begin{enumerate}
    \item Fixer les ANC cibles (énergie, G, L, P, fibres, eau, micronutriments).
    \item Répartir sur 3-4 repas : Petit-déj 25\%, Déjeuner 35-40\%, Collation 5-10\%, Dîner 25-30\%.
    \item Choisir au moins 1 aliment de \textbf{chaque groupe} par repas principal (Céréales/Tubercules + Légumineuses/Produits animaux + Légumes/Fruits + Matière grasse + Eau).
    \item Varier les sources (viande/poisson/œuf/légumineuses ; fruits/légumes de saison locaux : papaye, mangue, gombo, folong, manioc, maïs, niébé).
    \item Calculer la VE et les apports en macronutriments du menu proposé ; ajuster les portions si écart > 10\%.
  \end{enumerate}
\end{methode}

\courssub{Classification des aliments simples (3 groupes fonctionnels)}

\begin{center}
\begin{tabularx}{\linewidth}{|>{\bfseries}l|X|X|}\hline
\rowcolor{popgreenL}
\textbf{Groupe} & \textbf{Aliments simples représentatifs} & \textbf{Fonction dominante} \\ \hline
\textbf{Énergétiques} & Sucre, miel, huile, beurre, farine, riz, maïs, manioc, plantain & Fournir de l'ATP (G, L) \\ \hline
\textbf{Constructeurs} & Viande, poisson, œuf, lait, fromage, légumineuses (niébé, soja, arachide) & Construire/réparer (P) \\ \hline
\textbf{Protecteurs / Régulateurs} & Légumes (feuilles, fruits), fruits frais, eau, sel iodé & Réguler, protéger (V, sels, eau, fibres) \\ \hline
\end{tabularx}
\end{center}

\courssub{Repères ANC (ordre de grandeur, adulte 20-40 ans, activité modérée)}
\begin{itemize}
  \item Énergie : \textasciitilde 2200-2600 kcal/j (H), 1800-2200 kcal/j (F).
  \item Protéines : 0,83 g/kg/j (ex: 70 kg $\rightarrow$ 58 g/j).
  \item Lipides : 35-40\% VE totale ; AG saturés < 12\% ; AG trans < 1\%.
  \item Glucides : 45-55\% VE ; Sucres ajoutés < 10\% (idéal < 5\%).
  \item Fibres : 25-30 g/j. Eau : 1,5-2 L/j (boissons + alimentation).
  \item Fer : 11 mg (H), 16 mg (F ménopausée) ; Calcium : 950-1000 mg ; Vitamine D : 10-15 µg.
\end{itemize}

\end{bilan}

\begin{aretenir}[Checklist avant le Probatoire]
  \begin{itemize}
    \item $\square$ Je sais lister les 4 grands constituants organiques + eau + sels minéraux de la matière vivante.
    \item $\square$ Je distingue aliment simple / aliment composé et je donne 3 exemples de chaque.
    \item $\square$ Je connais les 4 tests caractéristiques (réactif, condition, résultat positif) pour G réducteurs, amidon, protéines, lipides.
    \item $\square$ Je récite par cœur les VE : G=17, L=38, P=17 (kJ/g) et je convertis en kcal (÷ 4,184).
    \item $\square$ Je sais calculer la VE totale d'un repas à partir des masses de G, L, P.
    \item $\square$ Je cite les 3 groupes fonctionnels (Énergétiques, Constructeurs, Protecteurs) et leur rôle.
    \item $\square$ Je définis l'alimentation équilibrée en citant les 4 conditions d'adaptation (âge, état physiologique, masse, santé).
    \item $\square$ Je connais la répartition énergétique de référence (G 50-55\%, L 30-35\%, P 12-15\%).
    \item $\square$ Je sais élaborer un menu équilibré sur une journée (3-4 repas) en respectant la diversité des groupes.
    \item $\square$ J'utilise correctement les unités (kJ, kcal, g, \%) et j'écris les formules avec \ce{} (mhchem) si réaction.
    \item $\square$ Je cite 2 exemples d'aliments locaux camerounais par groupe fonctionnel (ex: ndolé/feuilles de manioc = protecteur ; poisson braisé = constructeur ; bâton de manioc = énergétique).
  \end{itemize}
\end{aretenir}

\findecours

\end{document}
